\documentclass[tikz, a4paper,12pt]{book} 
\usepackage[spanish,es-tabla]{babel}	
\usepackage[utf8]{inputenc}
\usepackage{fancyhdr}
\usepackage[dvips]{graphicx}      
\usepackage{float}	  
\usepackage{biblatex}
\usepackage[bf]{caption}
\usepackage{subcaption}
\usepackage{amsmath}
\usepackage{graphicx}
\usepackage{array,tabularx}
\usepackage[table]{xcolor}
\usepackage[most]{tcolorbox}
\usepackage{setspace}
\usepackage[left=2cm, right=2cm, top=3.54cm, bottom=3cm]{geometry}
\usepackage{gensymb}
\usepackage{glossaries}
\usepackage{wasysym}
\usepackage[nice]{nicefrac}
\usepackage{units}
\usepackage{diagbox}
\usepackage{wrapfig}
\usepackage{dsfont}
\usepackage{cancel}
\usepackage{blkarray}
\usepackage{hyperref}
\usepackage{hyperxmp} % Este paquete gestiona los metadatos avanzados del PDF

\hypersetup{
    pdfauthor={Alfonso Fernández Martín},
    pdftitle={Electrónica Física},
    pdfcopyright={Copyright (c) 2026 Alfonso Fernández Martín. Obra bajo licencia CC BY-NC-SA 4.0},
    pdflicenseurl={http://creativecommons.org}
}
\usepackage{tikz}
\usepackage{pgfplots}
\usepackage{circuitikz}
\usepackage{appendix}
\usetikzlibrary{circuits.ee}


%Colores
\definecolor{applegreen}{rgb}{0.55, 0.71, 0.0}
\definecolor{airforceblue}{rgb}{0.36, 0.54, 0.66}
\definecolor{amber}{rgb}{1.0, 0.49, 0.0}

%Redefiniciones y nuevos comandos
\newcommand{\diff}[3][1]{%
	\IfBlankTF{#1}{\frac{\mathrm{d}{#2}}{\mathrm{d}{#3}}}%
	{\frac{\mathrm{d}^{#1} {#2}}{\mathrm{d}{#3}^{#1}}}%
}%

\newcommand{\pdiff}[3][1]{%
	\IfBlankTF{#1}{\frac{\partial {#2}}{\partial {#3}}}%
	{\frac{\partial^{#1} {#2}}{\partial {#3}^{#1}}}%
}%


\graphicspath{{Imagenes/}}

\addbibresource{references.bib}

\bibliography{references}

\setlength{\intextsep}{10pt}

\begin{document}
%Encabezado
\frontmatter
\begin{titlepage}
\begin{center}


\LARGE\textbf{Facultad de Ciencias. Grado en Física.} \vspace{2cm}

\textcolor{airforceblue}{\rule {\linewidth}{0.75mm}}\\
{\Large \textsc{Electrónica Física}}
\textcolor{airforceblue}{\rule {\linewidth}{0.75mm}}\\
\LARGE \textbf{Alfonso Fernández Martín} \vspace{4mm}
\vfill 
\begin{figure}[h]
      \centering
      \includegraphics[width=\linewidth]{imagenportada.jpg}
\end{figure}

\end{center}
\thispagestyle{empty}  %QUITAR NUMERACION
\end{titlepage}\thispagestyle{empty}  %QUITAR NUMERACION
\chapter*{Nota}
El objetivo de hacer estos apuntes no es otro sino el de darles un texto de apoyo a las diapositivas de la asignatura de Electrónica Física, a priori no debería haber erratas ya que he corregido todas las que encontrábamos tanto yo como otros compañeros. En principio están todas las ecuaciones corregidas aunque siempre es bueno contrastar con las diapositivas de clase, también es recomendable añadir comentarios que se hayan hecho en clase ya que alguno se me ha podido escapar y son puntos que restan en el examen.\\\\
Las preguntas principales de teoría (las que han caído otros años) están marcadas con los títulos en \textcolor{red}{rojo} para que sea más fácil identificarlas. Que sean preguntas individuales no quiere decir que no puedan pedirse varias en una misma pregunta de teoría en el examen.\\\\
Los apuntes están divididos en Física de Semiconductores, que gran parte se ve en Física del Estado Sólido I; componentes electrónicos básicos, que es donde mayor cantidad de ejercicios diferentes hay; y dispositivos electrónicos especializados que únicamente es teoría y no hay ejercicios (suelen caer dos puntos de la parte de teoría del examen).\\\\

\begin{flushright}
    Mucha suerte a todos y espero que os sirvan por mucho tiempo.
\end{flushright}

\begin{flushright}
    Chinri
\end{flushright}

\vfill
\begin{flushleft}
        % Nota de derechos y autoría
    \noindent\textbf{Licencia y Derechos de Autor} \\
    \vspace{0.2cm}
    \small
    Este documento y su código fuente son de elaboración propia. 
    Se distribuyen bajo una licencia \\ \textbf{Creative Commons Attribution-NonCommercial-ShareAlike 4.0 International} (CC BY-NC-SA 4.0).
    
    \vspace{0.4cm}
    
    % Esto inserta los iconos típicos de Creative Commons de forma limpia usando texto
    \includegraphics[width=0.25\textwidth]{by-nc-sa.pdf}
    
    \vspace{0.4cm}
    \footnotesize
    Eres libre de compartir y adaptar este material siempre que se reconozca la autoría original, 
    no se utilice con fines comerciales y se distribuya bajo la misma licencia.
\end{flushleft}
\tableofcontents 
\thispagestyle{empty}  %QUITAR NUMERATION
\mainmatter
\setcounter{page}{1}
\part{Propiedades de los materiales Semiconductores}
\pagestyle{fancy}
\fancyhead[R]{Electrónica Física: Tema 1}
\fancyhead[L]{4º Grado en Física}
\chapter{Propiedades Físicas de los Semiconductores}
Vamos a estudiar las propiedades eléctricas de los semiconductores, para ello utilizaremos la ecuación de Schrodinger y los potenciales asociados a las redes cristalinas que forman los Semiconductores. Estudiar estas propiedades a nivel microscópico nos permite: Determinar las propiedades de los electrones dentro de la red cristalina y, a través de la física estadística, entender las propiedades de un conjunto muy grande de los mismos.
\section{Bandas de energía} \label{sec: Bandas de energia}
\subsection{Formación de Bandas} \label{sec: Formacion de bandas}
Para hacernos una idea intuitiva de \textbf{banda de energía}, supongamos que tenemos átomo hidrogenoide aislado, el electrón de este átomo tiene una función de onda asociada $\psi(r)$, que nos permite estudiar la probabilidad que tenemos de encontrar al electrón en un cierto lugar. Si ahora cogemos otro átomo hidrogenoide y lo acercamos lo suficiente como para que la interacción comience a tener relevancia veremos que las funciones de onda comienzan a solaparse. Como nuestro sistema es ligado, la energía del sistema está cuantizada pero al ser fermiones el principio de exclusión de Pauli hace que el nivel de energía donde se encontraban los electrones se divide en dos niveles discretos. Si vamos añadiendo átomos y formamos una red cristalina (periódica en el espacio) los electrones irán ocupando poco a poco distintos niveles de energía que a su vez se irán dividiendo en subniveles discretos.\\

\noindent Esta idea la podemos aplicar a cualquier tipo de átomo que forma nuestra red cristalina. Como el número de electrones es significativamente grande respecto al valor de las magnitudes que tienen los niveles energético, por ejemplo, si tenemos $10^{19}$ electrones en un nivel con energía 1 eV, cada subnivel se diferenciará en $10^{-19}$ eV, por lo que podemos tratar esta banda como una región continua de energía.
\newpage
\subsection{Bandas de valencia y de condución} \label{sec: Bandas de valencia y de conduccion}
\begin{wrapfigure}{r}{0.3\textwidth}
	\centering
	\includegraphics{Standalone/T1_01_Banda_de_valencia_y_conduccion.pdf}
	\caption{Diagramas de la banda de valencia y la banda de conducción}
	\label{fig: Banda de valencia y conduccion}
\end{wrapfigure}


Para átomos multielectrónicos como el Si=[Ne] $3s^2 \ 3p^2$, los electrones que nos interesa estudiar son los de la última capa, conocida como \textbf{capa de valencia}, en este caso están en el nivel energético n=3, vemos que existen dos electrones en el nivel \textit{s} que solo pueden ocupar dos estados posibles (\textit{l}=0), mientras los otros dos electrones que están en el nivel \textit{p} (\textit{l}=1) pueden ocupar 6 posibles estados. Volviendo a la idea anterior, cuando acercamos más átomos de Si y formamos una red, los electrones que se encuentran en la capa de valencia interaccionan entre sí se forman dos bandas, denominadas \textbf{banda de valencia} y \textbf{banda de conducción}.
\noindent Supongamos ahora que la red cristalina está formada por $N$ átomos, las bandas de energía, se volverán dividir en niveles discretos, pero no tendrán $2N$ y $6N$ posibles estados, sino la banda más baja (valencia) tendrá $4 N$  estados y la banda superior (conducción) tendrá $4 N$ estados. La diferencia de enrgía que presentan las barras es lo que se conoce como GAP del semiconductor y es el ancho de una \textbf{banda de energía prohibida}.\\

\noindent Para acabar si hacemos que la temperatura de la red sea T=0 K, los electrones estarán quietos en el mínimo nivel de energía por lo tanto, todos estarán en la banda de valencia dejando vacía la banda de conducción., esto lo volveremos a utilizar cuando definamos el nivel de Fermi.
\subsubsection{\textcolor{red}{Cálculo ``estricto'' de las bandas de energía}} \label{sec: Calculo de bandas}
Para poder calcular las bandas de energía que presenta una red cristalina utilizamos dos aproximaciones:
\begin{itemize}
	\item \textbf{Hipótesis adiabática}: los iones de la red tienen posiciones fijas, únicamente se pueden mover los electrones de la capa de valencia.
	\item \textbf{Aproximación de partícula única}: se estudia el la dinámica de un électrón sometido a un potencial eficaz global que es periódico $V(\mathbf{r})=V(\mathbf{r+R})$
\end{itemize}
Cada elemento del potencial periódico es el potencial de Coulomb \footnote{Modelo de Kronig-Penney} y la ecuación de Schrodinger tiene como solución las conocidas como fucniones de Bloch
\begin{equation}
	\left[-\frac{\hbar}{2\,m} \nabla^2+V(\mathbf{r})\right]\psi(\mathbf{r})=E\psi(\mathbf{r}) \xrightarrow{\mathrm{Solucion}} \psi_{n\mathbf{k}}(\mathbf{r})=e^{i\mathbf{k \, r}}u_{n \mathbf{k} \,\sigma}(\mathbf{r})
	\label{eq: Schrodinger y Bloch}
\end{equation}
Como bien sabemos, las celdas del cristal no tienen extensión infinita, por lo que podemos establecer condiciones de contorno (es lo que nos permite obtener estados ligados), de la solución de \ref{eq: Schrodinger y Bloch}, y suponiendo que nos encontramos en un cristal cúbico simple de arista $a$ (es decir una partícula en una caja) la forma del vector $$\mathbf{k}=\frac{\mathbf{p}\, \pi}{N\, a}$$
\subsubsection{\textcolor{red}{Propiedades de las bandas de energía}}

Vamos a determinar algunas propiedades de los autovalores $E_n(\mathbf{k})$, como únicamente vamos a estudiar la celda unidad, por lo que tiene la periodicidd de la 1ª Zona de Brillouin, además es simétrica en \textbf{k} y en los bordes de la ZB $\nabla_\mathbf{k} E_n(\mathbf{k})=0$.\\

Los autovalores de la energía son formas cuadráticas:
\begin{equation}
	E_n(\mathbf{k})=	E_n(\mathbf{k_0})+\frac{1}{2}\left. \sum_{i, \ j} \frac{\partial ^2 E}{\partial k_i \partial k_j}\right |_\mathbf{k_0} (k_i-k_{i0})(k_j-k_{j0})
 \label{eq: Forma bandas}
\end{equation}
Esta expresión corresponde a un elipsoide de revolución (2 coeficientes iguales y uno distinto). Si la forma cuadrática sólo tiene términos diagonales no nulos, adquiere forma de esfera (3 coeficientes iguales). Además la dependencia $E_n \propto \mathbf{k}^2$.
\subsubsection{\textcolor{red}{Propiedades de las bandas del Si y GaAs}} 
\begin{wrapfigure}{r}{0.4\textwidth}
	\centering
	\includegraphics[width=1\linewidth]{Standalone/T1_11_Estructura_de_bandas_GaAs_Si.pdf}
	\caption{Estructura de bandas del GaAs (a) y Si (b)}
	\label{fig: estructuras de bandas}
\end{wrapfigure}
Si analizamos las gráficas de la figura \ref{fig: estructuras de bandas} vemos que hay una clara diferencia entre ambos semiconductores.\\
Estas figuras se representan sobre el espacio recíproco \textbf{k}. Por lo los desplazamientos horizontales implican que el momento no es constante dentro del cristal.
\begin{itemize}
	\item Para el caso del GaAs decimos que el semiconductor es de GAP \textbf{directo}, mientras que para el caso del Si decimos que es de GAP \textbf{indirecto}.
	\item El GAP directo implica que no hay cambios en el momento del cristal lo que le confiere propiedades ópticas muy interesantes.
    \item El GAP indirecto implica que debemos tener en cuenta la interacción de los electrones con el crisatal al cambiar de banda entre la banda de condcucción y la de valencia.
\end{itemize}

\begin{itemize}
    \item Para el caso del Si, tenemos que el mínimo en la capa de conducción se encuentra en la dirección $\Gamma-\mathrm{X}[1, \ 0, \ 0]$. Y por lo tando tienen forma de elipsoides de revolución (en este caso 6) 
    \item Para el GaAs tenemos un único mínimo en el punto $\Gamma$, 4 en el punto L y 3 en el punto X que son equivalentes en cada punto, su banda de conducción es esférica.
    \item Para ambos SC las bandas de valencia tienen forma esférica.
\end{itemize}
\noindent Otro ejemplo de GAP indirecto es el Germanio, mientras que otros ejemplos de GAP directo son el GaP y el AlAs.

\section{Conducción eléctrica en sólidos} \label{sec: Conduccion electrica en solidos}
\subsection{Modelo de enlace} \label{sec: Modelo de enlace}
Para explicar el modelo de enlace volvemos a la red de Si de la sección 1.1.2 representada en la figura, para ello supongamos que nos encontramos con una temperatura T=0 K. Ahora supongamos que la temperatura aumenta ligeramente, esto hará que alguno de los electrones que se encuentran en la capa de valencia, ganen energía por excitación térmica, esto genera un hueco como muestra en la figura \ref{fig: Red cristalina sin e} lo que los electrones de la capa de valencia ocupen dicho hueco y globalmente observemos que el hueco se desplaza a lo largo de la red en el sentido contrario en el que se desplazan los electrones.

\begin{figure}[h]
	\centering
	\begin{subfigure}[h]{0.3\textwidth}
		\includegraphics{Standalone/T1_02_Red_cristalina.pdf}
		\caption{La red está a T=0 K y los electrones están en la CV.}
		\label{fig:Red cristalina}
	\end{subfigure}
	\hspace{2cm}
	\begin{subfigure}[h]{0.3\textwidth}
		\includegraphics{Standalone/T1_03_Red_cristalina_sin_e.pdf}
		\caption{La red está a $\unit[T>0]{K}$ y hay un hueco en la CV.}
		\label{fig: Red cristalina sin e}
	\end{subfigure}
	\caption{Representación bidimensional de las bandas de valencia de una red de Si.}
\end{figure}
\subsection{\textcolor{red}{Dinámica del electrón en el sólido}} \label{sec: Dinamica del electron en el solido}

Como hemos visto en la ecuación \ref{eq: Schrodinger y Bloch}, las funciones de Bloch son solución a dicha ecuación diferencial. Ahora vamos a estudiar la dinámica del electrón en presencia de otras fuerzas, para ello utilizaremos Ecuación de Schrödinger Dependiente del Tiempo
\begin{equation}
    (H_i+U(\mathbf{r}))\psi(\mathbf{r},t)=-i\hbar \pdiff[]{\psi(\mathrm{r},t)}{t}
    \label{eq: ESDT con fuerza externas}
\end{equation}
Vamos a utilizar la conocida como \textbf{aproximación semiclásica}
\begin{itemize}
    \item Las funciones de Bloch, al formar una base del Espacio de Hilbert, nos permiten expresar la solución de la ecuación \ref{eq: ESDT con fuerza externas} como $$\psi(\mathbf{r},t)=\sum_n \sum_\mathbf{k} \beta (n,\mathbf{k},t) \, \psi_{n\mathbf{k}} (\mathbf{r}) $$
    \item Supongamos que la fuerza externa no es lo suficientemente fuerte como para que el electrón pueda cambiar de banda.
    $$\psi(\mathbf{r},t)=\sum_\mathbf{k} \beta (\mathbf{k},t) \, \psi_{n\mathbf{k}} (\mathbf{r}) $$
    \item Si además suopnemos que la dependencia temporal para las perturbaciones débiles se encuentra de manera implícita en \textbf{k} llegamos a $$\hbar \pdiff[]{\mathbf{k}}{t}=\mathbf{F}=-\nabla U$$
\end{itemize}

A partir de dicha ecuación podemos calcular los valores esperados de \textbf{v}, \textbf{k} y \textbf{r}.
\subsubsection{\textcolor{red}{Masa efectiva del electrón. Relación con las bandas de energía entorno a sus extremos.}} \label{sec: Masa efectiva}
Siguiendo con la aproximación semiclásica, utilicemos la Segunda Ley de Newton
\begin{equation}
    \mathbf{F}=m \diff[]{\mathbf{v}}{t} \rightarrow \langle \mathbf{v}\rangle=\frac{1}{\hbar} \nabla_\mathbf{k} E_{n \mathbf{k}} \rightarrow{}\diff[]{}{t} \left[\frac{1}{\hbar}  \pdiff[]{E}{k_i} \right]_i=\frac{1}{\hbar} \sum_j \frac{\partial ^2 E}{\partial k_i \partial k_j} \, \frac{\partial k_j}{\partial t}=\frac{1}{\hbar^2} \sum_j \frac{\partial ^2 E}{\partial k_i \partial k_j} \, F_j
\end{equation}
El último término, lo identificamos como un tensor de orden 2, y definimos el tensor de masa efectiva inversa como
\begin{equation}
    \left(\frac{1}{m^*}\right)_{i,j}=\frac{1}{\hbar^2} \frac{\partial ^2 E}{\partial k_i \partial k_j} 
    \label{eq: masa efectiva}
\end{equation}
Es decir, de manera general, los electrones se comportn como partículas libres pero su masa depende del estado en el que se encuentran. Físicamente lo interpretamos como la influencia del potencial periódico del cristal en el movimiento de los portadores.
Además si el tensor es diagonal y sus componentes son iguales se convierte en un escalar.\\

La ecuación \ref{eq: Forma bandas} podemos reescribirla 
\begin{equation}
	E_n(\mathbf{k})=	E_n(\mathbf{k_0})\pm\frac{\hbar^2}{2m_0} \sum_{i, \ j} \alpha_{ij} (k_i-k_{i0})(k_j-k_{j0}) 
 \label{eq: Forma bandas 2}
\end{equation}
El signo $\pm$ se utiliza si nos encontramos en el mínimo de BC o en el máximo de BV respectivamente.\\\\
Donde el nuevo coeficiente $\alpha_{ij}$ lo definimos como el tensor masa efectiva inversa reducida$$\alpha_{ij}=\left |\frac{m_0}{\hbar} \left[  \frac{\partial ^2 E}{\partial k_i \partial k_j}\right]\right|= m_0 \left|\left(\frac{1}{m^*}\right)_{ij}\right| $$
Que es un tensor simétrico, por lo tanto si escogemos los ejes principales, el tensor se vuelve diagonal y si definimos $m_{ii}^*=m_0\, (a_{ii})^{-1}$, la ecuación \ref{eq: Forma bandas 2}, tiene la forma:
\begin{equation}
    E(\mathbf{k})=\sum_i \pm \frac{\hbar^2 \mathbf{k}^2}{2m_{ii}^*}
\end{equation}
Es decir, que la energía tiene la misma forma que la partícula libre con $\mathbf{p}=\hbar\, \mathbf{k}$. Si las bandas son esféricas, la masa se vuelve un escalar y la solución para la energía a la ecuación \ref{eq: ESDT con fuerza externas} es la de partícula libre.
\newpage
\subsection{\textcolor{red}{Concepto de Hueco}} \label{sec: Concepto de Hueco}
Volvamos a fijarnos en la Figura \ref{fig: Red cristalina sin e}, y en lo comentado en la sección 1.2.1, este hueco lo podemos modelar utilizando mecánica newtoniana. Supongamos que aplicamos un campo eléctrico en el semiconductor, este campo eléctrico genera una corriente en la banda de valencia $$J=-e \sum_{\mathrm{ocupados}} n_iv_i=-e\sum_{\mathrm{todos}} n_j v_j +e\sum_{\mathrm{vacios}} n_k v_k$$ Pero en la banda de valencia, la densidad de corriente debe ser \textbf{nula} por simetría de bandas\footnote{Para cada \textit{v} existe una \textit{-v}.}, por lo tanto obtenemos que los estados vacíos generan una corriente $$J=+e\sum_{\mathrm{vacios}} n_k v_k$$ Esta corriente se asocia a \textbf{cargas positivas} cuyo número es el de vacantes y tienen la misma velocidad que los electrones en los estados vacíos.\\\\
Aplicando la Segunda ley de Newton  vemos que ya tenemos signficado físico al signo negativo de la masa, la masa negativa $m_\mathrm{e}^*=-|m^*|$ que resulta de \ref{eq: masa efectiva}, se debe a que la aceleración que experimentan los electrones\footnote{Al tener carga negativa, se mueven en sentido contrario al campo.} y los huecos debe ser la misma, por lo que la dinámica de los huecos se estudiará como partículas de carga +\textit{e} y masa $|m^*|$, en cantidad igual a las vacantes; mientras que los electrones tendrán ambas negativas, en cantidad igual a los estados ocupados.\footnote{Esto siempre es referido al techo de la banda de valencia, en la banda de conducción, el electrón tiene masa positiva y no consideramos huecos.}

\section{\textcolor{red}{Semiconductores intrínsecos y extrínsecos}} \label{sec: Semiconductores intrinsecos y extrinsecos}
Una vez caracterizados los portadores de carga que existen en los semiconductores podemos diferenciar entre dos tipos de semiconductores en función de si presentan o no impurezas.
\begin{itemize}
    \item \textbf{Semiconductores intrínsecos}, son aquellos que son cristalográficamente perfectos ya que no presentan impurezas.
    \item \textbf{Semiconductores extrínsecos}, contienen impurezas para mejorar la conductividad. Dentro de estos semiconductores, según el tipo de impurezas podemos encontrar semiconductores Tipo N, con impurezas donadoras ($n>p$) y tipo P, con impurezas aceptoras ($n<p$)
\end{itemize}

 Un ejemplo de un semiconductor extrínseco tipo N es el Si (G IV) dopado con P (G V); al introducir algunos átomos de P dentro del cristal de Si, el 5º electrón del P se encuentra debilmente ligado por lo que la energía térmica es suficiente como para hacer que salte a la capa de conducción como muestra la figura \ref{fig: Impurezas N}, resaltar que cada impureza deja libre un electrón y fijo un ión. Definimos $N_\mathrm{D}=N_\mathrm{D}^++N_\mathrm{D}^0$, como la \textbf{densidad efectiva de donadores} $N_\mathrm{D}^+$ es la densidad de impurezas ionizadas y $N_\mathrm{D}^0$ la densidad de impurezas no ionizadas.\\\\
 Como ejemplo de semiconductor extrínseco tipo P, tenemos el Si (G IV) dopado con Al (G III), esto nos genera un hueco en el cristal, y por lo tanto nos contribuye a la conducción de manera análoga al semiconductor tipo N (figura \ref{fig: Impurezas P}), la única diferencia es que tratamos con huecos en lugar de electrones.
 Análogamente definimos $N_\mathrm{A}=N_\mathrm{A}^-+N_\mathrm{A}^0$, como la \textbf{densidad efectiva de aceptadores} $N_\mathrm{A}^+.$ es la densidad de impurezas ionizadas y $N_\mathrm{A}^0$ la densidad de impurezas no ionizadas.
 \begin{figure}[h!]
    
	\centering
    \begin{subfigure}[h]{0.3\textwidth}
	\includegraphics{Standalone/T1_04_Impurezas_N.pdf}
	\caption{Efecto de la ionización en una impureza donadora.}
        \label{fig: Impurezas N}
    \end{subfigure}
    \hspace{2cm}
    \begin{subfigure}[h]{0.3\textwidth}
        	\includegraphics{Standalone/T1_05_Impurezas_P.pdf}
        \caption{Efecto de la ionización en una impureza aceptadora.}
        \label{fig: Impurezas P}
    \end{subfigure}
        \caption{Efecto de las impurezas en las bandas del cristal.}
        \label{fig: Impurezas}
\end{figure}
\subsubsection{\textcolor{red}{Densidad intrínseca de portadores}} \label{sec: Densidad intrinseca de portadores}
En situaciones reales, el semiconductor estará a una cierta temperatura $T$ y por lo tanto tendrá una contribución de energía térmica $E_\mathrm{th}=\frac{3}{2}k_\mathrm{B}T$, esta energía puede ser lo suficientemente alta como para hacer que algunos electrones que se encuentran en la banda de valencia salten a la banda de conducción y por lo tanto podamos producir corriente eléctrica. Por lo tanto definimos la \textbf{densidad intrínseca} de portadores como el número de electrones que se encuentran en la capa de conducción únicamente generados por la energía térmica. Para semiconductores intrínsecos $n=p=n_i$ mientras que para semiconductores extrínsecos se rige por $n\,p=n_i^2$ \footnote{Las definiciones matemáticas de densidades, así como la ecuación $n\,p=n_i^2$ lo desarrollamos cuando estudiemos las distribuciones de estados.} \newpage
\subsection{\textcolor{red}{Caracterización energética de los estados.}} \label{sec: Caracterizacion g de los estados}
Una vez introducidos los dos tipos de semiconductores que existen vamos entonces a caracterizar matemáticamente y a observar cómo se distribuyen los estados en la red cristalina.\\\\
Para ello vamos a modelar los electrones como si estuvieran confinados en un pozo cuadrado infinito tridimensional.
La densidad de estados en el recíproco es $g(\mathbf{k})=(4\pi^3)^{-1}$ mientras que en energía será $$g(E)=\frac{1}{4\pi^3} \int_{S(E)} \frac{\mathrm{d}S}{|\nabla_\mathbf{k} E|}$$ Donde $S(E)$ son superficies isoenergéticas esféricas. Además conocemos la forma que tiene la energía dentro del pozo
\begin{equation}
    E(\mathbf{k})=E_0\pm \frac{\hbar^2\, \mathbf{k}^2}{2m} \Rightarrow g(E)=\frac{4\pi}
{h^3} (2m^*)^{3/2}(\pm E \mp E_0)^{1/2}
\label{eq: Densidad de estados energetica}
\end{equation}
El $\pm$ de la energía depende de si nos encontramos en el fondo de la banda de conducción (+) o en el techo de la banda de valencia (-). Este modelo se conoce como \textbf{Estados de Bloch} y corresponde a semiconductores intrínsecos ya que son los estados generados por la periodicidad del cristal (ya lo vimos en 1.1.2).\\\\
Por otro lado tenemos si introducimos \textbf{impurezas}, por ejemplo fosforo en un cristal de silicio, debemos utilizar la ESIT para poder modelarlo, aquí utilizarmemos que la constante dieléctrica cambia (por aproximación de partícula única)
\begin{equation}
    \frac{-\hbar^2}{2m^*}\nabla^2-\frac{e^2}{4\pi \epsilon_\mathrm{SC}r}\psi(\textbf{r})=E\psi(\textbf{r})
\end{equation}
Que tiene la forma de un átomo de hidrógeno, por lo que obtenemos que la energía $$E_n=\frac{13.6}{\epsilon_r^2} \frac{m^*}{m_0} \frac{1}{n^2} \ \unit{eV}$$ es aquella que liga el electrón al átomo de la impureza. Por lo que encontramos que estados discretos de energía. Aunque en la práctica se comportan como un solo nivel (estado fundamental) $E_1=E_d$.
\subsubsection{Resumen}

\begin{table}[h!]
    \centering
      \extrarowheight = -0.5ex
\renewcommand{\arraystretch}{1.43}
    \begin{tabular}{|c|c|}
    \hline
        Estados de Bloch & Estados de impurezas hidrogenoides \\ \hline \hline
        No están localizados & Localizados espacialmente \\
        \small{(Generados por la periodicidad del cristal)} & \small{(Localizados con la impureza)}\\ \hline
        Continuo de estados de energía&  Estados discretos de energía\\\hline
         Ocupación independiente& Ocupación dependiente \\\hline
    \end{tabular}
    \label{tab:Comparacion bloch impurezas 1}
\end{table}
\newpage
\subsection{\textcolor{red}{Estadística de ocupación}} \label{sec: Estadistica de ocupacion}

\begin{wrapfigure}{l}{0.4\textwidth}
    \centering
\includegraphics{Standalone/T1_06_representacion_fd.pdf}
 \caption{Representación de la distribución de Fermi-Dirac. El eje de abcisas tiene como unidad $E_\mathrm{F}$.}
    \label{fig: representacion fd}
\end{wrapfigure}
\noindent Las partículas estudiadas son electrones, es decir, \textit{fermiones} por lo que se rigen por la \textbf{estadística de Fermi-Dirac}, para los \textbf{estados de Bloch} esta distribución en energías
\begin{equation}
    f(E)=\frac{1}{1+\exp\left( \frac{E-E_\mathrm{F}}{k_\mathrm{B}T}\right)}
    \label{eq: Distribucion FD}
\end{equation}
Podemos estudiar distintas situaciones en función de la temperatura fijando la energía. Si hacemos $T=0\unit{K}$ observamos que la gráfica tiene la forma punteada mientras que si la energía coincide exactamente con la del nivel de Fermi, y la temperatura está por encima de 0, la distribución toma valor $f(E)=0.5$. \newline
\\Destacar si la energía de Fermi se encuentra dentro del GAP, no tiene sentido físico y es sólo un parámetro en función del cual se expresa la ocupación.
Por último podemos expresarla probabilidad de \textbf{no ocupación} como
\begin{equation}
    1-f(E)=\frac{1}{1+\exp\left( \frac{-E+E_\mathrm{F}}{k_\mathrm{B}T}\right)}
    \label{eq: no ocupacion FD}
\end{equation}

\begin{tcolorbox}[colback=airforceblue!20, colframe=airforceblue, title=Ejemplo]
Supongamos que tenemos una diferencia entre energías de nivel y Fermi $E-E_\mathrm{F}=3k_\mathrm{B}T$, con $T=\unit[300]{K}$.\\\\
La ecuación \ref{eq: Distribucion FD}, toma valor $f(E)=0.0474= 4.74 \%$
\\\\
\textbf{\textcolor{airforceblue}{Comentario}} Para energías superiores a $E_\mathrm{F}$ la probabilidad de que se produzca un salto entre niveles es residual por lo tanto la razón de electrones que ocupan los estados disponibles es residual.
\end{tcolorbox}
\noindent El ejemplo anterior nos permite ``fijar'' un límite para pasar de la estadística de Fermi-Dirac, que es más compleja de utlizar matemáticamente; a la estadística de Maxwell-Boltzmann
\begin{equation}
    f(E)\approx \exp \left[\frac{E_\mathrm{F}-E}{k_\mathrm{B}T}\right]
\end{equation}
Esta distribución es una distribución clásica (rige la cinética de los gases clásicos). 
\subsubsection{\textcolor{red}{Nivel de Fermi}} \label{sec: Nivel de Fermi}
\noindent Definimos la energía de Fermi $E_\mathrm{F}$ como la energía a la que se encuentra el electrón más energético a una temperatura $T=\unit[0]{K}$.
\newpage
Ahora bien, no solo tenemos estados de Bloch en nuestros SC sino que también tenemos \textbf{impurezas}, podemos estudiar las densidades de las impurezas en función de las distribuciones de energía.

\subsubsection{Donadoras} \label{sec: Impurezas dondadoras}
\begin{itemize}
    \item Mediante la distribución de FD (Ec\ref{eq: Distribucion FD}), vamos a determinar las concentraciones de las impurezas neutras y las ionizadas, como la energía de las impurezas nomrmalmente se comporta como un solo nivel, las impurezas no ionizadas siguen la función de probabilidad
    \begin{equation}
        \frac{N_\mathrm{D}^0}{N_\mathrm{D}}=\frac{1}{1+\frac{1}{g_d}\exp\left( \frac{E_\mathrm{D}-E_\mathrm{F}}{k_\mathrm{B}T}\right)}
        \label{eq: probabilidad no ionizadas d}
    \end{equation}
    Donde $g_d$ es el nivel de degeneración. Para el caso de las impurezas donadoras ionizadas tenemos
    \begin{equation}
          \frac{N_\mathrm{D}^{+}}{N_\mathrm{D}}=\frac{N_\mathrm{D}-N_\mathrm{D}^0}{N_\mathrm{D}}=\frac{\frac{1}{g_d} \exp \frac{E_\mathrm{D}-E_\mathrm{F}}{k_\mathrm{B} T}}{1+\frac{1}{g_d} \exp \frac{E_\mathrm{D}-E_\mathrm{F}}{k_\mathrm{B} T}}=\frac{1}{1+g_d\exp{\left( \frac{E_\mathrm{F}-E_\mathrm{D}}{k_\mathrm{B}T}\right)}}
          \label{eq: probabiliad ionizadas d}
    \end{equation}
    Por último podemos establecer la relación que hay entre impurezas ionizadas y no ionizadas
    \begin{equation}
        \frac{N_\mathrm{D}^+}{N_\mathrm{D}^0}=\frac{1}{g_d} \exp\frac{E_\mathrm{D}-E_\mathrm{F}}{k_\mathrm{B}T}
        \label{eq: Razon N+N0}
    \end{equation}
\end{itemize}

\begin{wrapfigure}{r}{0.35\textwidth}
	\centering
    \begin{subfigure}[h]{0.3\textwidth}
	\includegraphics{Standalone/T1_07_Niveles_N.pdf}
	\caption{Distribución de niveles de energía en SCE tipo N.}
        \label{fig: Niveles  N}
    \end{subfigure}
    \begin{subfigure}[h]{0.3\textwidth}
        	\includegraphics{Standalone/T1_08_Niveles_P.pdf}
        \caption{Distribución de niveles de energía en SCE tipo P.}
        \label{fig: Niveles P}
    \end{subfigure}
        \caption{Niveles de energía de Fermi en las bandas del cristal a $T=\unit[0]{K}$.}
        \label{fig: Impurezas}
\end{wrapfigure}


\subsubsection{Aceptadoras} \label{sec: Aceptadoras}
\begin{itemize}
    \item Análogamente para las impurezas aceptadoras, las expresiones recogen una forma similar.
    Para las impurezas no ionizadas
    \begin{equation}
        \frac{N_\mathrm{A}^0}{N_\mathrm{A}}=\frac{1}{1+\frac{1}{g_d}\exp\left( \frac{E_\mathrm{F}-E_\mathrm{A}}{k_\mathrm{B}T}\right)}
        \label{eq: probabilidad no ionizadas A}
    \end{equation}
    Para las impurezas ionizadas
    \begin{equation}
           \frac{N_\mathrm{A}^{-}}{N_A}=\frac{N_\mathrm{A}-N_\mathrm{A}^0}{N_\mathrm{A}}=\frac{\frac{1}{g_d} \exp \frac{E_\mathrm{F}-E_\mathrm{A}}{K_\mathrm{B} T}}{1+\frac{1}{g_d} \exp \frac{E_\mathrm{F}-E_\mathrm{A}}{K_\mathrm{B} T}}=\frac{1}{1+g_d\exp{\left( \frac{E_\mathrm{A}-E_\mathrm{F}}{k_\mathrm{B}T}\right)}}
          \label{eq: probabiliad ionizadas A}
    \end{equation}
    Y la razón entre ionizadas y no ionizadas
    \begin{equation}
        \frac{N_\mathrm{A}^-}{N_\mathrm{A}^0}=\frac{1}{g_d} \exp\frac{E_\mathrm{F}-E_\mathrm{A}}{k_\mathrm{B}T}
        \label{eq: Razon N-N0}
    \end{equation}
\end{itemize}
Debido a que contiene exponenciales, podemos definir la \textbf{ionización total de impurezas} como aquellas en los que los denominadores de las exponenciales son  $$E_\mathrm{D}-E_\mathrm{F}\gg k_\mathrm{B}T \hspace{2cm} E_\mathrm{F}-E_\mathrm{A}\gg k_\mathrm{B}T$$
\clearpage



Esto también nos permite determinar la posición del nivel de Fermi
\begin{itemize}
    \item Para las impurezas donadoras, $E_\mathrm{F}$ se encuenta por encima de $E_\mathrm{D}$.
    \item Para las impurezas aceptadoras $E_\mathrm{F}$ se encuenta por debajo de $E_\mathrm{A}$.
\end{itemize}
\section{Densidades de portadores libres en equilibrio} \label{sec: Densidaddes de portadores en equilibrio}
Ya hemos conseguido determinar $N_\mathrm{D}^0(E_\mathrm{F})$ y $N_\mathrm{A}^0(E_\mathrm{F})$, ahora vamos a determinar las \textbf{densidades en la banda de conducción} $n(E_\mathrm{F})$ y $p(E_\mathrm{F})$.
\subsubsection{Densidad de electrones en la banda de conducción, \textit{n}} \label{sec: Densidad de electrones en la banda de conduccion}
Si utilizamos la ecuación \ref{eq: Densidad de estados energetica}, y la distribuicón de FD podemos determinar el número de electrones en la banda de conducción, recordar que en este caso cogemos el + de la ecuación.
\begin{equation}
    n=\int_{E_\mathrm{C}}^{\infty}g(E)f(E)=\int_{E_\mathrm{C}}^{\infty} \frac{4\pi}
{h^3} (2m^*)^{3/2}\frac{(E- E_0)^{1/2}}{1+\exp\left( \frac{E-E_\mathrm{F}}{k_\mathrm{B}T}\right)}=N_\mathrm{c} \frac{2}{\sqrt{\pi}} F_{1/2} \left(\frac{E_\mathrm{F}-E_\mathrm{C}}{k_\mathrm{B}T}\right)
\label{eq: Densidad de portadores}
\end{equation}
Con $N_\mathrm{c}=2\left(\frac{2\pi m^* k_\mathrm{B} T}{h^2}\right)^{3/2}$ la densidad equivalente de estados en la banda de conducción y $F_{1/2}$ la intergral de Fermi-Dirac de orden $1/2$
Gracias a que la densidad de portadores libres incluye la integral de Fermi Dirac podemos determinar algunas propiedades de la BC
\begin{equation}
    F_{1/2}=\int_0^\infty \frac{\eta ^{1/2}}{1+\exp{\eta-\eta_\mathrm{F}}} \mathrm{d}\eta \text{   con } \eta_\mathrm{F}=\frac{E_\mathrm{F}-E_\mathrm{C}}{k_\mathrm{B}T}
\end{equation}
\begin{itemize}
    \item Si $\eta_\mathrm{F}<-3$ la integral tiene la forma $F_{1/2}=\frac{\sqrt{\pi}}{2} \exp{\eta_\mathrm{F}}$ y entonces $E_\mathrm{C}-E_\mathrm{F}>3k_\mathrm{B}T \rightarrow E-E_\mathrm{F}>3k_\mathrm{B}T$ para todas las energías en la banda de conducción y como hemos dicho en la sección 1.3.2 es posible aproximar la estadísitca de FD a la de MB.
\end{itemize}
Si aplicamos a la ecuación \ref{eq: Densidad de portadores} la condición anterior, obtenemos
\begin{equation}
    n=N_\mathrm{c} \exp{\frac{E_\mathrm{F}-E_\mathrm{C}}{k_\mathrm{B}T}}
    \label{eq: densidad de portadores MB}
\end{equation}
\newpage
\subsubsection{Densidad de huecos en la banda de valencia, \textit{p}} \label{sec: Densidad de huecos enn la banda de valencia}
Seguiremos considerando superficies isoenergéticas esféricas pero esta vez utilizaremos distribución de probabilidad de \textbf{no ocupación} dada por la equación \ref{eq: no ocupacion FD}, la densidad de huecos tiene la forma
\begin{equation}
p=\int_{-\infty}^{E-\mathrm{V}}(1-g(E))f(E)=\int_{-\infty}^{E-\mathrm{V}} \frac{4\pi}
{h^3} (2m^*)^{3/2}\frac{(E- E_0)^{1/2}}{1+\exp\left( \frac{E_\mathrm{F}-E}{k_\mathrm{B}T}\right)}=N_\mathrm{v} \frac{2}{\sqrt{\pi}} F_{1/2} \left(\frac{E_\mathrm{V}-E_\mathrm{F}}{k_\mathrm{B}T}\right)
\end{equation}
\begin{itemize}
    \item En este caso si el semiconductor es no degenerado tenemos que si $E_\mathrm{V}-E_\mathrm{F}<-3k_\mathrm{B}T$ podemos hacer la aproximación que hemos desarrollado anteriormente y podemos utilizar la estadística de MB en lugar de FD. La densidad equivalente de estados de la banda de valencia es $N_\mathrm{v}=2\left(\frac{2\pi m^* k_\mathrm{B} T}{h^2}\right)^{3/2}$ y por lo tanto
    \begin{equation}
           p=N_\mathrm{v} \exp{\frac{E_\mathrm{V}-E_\mathrm{F}}{k_\mathrm{B}T}} 
           \label{eq: densidad de huecos MB}
    \end{equation}
\end{itemize}

\subsection{\textcolor{red}{Densidad intrínseca de portadores}} \label{sec: Densidad intrinseca de portadores}
Una vez definidas las densidades de huecos y electrones en sus respectivas capas como funciones de la energía de Fermi, vamos a estudiar cómo se relacionan ambas cantidades, esto variará en función de si el semiconductor es intrínseco o extrínseco.Si definimos \begin{equation}
    n_\mathrm{i}^2=n\cdot p=N_\mathrm{v} N_\mathrm{c} \exp{\frac{E_\mathrm{V}-E_\mathrm{F}}{k_\mathrm{B}T}} \exp{\frac{E_\mathrm{F}-E_\mathrm{C}}{k_\mathrm{B}T}}=N_\mathrm{v} N_\mathrm{c} \exp{\frac{E_\mathrm{V}-E_\mathrm{C}}{k_\mathrm{B}T}}
    \label{eq: Densidad intrinseca de portadores}
\end{equation}
Por lo tanto obtenemos $$n_\mathrm{i}=\sqrt{N_\mathrm{v} N_\mathrm{c}} \exp{\frac{-\mathrm{GAP}}{2k_\mathrm{B}T}}$$

\subsubsection{Semiconductor intrínseco}
\begin{wrapfigure}{r}{0.3\textwidth}
    \includegraphics{Standalone/T1_09_Gen_Rec.pdf}
	\caption{Proceso de generación-recombinación.}
        \label{fig: Gen-Rec}
    \end{wrapfigure}
\noindent En los semiconductores intrínsecos, cuando generamos portadores térmicamente se da un proceso que se conoce como \textbf{generación-recombinación} en el cual cuando un electrón deja un hueco en la capa de valencia, es ocupado por otro electrón, a su vez ese electrón deja un hueco en esa posición que puede ser ocupado por un electrón de la banda de conducción, como vemos en la figura.\newline

\noindent Volver a destacar que el nivel de Fermi se encuentra cerca del MIDGAP.

\newpage
\subsubsection{Semiconductor extrínseco}
\noindent Supongamos que el semiconductor se encuentra en equilibrio, para el caso de los semiconductores tipo N $n \uparrow \Rightarrow{} p \downarrow$ mientras que para los de tipo P $p \uparrow \Rightarrow{} n \downarrow$. Lo que vamos a desarrollar ahora será con un semiconductor tipo N pero es análogo a un tipo P.

\begin{figure}[h]
    \centering
  \includegraphics{Standalone/T1_10_Gen_recom_tipo_N.pdf}
    \caption{Generación-recombinación en un semiconductor intrínseco tipo N.}
    \label{fig:Gen-recom tipo N}
\end{figure}

\noindent Vemos en la figura \ref{fig:Gen-recom tipo N} que en este caso algunos de los huecos que existen en la banda de las impurezas, generan huecos en la banda de valencia, aunque algunos de los electrones destruyen algunos huecos intrínsecos por lo tanto vemos que la densidad de estados equivalente debe tener la forma $N_\mathrm{D}^+=n-p$, si utilizamos la ecuación \ref{eq: Densidad de portadores} queda como $$p+N_\mathrm{D}^+=n=N_\mathrm{c} \exp{\frac{E_\mathrm{F}-E_\mathrm{C}}{k_\mathrm{B}T}}$$
Por lo tanto cuanto más cerca esté $E_\mathrm{F}$ de la banda de conducción mayor será \textit{n}.
\subsubsection{Semiconductores degenerados y no degenerados. Localización del nivel de Fermi}
\noindent A modo de resumen vamos a interpretar físicamente la localización del nivel de Fermi en las bandas de energía.\\


\indent   \textbf{Semiconductores no degenerados}
\begin{itemize}
    \item Si en la banda de conducción hay pocos estados ocupados, es decir existen menos electrones libres, el nivel de Fermi se encuentra cerca del MIDGAP sólo que desplazado hacia la banda de conducción. Físicamente lo podemos interpretar como $N_\mathrm{D}>N_\mathrm{V}$.
    \item Si en la banda de valencia hay pocos estados ocupados, es decir existen menos huecos, el nivel de Fermi se encuentra cerca del MIDGAP sólo que desplazado hacia la banda de valencia. Físicamente lo podemos interpretar como $N_\mathrm{A}>N_\mathrm{V}$
\end{itemize}
\indent \ \ \  \textbf{Semiconductores degenerados}
\begin{itemize}
    \item En este caso el nivel de Fermi puede estar dentro de las bandas de conducción o de valencia, esto se debe a que existe un gran número de electrones o huecos respectivamente.
    \item Si el nivel de Fermi se localiza dentro de las bandas, lo interpretaremos como que todos los estados están ocupados/vacíos hasta el nivel de Fermi.
\end{itemize} Por lo que podemos concluir que \textbf{el nivel de Fermi nos indica el tipo de semiconductor}.
 \newpage
 \subsection{Densidad de portadores y nivel de Fermi en función de la Temperatura.} \label{sec: Densidad de portadores y nivel de fermi en funcion de la temperatura}
Como hemos visto, todas las ecuaciones referidas a las densidades de portadores presentan dependencia con la temperatura. Dicha dependencia nos permite estudiar los tipos de semiconductores si son intrínsecos o no o si son tipo N o tipo P y los regímenes en los que nos movemos.
\subsubsection{Semiconductor intrínseco}
\noindent Los semiconductores intrínsecos como hemos visto con anterioridad en los SCI la densidad intrínseca de portadores tiene la forma
\begin{equation}
    n_\mathrm{i}=\sqrt{N_\mathrm{v} N_\mathrm{c}} \exp{\frac{E_\mathrm{V}-E_\mathrm{C}}{k_\mathrm{B}T}} \xrightarrow{\text{Despejando}} E_\mathrm{Fi}-E_\mathrm{C}=\frac{E_\mathrm{V}-E_\mathrm{C}}{2}+\frac{k_\mathrm{B}T}{2} \ln{\frac{N_\mathrm{V}}{N_\mathrm{C}}}
    \label{eq: Nivel de Fermi intrínseco}
\end{equation}
Si las masas de los portadores y los huecos son iguales entonces el \textbf{nivel de Fermi intrínseco} se encontrará en el centro del GAP.
\subsubsection{Semiconductor extrínseco}
\begin{wrapfigure}{r}{0.4\textwidth}
    \centering
    \includegraphics{Standalone/T1_12_Concentracion_electrones_frente_a_T.pdf}
    \caption{Concentración de electrones frente a la temperatura.}
    \label{fig:Representacion densidad frente temperatura}
\end{wrapfigure}
\noindent Si utilizamos las diferentes ecuaciones \ref{eq: Razon N-N0} y \ref{eq: Razon N+N0} podemos construir lo que se conoce como \textbf{condición de neutralidad}:
\begin{equation}
    p+N_\mathrm{D}^+=n+N_\mathrm{A}^-
    \label{eq: condicion de neutralidad}
\end{equation}
que además si utilizamos la ecuación \ref{eq: Densidad intrinseca de portadores} podemos determinar lo que se conoce como \textbf{ionización total de impurezas} combinando ambas ecuaciones obtenemos
\begin{equation}
        n=\frac{N_\mathrm{D}-N_\mathrm{A}}{2}+ \sqrt{\left(\frac{N_\mathrm{D}-N_\mathrm{A}}{2} \right)+n_\mathrm{i}^2}
        \label{eq: SCE ionizacion total}
\end{equation}
Y tenemos un criterio para determinar si los minoritarios son despreciables frente a las impurezas. 
\begin{equation}
  n=p+N_\mathrm{D}^+-+N_\mathrm{A}^- \Rightarrow n>N_\mathrm{D}-N_\mathrm{A} \Rightarrow p=\frac{n_\mathrm{i}^2}{n}< \frac{n_\mathrm{i}^2}{N_\mathrm{D}-N_\mathrm{A}}
\end{equation}
\subsubsection{Concentración de portadores en función de la temperatura}
\noindent Si analizamos las ecuaciones \ref{eq: condicion de neutralidad} y \ref{eq: SCE ionizacion total}, las expresamos como funciones de la temperatura y las representamos en la Figura \ref{fig:Representacion densidad frente temperatura}, la región extrínseca también se conoce como región de agotamiento. Vemos que cuando se alcanza la llanura, consideramos que todas las temperaturas están ionizadas. Como la densidad intrínseca crece como $n_\mathrm{i}\propto T^{3/2}$ llegará un punto en el que comenzará a dominar sobre el resto de portadores. Las regiones estan gobernadas por $$\text{Ioniz Par:} \ \boxed{ n=N_\mathrm{D}^+ \ \ p=\frac{n_\mathrm{i}^2}{N_\mathrm{D}^+}} \ \hfill \text{Extr:} \ \boxed{n=N_\mathrm{D} \ \ p=\frac{n_\mathrm{i}^2}{N_\mathrm{D}}} \ \hfill \text{Intr:} \ \boxed{n=p=n_\mathrm{i}=\sqrt{N_\mathrm{v} N_\mathrm{c}} \exp{\frac{E_\mathrm{V}-E_\mathrm{C}}{2k_\mathrm{B}T}}}$$
\newpage
\subsubsection{\textcolor{red}{Nivel de Fermi en función de la temperatura}}
\noindent Si vemos las ecuaciones \ref{eq: Nivel de Fermi intrínseco} (cambiando los términos de conducción por los de valencia y obtenemos dos ecuaciones), estas ecuaciones como vemos son ecuaciones logarítimicas en la concentración y nos permite establecer las separaciones que tiene el nivel de Fermi en función del tipo de Semiconductor (N ó P) y la temperatura.
\begin{figure}[h]
\begin{subfigure}[h]{0.3\textwidth}
    \includegraphics{Standalone/T1_13_Nivel_de_Fermi_frente_a_T.pdf}
    \caption{Representación del nivel de Fermi frente a T.}
\end{subfigure}
\hspace{3cm}
\begin{subfigure}[h]{0.3\textwidth}
    \includegraphics{Standalone/T1_14_Nivel_de_Fermi_frente_a_N.pdf}
    \caption{Representación del nivel de Fermi frente a n}
\end{subfigure}
\end{figure}
 \section{Propiedades ópticas de los semiconductores} \label{sec: propiedades opticas de los semiconductores}
 \begin{wrapfigure}{r}{0.4\textwidth}
    \centering
    \includegraphics{Standalone/T1_15_Espectro_electromagnetico.pdf}
    \caption{Espectro de electromagnético frente a long. de onda.}
    \label{fig:Espectro electromagnetico}
\end{wrapfigure}
 Como sabemos el efecto fotoeléctrico nos permite generar electrones libres mediante la irradiación de ondas electromagnéticas en la superficie de un metal. La energía de los fotones que inciden en el metal es $$E=h\nu=h\frac{c}{\lambda}$$ además, recordar que por cada electrón generado, se creará un hueco en la BV.

\subsubsection{Espectro electromagnético y GAP de los principales semiconductores}
\noindent Como el Si y el GaAs tienen que ser iluminados con luz infraroja, el GaAs$_{1-y}$P$_{y}$ conocido es un semiconductor ternario, admite un amplio rango del espectro; para el resto de semiconductores podemos introducir indio (In) para reducir el tamaño del GAP.
\clearpage
\subsubsection{Interacción radiación-semiconductor}
\noindent La absorción de radiación puede darse por absorción de la red o de los portadores, lo que modifica sus propiedades eléctricas y de emisión de energía. La absorción produce los siguientes efectos:
\begin{itemize}
    \item Fotoconductividad, si el semiconductor es homogéneo.
    \item Efecto fotovoltaico, si el semiconductor es no homogéneo.
    \item Efecto fotoelectromagnético.
\end{itemize}

Otra forma de interactuar la radiación y el semiconductor es por recombinación, generando calor (fonones) o radiación óptica (\textbf{luminiscencia}). La recombinación estimulada por fotones nos permite también desarrollar láseres.
\subsection{Eventos ópticos clásicos. Reflexión, absorción y transmisión} \label{sec: eventos opticos clasicos}
Cuando incidimos con una onda electromagnética sobre la superficie de un metal, vemos que se producen 3 procesos Reflexión, Transmisión y Absorción $$\frac{I_\mathrm{R}}{I_0}+\frac{I_\mathrm{T}}{I_0}+\frac{I}{I_0}=1$$ para aplicaciones en semiconductores nos interesa que exista únicamente \textbf{absorción}.
\subsubsection{Reflexión y transmisión}
\noindent La reflexión reduce la eficiencia de los dispositivos basados en absorción, además hace que la sensibilidad dependa de la frecuencia (óptica I).
\\\\ \noindent Por otra parte la transmisión depende también de la frecuencia (ópitica I).
\subsubsection{Absorción de la radiación por un semiconductor. Mecansimos de absorción}
La absorción se rige por la \textbf{Ley de Bouger-Lambert} $$\mathrm{d} I(x)=-\alpha I(x) \mathrm{d}x \Rightarrow I(x)=I_0\exp{-\alpha\,x}$$Si la representamos, vemos que hay un incremento abrupto cuando $h\nu \geq E_\mathrm{g}$.

\subsubsection{Absorción intrínseca}
La absorción intrínseca se produce por una fotoexcitación de un electrón desde la banda de valencia hasta la banda de conducción. Si nos fijamos en la Figura \ref{fig: estructuras de bandas} los procesos de transición directa, son los más probables mientras que los de transición indirecta son los menos probables.\\\\
Consideraciones de diseño:
\begin{itemize}
    \item GAP directo: Absorción superficial, zona activa próxima a la superficie.
    \item GAP indirecto: Absorción en volumen, la zona activa amplia en volumen.
    \item Nos interesa introducir la \textbf{vida media} $\tau$, que se define como el tiempo medio que los portadores pasan en una banda antes de recombinarse.
    \item Otro término interesante es la \textbf{Longitud de difusión} $L_\mathrm{p}=\sqrt{D_\mathrm{p} \tau}$, como la longitud media que los portadores recorren antes de recombinarse.
\end{itemize}
\subsubsection{Absorción extrínseca}
Las impurezas absorben los fotones y los picos del espectro de absorción están entorno a la energía de ionización de las impurezas.
\subsubsection{Absorción por portadores libres}
Las transiciones de portadores entre niveles dentro de su propia banda, el continuo de niveles energéticos en la banda de conducción y en la banda de valencia generan un nivel de fondo en el espectro.
\subsubsection{Absorción excitónica}
Líneas en el espectro correspondients a niveles de excitación del excitón\footnote{Cuasipartícula o excitación elemental de los sólidos formada por un electrón y un hueco ligados a través de la interacción coulombiana.}.
\subsubsection{Absorción por vibraciones de la red}
Los fotones son absorbidos por parte de la red en forma de fonones, los picos en el espectro de absorción correspondientes a múltiplos enteros de la energía de los fonones.
\subsection{Fotoconductividad en semiconductores} \label{sec: fotoconductividad en semiconductores}
La fotoconductividad se desarrollará más en el siguiente tema, de manera introductoria, observamos un incremento de la conductividad originado por la absorción de radiación.
\begin{equation}
    \sigma_\mathrm{L}=\sigma_0 +\Delta \sigma \ \text{con } \sigma_0=q(n_0\mu_{n0}+p\mu_{p0})
\end{equation}
Para un semiconductor tipo N tenemos
\begin{equation*}
    \sigma_0=q n_0 \mu_{n0} \hspace{1cm} \sigma_0+ \Delta \sigma=q (n_0+\Delta n) (\mu_{n0}+\Delta \mu) \hspace{1cm} \Delta \sigma= q \mu_{n0} \Delta n+ q (n_0+\Delta n)\Delta \mu
\end{equation*}
Donde $\Delta n= G\tau $ son los portadores fotogenerados y $\Delta \mu $ es el estado de ionización de impurezas de portadores en sub-bandas con diferente $\mu$.\\\\
Si hacemos pasar una corriente por un semiconductor que esta siendo iluminado con una radiación capaz de generar electrones. $$I=A\sigma E= A\sigma_0 E+ A\Delta E=I_0+\Delta I$$ que combinando con lo que hemos desarrollado anteriormente, obtenemos $$\Delta I=Aq(n_0\mu_{n0}+p\mu_{p0}) G \tau \frac{V}{L}$$ 
\subsection{\textcolor{red}{Luminiscencia en semiconductores}} \label{sec: luminiscencia en semiconductores}
Este proceso se da por pérdida de energía de los electrones excitados, tenemos tres posibles transiciones de desexcitación: interbanda, por impureas o intrabanda.\\\\
Cuando se emiten electrones estos pueden saltar de una banda a otra, la más probable es la que salta una distancia de energías del GAP. Los semiconductores  más eficientes son intrínsecos y de GAP directo tales que el electrón salte entre bandas de energía y emita un fotón lo más energético posible.\\\\
Otro aspecto que nos interesa es que queremos que los semiconductores estén entre el grupo III y el V puesto que el GAP es mayor. Para conseguir entrar en el visible, lo que hacemos es dopar con otros semimetales como el Indio para reducir el tamaño del GAP.
\pagestyle{fancy}
\fancyhead[R]{Electrónica Física: Tema 2}
\fancyhead[L]{4º Grado en Física}
\chapter{Fenómenos de transporte de carga en Semiconductores}
Hasta ahora hemos estudiado semiconductores homogéneos en equilibrio termodinámico, caracterizados por $n_0\cdot p_0=n_\mathrm{i}^2$ e $I=0$. Si nos vamos a situaciones fuera del equilibrio (ya sea homogéneo o no), cuando haya una presencia de \textit{I}. Fuera del equilibrio, las concentraciones de portadores ya no son constantes, por lo tanto $$n_0 \xrightarrow{} n(\mathbf{r},t) \hspace{2cm} p_0 \xrightarrow{} p(\mathbf{r},t)$$.
\section{Corrientes de Arrastre y Difusión. Estructura de los niveles de Fermi en el no-equilibrio} \label{sec: Corrientes de difusión y arrastre}
\subsection{\textcolor{red}{Corriente de Arrastre}} \label{sec: Corriente de arrastre}
Supongamos que aplicamos un campo $\mathbf{E}=E_x \ \mathbf{\hat{\imath}}$ en la red cristalina. Dentro de la red, la concentración de portadores de carga $n$\footnote{Podríamos poner p} se mueve con una velocidad media $\langle v_x \rangle$, por lo tanto podemos definir una densidad de corriente de portadores y una intensidad de corriente
\begin{equation}
    J_x=qn\langle v_x \rangle \rightarrow{} I=J_xA
    \label{eq: Densidad de corriente de portadores}
\end{equation}
Si recordamos la sección 1.2.2, en ella vimos que podemos considerar el movimiento del electrón dentro de la red cristalina como si fuera una única partícula libre pero con masa $m^*$ sin embargo, las redes cristalinas pueden presentar \textbf{impurezas} o incluso por el mero hecho de estar a una cierta temperatura encontramos \textbf{fonones}, tanto las impurezas como los fonones interactúan con los electrones (en este caso tratados como partículas libres) lo cual hace que no puedan ser aceleradas indefinidamente debido al campo $E_x$, lo que se conocen com \textbf{mecanismos de scattering} que se dan cada tiempo medio $\tau$.\\\\
Por lo tanto el movimiento de los portadores serán una secuencia de recorridos de partículas libres y mecanismos de scattering que se propagan en la dirección del campo eléctrico.
\newpage
En un recorrido libre de duación t, la velocidad instantánea será 
\begin{equation}
    v_x=v_0+\frac{qE_xt}{m^*} \xrightarrow{\text{Promedio}} \langle v_x\rangle=v_\mathrm{d}=0+\frac{q \tau}{m^*} E_x  
    \label{eq: Velocidad de desplazamiento}
\end{equation}
en este caso por scattering $\langle v_\mathrm{d} \rangle=0$. El factor de proporcionalidad que existe entre la velocidad de arrastre $v_\mathrm{d}$ y el campo eléctrico lo definimos como \textbf{movilidad} $\unit{[cm^2 \, V^{-1} \, s^{-1}]}$. En términos generales, la movilidad es un tensor (aunque solo lo vamos a considerar un escalar)
\begin{equation}
   \mathbf{v}=\mu_{ij} \mathbf{E}=\left (\frac{q\tau}{m^*}\right)_ {ij} \mathbf{E}
    \label{eq: movilidad tensor}
\end{equation}
Donde $q=\pm e$ \footnote{El signo - no lo tenemos en cuenta dentro del escalar sino fuera, ya que representa el sentido en el que se mueven los electrones.}en función de si son huecos o electrones y etiquetaremos con \textit{p} o \textit{n} a la movilidad para referirnos a huecos o electrones, respectivamente A partir de ahora consideraremos que la corriente es solo en la dirección $x$ salvo que se indique lo contrario.\\\\
Vamos a estudiar los efectos que produce en la corriente, de las ecuaciones \ref{eq: Densidad de corriente de portadores} y \ref{eq: movilidad tensor} obtenemos
\begin{equation}
    J_{\mathrm{a}n x}=en\mu_n^{(xx)}E_x=\sigma_n^{(xx)}E_x \hspace{2cm}J_{\mathrm{a}p x}=en\mu_p^{(xx)}E_x=\sigma_p^{(xx)}E_x 
\end{equation}
Donde $\sigma_{n/p}$ es la \textbf{conductividad} $[\unit{(\Omega \, cm)^-1}]$ que de nuevo es un tensor por el carácter tensorial de la movilidad. Por lo que podemos definir la densidad de corriente de arrastre como  \footnote{En este caso da igual donde situemos los índices, únicamente están para indicar que son tensores, de aquí en adelante los omitiremos ya que los trataremos como escalares.}
\begin{equation}
    \mathbf{J}_\mathrm{a}=\mathbf{J}_{\mathrm{a}n}+\mathbf{J}_{\mathrm{a}p}=(\sigma_n^{ij}+\sigma_p^{ij})\textbf{E}=e(n\mu_n^{ij}+p\mu_p^{ij})\mathbf{E}=\sigma_{ij}\mathbf{E}
    \label{eq: conductividad de semiconductor}
\end{equation}
La \textbf{conductividad del semiconductor} $\sigma$ es el que nos indica si el semiconductor es buen conductor o no.
\subsubsection{\textcolor{red}{Ley de Ohm. Origen microscópico}}
Supongamos que tenemos una barra paralelepípeda, por la que hacemos pasar una corriente.
\begin{wrapfigure}{r}{0.4\textwidth}
    \centering
    \includegraphics{Standalone/T2_01_barra_de_semiconducor.pdf}
    \caption{Barra de un semiconductor atravesada como una corriente. Funciona como una resistencia}
    \label{fig: barra de semiconducor}
\end{wrapfigure}
\\ En las ecuaciones \ref{eq: movilidad tensor} y \ref{eq: conductividad de semiconductor} observamos relaciones lineales. Lo que conocemos \textbf{\textcolor{blue}{Ley de Ohm microscópica}} de donde podemos deducir la \textbf{\textcolor{red}{Ley de Ohm macroscópica}} (trabajamos en módulo)
\begin{equation}
\begin{aligned}
    &\textcolor{blue}{J=\frac{I}{W\,H}=\sigma E=\sigma \frac{V}{L}} \Rightarrow{}\\ &\textcolor{red}{V=\frac{1}{\sigma} \frac{L}{W\,H}I=\rho \frac{L}{W\,H}I=R\,I}
\end{aligned}
\end{equation}
Estas ecuaciones sólo y las ecuaciones \ref{eq: movilidad tensor} y \ref{eq: conductividad de semiconductor} (tratadas como escalares) solo son válidas para campos bajos.
\clearpage
\subsubsection{\textcolor{red}{Comparación entre los semiconductores de los grupos IV, III-V en un diagrama v-E}}
\noindent Una vez determinada las ecuaciones de la corriente tanto microscópicas como macroscópicas, para campos bajos, veamos cómo afecta el aumento del campo eléctrioco a la velocidad de la corriente en semiconductores intrínsecos.\\\\Supongamos que tenemos un cristal de Silicio (G IV) a una temperatura $T=\unit[300]{K}$, si le aplicamos un campo eléctrico, vemos que la velocidad aumenta de manera lineal, esto hace que la energía de los $e^-$ aumente por lo que aumentan los procesos de scattering dentro de la red, lo que hace que la energía se pierda en vibraciones (fonones). Que aumenten los procesos de scattering implica que el tiempo de relajación $\tau$ sea cada vez menor hasta que se queda estático la velocidad se satura y se mantiene constante como vemos en la relación de la ecuación \ref{eq: Velocidad de desplazamiento}.\\\\
Vamos a fijarnos también en un cristal GaAs (G III-V), vemos que exsite una movilidad diferencial negativa, es decir, $\mu =\partial_E\,v<0$ cuando pasamos del punto de saturación. En este caso los elecrones ganan energía y pasan a otros mínimos de la BC que tienen mayor masa, en la Figura \ref{fig: estructuras de bandas} (a) en lugar de estar en el mínimo central estará en uno de los mínimos de los lados, Es decir hay transiciones a mínimos superiores de la banda de conducción, que tienen mayor masa efectiva (Ecuación \ref{eq: masa efectiva}).

\begin{figure}[h]
    \centering
    \includegraphics{Standalone/T2_11_Velocidad_frente_a_campo.pdf}
    \caption{Velocidad de la corriente de arrastre frente al campo eléctrico para Si, Ge y GaAs hiperpuros.}
    \label{fig: Velocidad vs E}
\end{figure}
\newpage
\subsubsection{\textcolor{red}{Dependencia de la movilidad con la concentración de impurezas}}

\begin{wrapfigure}{l}{0.5\textwidth}
    \centering
    \includegraphics{Standalone/T2_12_Movilidad_frente_a_impurezas.pdf}
    \caption{Movilidad de los electrones y los huecos frente a la concentracion de impurezas para Ge, Si y GaAs. Con $T=\unit[300]{K}$ }
    \label{fig: movilidad vs impurezas}
\end{wrapfigure}
\noindent Como hemos visto con anterioridad, los portadores sufren dos mecanismos de scattering, el scattering generado por la red y por las impurezas, si los estudiamos por separados tenemos

\begin{itemize}
    \item \textbf{Scattering de red / Scattering de fonones}. Vamos a suponer que tenemos una red cristalina en la que sólo tenemos scattering de fonones. Como dichas interacciones son funciones de la temperatrua, podemos denotar la movilidad $\mu_\mathrm{F}$ como la movilidad que observamos debida a la interacción con los fonones. A primer orden (según la teoría de scattering)$$\mu_\mathrm{F} \propto T^{-3/2}$$ Por lo que vemos que esta movilidad aumenta conforme disminuye la temperatura. si nos extendemos a ordenes superiores el exponente sigue siendo negativo pero distinto de $\nicefrac{-3}{2}$. 
    \item \textbf{Scattering de impurezas ionizadas}. Del mismo modo que hemos modelizado para las interacciones de fonones, las interacciones con las impurezas ionizadas (interacciones coulombianas) tienen asociada una movilidad $\mu_\mathrm{I}$, a primer orden \begin{equation}
        \mu_\mathrm{I}\propto \frac{T^{3/2}}{N_\mathrm{D}^++N_\mathrm{A}^-}
    \end{equation} 
\end{itemize}
\noindent Por lo que si aumentamos la temperatura, aumentamos la velocidad de la corriente, lo que nos reduce el tiempo entre interacciones, esto hace que se reduzca el efecto de scattering y que sea un valor mayor al esperado. Si por otro lado aumentamos el número de impurezas, tendremos un mayor número de interacciones por lo que se reduce el valor de $\mu_I$.\\\\
Si denotamos como $\tau_i$ al tiempo entre interacciones (ya sean con fonones o impurezas), tenemos que $\mathrm{d}t/\tau_i$ es la probabilidad de que ocurra una de las interacciones. Por lo tanto definimos el tiempo de relajación $\tau$ como la probabilidad de que se de una interacción $$\frac{\mathrm{d}t}{\tau}=\frac{\mathrm{d}t}{\tau_\mathrm{I}}+\frac{\mathrm{d}t}{\tau_\mathrm{F}}$$ y de la ecuación \ref{eq: movilidad tensor}
\begin{equation}
    \frac{1}{\mu}=\frac{1}{\mu_\mathrm{I}}+\frac{1}{\mu_\mathrm{F}}
    \label{eq: Movilidad total}
\end{equation}
\subsubsection{\textcolor{red}{Dependencia de la movilidad con la termperatura}}
Si analizamos la ecuación \ref{eq: Movilidad total} debido al gran número de impurezas, la movilidad se mantiene casi invariante con \textit{T} por lo que mientras sólo tengamos términos de primer orden, encontraremos que hasta la región de saturación es casi constante por ello vemos la aproximación de la recta en la figura \ref{fig: Velocidad vs E}.
\subsection{\textcolor{red}{Corriente de difusión}} \label{sec: Corrientes de difusión}
Ya hemos estudiado qué ocurre si aplicamos un campo eléctrico en un semiconductor, ahora vamos a estudiar qué implica que el semiconductor sea \textbf{no homogéneo}, es decir que $n=n(\mathbf{r})$ y $p=p(\mathbf{r})$, una analogía clásica la encontramos en un volumen $V$ subdividido en dos volúmenes  $V/2$ por una membrana. En el que uno de ellos, tiene $n$ partículas en estado gaseoso y el otro está vacío, retiramos la membrana, observamos que hay un flujo de partículas hacia el lado vacío ocupando todo el volumen. En el caso de los semiconductores no homogéneos, tienden a crear flujos de partículas buscando que la concentración de las mismas sea homogénea.\\\\
Los fenómenos de difusión están modelizados por la \textbf{ley de Fick} 
\begin{equation}
    F_x=-D\diff[]{N}{x}  \xrightarrow{\text{En 3D}} \mathbf{F}=-D\nabla N
    \label{eq: difusion}
\end{equation}
Donde $\textbf{F}$ es el flujo, $D$ es el coeficiente de difusión y $N$ la concentración de partículas ($e^-$ o $h^+$), dicho flujo induce una corriente $\mathbf{J}=q\mathbf{F}$ que en función de si son electrones o huecos será:
\begin{equation}
    \mathbf{J}_{\mathrm{d}n}=eD_n\nabla n \hspace{2cm} \mathbf{J}_{\mathrm{d}p}=-eD_n\nabla p
    \label{eq: correinte de difusion}
\end{equation}
Donde los coeficientes de difusión están definidos por lo que se conoce como \textbf{relaciones de Einstein}.
\begin{equation}
    D_n=\frac{k_\mathrm{B}T}{e} \mu_n \hspace{2cm} D_p=\frac{k_\mathrm{B}T}{e} \mu_p
\end{equation}
Destacar que en este caso el signo - está cambiado frente a las ecuaciones \ref{eq: movilidad tensor} ya que la difusión tiene sentido contrario al gradiente (menor concentración a mayor concentración), este signo implica que si tenemos un mismo gradiente de concentración de electrones y huecos, se generará una corriente de difusión en el mismo sentido.\\\\
Si sumamos las contribuciones de la corriente de arrastre y la de difusión, tenemos que la corriente total será 
\begin{equation}
\boxed{
\mathbf{J}=\mathbf{J}_n+\mathbf{J}_p\left\{\begin{array}{l}
\mathbf{J}_n=\mathbf{J}_{\mathrm{a} n}+\mathbf{J}_{\mathrm{d} n}=e \mu_n n \mathbf{E}+e D_n {\nabla} n \\
\mathbf{J}_p=\mathbf{J}_{\mathrm{a} p}+\mathbf{J}_{\mathrm{d} p}=e \mu_p p \mathbf{E}-e D_p {\nabla} p
\end{array}\right.
} %
\label{eq: Densidad de corriente total}
\end{equation}
\newpage
\subsection{Nivel de Fermi. Relación con las Corrientes de Difusión y Arrastre} \label{sec: Nivel de Fermi. Relacion con las corrientes de difusion y arrastre}
 En el capítulo anterior hemos observado la relevancia del nivel de Fermi a la hora de determinar los distintos niveles de energía. Una de las suposiciones que hicimos es que el nivel de Fermi era constante frente a la posición, (debida a la homogeneidad del cristal).
\begin{tcolorbox}[colback=applegreen!20, colframe=applegreen, title=Demostración: En equlibrio termodinámico el nivel de Fermi es constante con la posición.]
Supongamos que en el semiconductor existe un potencial eléctrico $V(x)$, por lo que la energía de la banda de conducción es $E_\mathrm{C}=E_\mathrm{C0}-eV(x)$. El término de la energía eléctrica está superpuesto al término de la del Hamiltoniano del cristal.\\\\
En un semiconductor tipo N, los electrones se encuentran innicialmente en $E_\mathrm{F}=E_\mathrm{F}(x)$, de la ecuación \ref{eq: densidad de portadores MB} obtenemos
$$n(x)=N_\mathrm{c} \exp{\frac{E_\mathrm{F}(x)-E_\mathrm{C}(x)}{k_\mathrm{B}T}}$$ El campo eléctrico, supongamos que es electroestático, tendrá la forma$$E_x(x)=\frac{1}{e}\diff[]{E_\mathrm{C}(x)}{x}$$
De la ecuación \ref{eq: Densidad de corriente total} para los portadores $$J_n^{(x)}=e \mu_n n E_x+e D_n \diff[]{n}{x}=e\mu_n n\frac{1}{e}\diff[]{E_\mathrm{C}(x)}{x} +e\frac{k_\mathrm{B}T}{e}\mu_n \frac{n}{k_\mathrm{B}T} \left( \diff[]{E_\mathrm{F}}{x}-\diff[]{E_\mathrm{C}}{x}\right)=n\mu_n \diff[]{E_\mathrm{F}}{x}$$
Con lo que tenemos que la corriente de portadores total es $$\mathbf{J_n}=n\mu _n \nabla E_\mathrm{F}$$ Como en el equilibrio térmico $J_n=0$ entonces $E_\mathrm{F}=\text{cte}$
\end{tcolorbox}
\noindent Si unimos dos materiales con diferente nivel de Fermi, se produce una \textbf{redistribución de portadores} entre ambos hasta que se igualan los niveles de Fermi y se alcanza el equilibrio térmico. Alcanzar dicho equilibrio genera un desequilibrio de carga y por ende, dependencia de los extremos de la banda con la posición. 
\newpage
\subsubsection{Cuasi-Niveles de Fermi}
\begin{itemize}
    \item \textbf{En el equilibrio} tenemos un \textbf{único nivel de Fermi} y las densidades de portadores vienen dados por las ecuaciones \ref{eq: densidad de portadores MB} y \ref{eq: densidad de huecos MB} y la densidad de portadores intrínseca es \ref{eq: Densidad intrinseca de portadores} que tienen como expresión alterna en función del nivel de Fermi intrínseco
\begin{equation}
    n_0=n_\mathrm{i}\exp{\frac{E_\mathrm{F}-E_\mathrm{Fi}}{k_\mathrm{B}T}} \hspace{2cm} p_0=n_\mathrm{i}\exp{\frac{E_\mathrm{Fi}-E_\mathrm{F}}{k_\mathrm{B}T}}
    \label{eq: densidad intrinseca, nivel de fermi}
\end{equation}
    \item \textbf{Fuera del equilibrio} no podemos definir $E_\mathrm{F}$ no es definible y las expresiones $n(E_\mathrm{F})$ y $p(E_\mathrm{F})$ de equilibrio no son válidas. En la figura \ref{fig: niveles de fermi tipo N y P} representamos dichos niveles cuando están en equilibrio.
    \begin{figure}[h]
	\centering
    \begin{subfigure}[h]{0.3\textwidth}
	\includegraphics{Standalone/T2_02_Niveles_N.pdf}
	\caption{Distribución de niveles de energía en SCE tipo N.}
        \label{fig: Niveles  N}
    \end{subfigure}
    \hspace{3cm}
    \begin{subfigure}[h]{0.3\textwidth}
        	\includegraphics{Standalone/T2_03_Niveles_P.pdf}
        \caption{Distribución de niveles de energía en SCE tipo P.}
        \label{fig: Niveles P}
    \end{subfigure}
        \caption{Diagrama de las bandas de energía en equilibrio térmico.}
        \label{fig: niveles de fermi tipo N y P}
\end{figure}
    
    Para poder estudiar cómo se comportan las concentraciones de portadores en el no equilibrio, definimos \textbf{dos cuasi-niveles de Fermi}, uno para $e^-$ y otro para $h^+$.
    \begin{equation}
        n_0+\delta n= n_\mathrm{i}\exp{\frac{E_{\mathrm{F}n}-E_{\mathrm{F}i}}{k_\mathrm{B}T}} \hspace{2cm} p_0+\delta p= n_\mathrm{i}\exp{\frac{E_{\mathrm{F}i}-E_{\mathrm{F}p}}{k_\mathrm{B}T}}
        \label{eq: cuasiniveles de fermi}
    \end{equation}
    La concentración de portadores total es entonces función de los cuasi-niveles de Fermi.
\end{itemize}
\section{Procesos de generación-recombinación} \label{sec: procesos de generacion-recombinacion}
Los procesos de \textbf{generación-recombinación}, provocan variación local de las densidades de portadores, como hemos visto en la sección \ref{sec: Densidad intrinseca de portadores} los procesos de generación se dan cuando aportamos energía; mientras que los de recombinación se dan cuando se cede energía.\\\\
Estos procesos los podemos identificar con la dependencia temporal de la densidad de portadores $n=n(t)$, $p=p(t)$
\begin{equation}
    \left. \pdiff[]{n}{t} \right)_{\mathrm{g-r}}=G_\mathrm{ext}-\frac{n-n_0}{\tau_n}
    \label{eq: generacion recombinacion}
\end{equation}
Donde $G_\mathrm{ext}$ es la excitación externa y $\tau_n$ es el tiempo de vida media del exceso de portadores. De la ecuación \ref{eq: generacion recombinacion}, obtenemos dos expresiones para $\delta n$.
\begin{equation}
    \underbrace{\delta n =G_\mathrm{ext}\tau_n}_{\text{condiciones estacionarias}} \hspace{1cm} \underbrace{\delta n=G_\mathrm{ext} \tau_n \exp - \frac{t}{\tau_n}}_{\text{al cesar la excitación}}
\end{equation}
\newpage
\subsection{Procesos banda a banda} \label{sec: Procesos banda a banda}
Al incidir con luz, podemos observar diferentes procesos mostrados en la Figura \ref{fig: Procesos de gen recomb}, la generación se puede dar por \textbf{procesos térmicos}, responsables de devolver las concentraciones a los valores de equilibrio, donde $G_\mathrm{th}, \ R_\mathrm{th} \Rightarrow n_0\cdot p_0>n_\mathrm{i}^2$. Por \textbf{perturbaciones externas, inyección (óptica, eléctrica)} $G_\mathrm{L} \Rightarrow n\cdot p = n_\mathrm{i}^2$, donde $$n=n_0+\delta n \hspace{2cm} p=p_0+\delta p \hspace{2cm} 
\text{   con } \delta n \cong \delta p $$
La última ecuación se da ya que por cada $e^-$ generado se genera un $h^+$. Otros procesos de recombinación se los incluimos en \textit{U}.
Destacamos que la generación no es superficial, sino que es en todo el volumen por lo que las unidades de $[G,R]=\unit{[cm^{-3} \, s^{-1}]}=[\delta N/\tau_N]$

\begin{figure}[h]
    \centering
  \includegraphics{Standalone/T2_04_Procesos_de_gen_recomb.pdf}
    \caption{Distintos procesos de generación recombinación.}
    \label{fig: Procesos de gen recomb}
\end{figure}
\noindent Las perturbaciones externas, están caracterizadas por el la densidad de estados \textbf{libres} $N_i-n$. Si suponemos que teenemos perturbaciones débiles, es decir $n\ll N_\mathrm{c}$ y $p \ll N_\mathrm{v}$ tenemos 
\begin{equation}
    \left. \pdiff[]{n}{t} \right)_{\mathrm{g-r}}=\left. \pdiff[]{p}{t} \right)_{\mathrm{g-r}}= G(N_\mathrm{v}-p)(N_\mathrm{c}-n)-cnp\approx e-cnp
    \label{eq: coef emision y captura}
\end{equation}
Donde \textit{c} y $e=GN_\mathrm{c}N_\mathrm{v}$ son los coeficientes de captura y emisión respectivamente y son independientes de \textit{n} y \textit{p}. Distinguimos dos casos
\begin{itemize}
    \item \textbf{En equilibrio}: La ecuación \ref{eq: coef emision y captura}, se hace nula y obtenemos
    \begin{equation*}
        e=cn_0p_0=cn_\mathrm{i}^2 \Rightarrow \left. \pdiff[]{n}{t} \right)_{\mathrm{g-r}}=\left. \pdiff[]{p}{t} \right)_{\mathrm{g-r}}=-c(np-n_\mathrm{i}^2)=-\frac{n-n_0}{\tau_n}=-\frac{p-p_0}{\tau_p}
    \end{equation*}
    Operando obtenemos el tiempo de vida medio del exceso de portadores en funcion de las concentraciones de portadores
    \begin{equation}
        \tau_n=\tau_p=\tau= \frac{1}{c} \frac{n-n_0}{np-n_\mathrm{i}}=\frac{1}{c} \frac{p-p_0}{np-n_\mathrm{i}^2}
    \end{equation}
    \item \textbf{Material extrínseco, débil nivel de inyección}: Esta situacion se da cuando $\delta n=\delta p \ll n_0,p_0$ y utilizamos la concentración de los mayoritarios. De la expresión anterior
    $$ n p-n_i^2=\left(n_o+\delta n\right)\left(p_o+\delta p\right)-n_i^2=\left(p_o+n_o\right) \delta p+\delta p^2 \cong\left(p_o+n_o\right) \delta p $$ Obteniendo que el tiempo de vida medio es
    \begin{equation}
        \tau =\frac{1}{c} \frac{1}{p_0+n_0} \xrightarrow{\text{Mayoritarios}}     \left\{\begin{aligned}
          &\tau_{np^+}=\frac{1}{cp_0} \ \ \text{SCE tipo P}\\
          &\tau_{pn^+}=\frac{1}{cn_0} \ \ \text{SCE tipo N}\end{aligned} \right.
    \end{equation}
\end{itemize}
\newpage
\subsection{Procesos a través de centros}
Cuando introducimos impurezas o algunas alteraciones en la red cristalina, podemos encontrarnos con que aparecen nuevos nieles de energía permitidos dentro del GAP, estos niveles se denominan \textbf{centros recombinantes} y nos facilitan los procesos de recombinación
\begin{figure}[h]
    \centering
  \includegraphics{Standalone/T2_05_procesos_a_traves_de_centros.pdf}
    \caption{Cuatro procesos básicos de atrapamiento y emsión para una trampa aceptadora.}
    \label{fig: procesos a traves de centros}
\end{figure}

\noindent En la figura \ref{fig: procesos a traves de centros} vemos que cuando se emite un hueco al nivel intermedio, podemos emitir un electrón al nivel de conducción y al revés, cuando el nivel intermedio captura un electrón genera que el nivel de valencia capture un hueco.\\\\
En condiciones estacionarias ($\delta n= \delta p$) tenemos 
\begin{equation}
    \left. \pdiff[]{n}{t} \right)_{\mathrm{g-r}|_{E_\mathrm{T}}}=\left. \pdiff[]{p}{t} \right)_{\mathrm{g-r} |_{E_\mathrm{T}}}= g(t)- 
\cfrac{(n_0+p_0)\delta n+\delta n^2}{\cfrac{n^1+n_0+ \delta n}{c_p}+\cfrac{p^1+p_0+ \delta p}{c_n}}N_\mathrm{T}
\ \ \text{con} \ \  \left \{\begin{aligned}
    &n^1=N_\mathrm{c}\exp{\frac{E_\mathrm{T}-E_\mathrm{c}}{k_\mathrm{B}T}}\\
    &p^1=N_\mathrm{v}\exp{\frac{E_\mathrm{v}-E_\mathrm{T}}{k_\mathrm{B}T}}
\end{aligned} \right.
\label{eq: Ecuacion de Shockley‐Read‐Hall}
\end{equation}
Que podemos escribirla como $$\left. \pdiff[]{n}{t} \right)_{\mathrm{g-r}|_{E_\mathrm{T}}}=\left. \pdiff[]{p}{t} \right)_{\mathrm{g-r} |_{E_\mathrm{T}}}=g(t)\frac{\delta n}{\tau}=g(t)-\frac{\delta p}{\tau}$$ Donde el tiempo de vida media es 
\begin{equation}
\tau=\textcolor{red}{\frac{1}{N_T c_p}} \frac{n^1+n_o+\delta n}{n_o+p_o+\delta n}+\textcolor{blue}{\frac{1}{N_T c_n}} \frac{p^1+p_o+\delta n}{n_o+p_o+\delta n}=\textcolor{red}{\tau_{p n^{+}}} \frac{n^1+n_o+\delta n}{n_o+p_o+\delta n}+\textcolor{blue}{\tau_{n p^{+}}} \frac{p^1+p_o+\delta n}{n_o+p_o+\delta n}
\label{eq: tiempo de vida medio proceso a traves de centro}
\end{equation}
Que volviendo a aplicar las condiciones de \textbf{débil nivel de inyección} ($\delta n=\delta p \ll n_0,p_0$) para un material extrínseco las concentraciones en semiconductores extrínsecos tienen vidas medias en este caso de los \textbf{portadores minoritarios}
\begin{equation}
    \underbrace{\tau_{np^+}=\frac{1}{N_\mathrm{T}c_p}}_{\text{SCE tipo N}} \hspace{2cm}
    \underbrace{\tau_{pn^+}=\frac{1}{N_\mathrm{T}c_n}}_{\text{SCE tipo P}}
\end{equation}
Por lo que para conseguir $\tau$ pequeñas necesitamos $N_\mathrm{T}\uparrow$ y $E_\mathrm{T} \approx \mathrm{MIDGAP}\Rightarrow n^1,p^1 \downarrow$
\newpage 
\subsection{Recombinación superficial}
Debido a su gran tamaño, los semiconductores se tratan como si no tuvieran límites físicos, aunque los estados ligados se crean por las condiciones de contorno que hemos visto en la sección \ref{sec: Modelo de enlace}\footnote{Para el modelo de Kroning-Pennney utilizamos la normalización de caja}. En esta sección, estudiamos la celda como un sistema y aplicamos el Teorema de Bloch, lo que hace que no tengamos en cuenta la finitud del cristal debido a que su tamaño es mucho mayor que cualquiera de las distancias intermoleculares de la red.\\

\noindent Como hemos en la sección anterior, introducir imperfecciones en la red cristalina nos permite crear niveles de energía en el interior del GAP, esto se debe a que las imperfecciones introducen perturbaciones en la periodicidad de la red y es cuando se interrumpe de manera abrupta esa periodicidad cuando se generan dichos centros recombinantes. Cuando llegamos a la superficie del cristal, pasa exáctamente lo mismo, el potencial periódico se interrumpe de manera abrupta lo que hace que aparezcan estados de energía dentro del GAP.
\begin{wrapfigure}{r}{0.4\textwidth}
    \centering\includegraphics{Standalone/T2_06_estados_superficie.pdf}
 \caption{Distribución de los estados de la superficie en la banda del GAP.}
 \label{fig: estados superficie}
\end{wrapfigure}

\noindent Las ecuaciones \ref{eq: Ecuacion de Shockley‐Read‐Hall} y \ref{eq: tiempo de vida medio proceso a traves de centro} nos muestran el ratio de recombinación y el tiempo de vida de los portadores minoritarios así que si consideramos un SCE tipo N, en este caso la densidad de supercficie de los minoritarios será $$p_\mathrm{s}=p_\mathrm{s0}+\delta$$ Por lo que si volvemos a utilizar el débil nivel de inyección, obtenemos que 
\begin{equation}
   R_\mathrm{s}=\frac{(n_0+p_0)\delta n}{\frac{n^1+n_0}{c_p}+\frac{p^1+p_0}{c_n}}N_\mathrm{s}
\end{equation}

\noindent Con lo que podemos definir la velocidad de recombinación superficial como $$s=\frac{R_\mathrm{s}}{\delta}=\frac{1}{\tau}$$ que para un semiconductor extrínseco corresponde
\begin{equation}
\underbrace{s_{np^+}=c_{n\mathrm{s}}N_\mathrm{s}}_{\text{SCE tipo P}} \hspace{2cm} \underbrace{s_{pn^+}=c_{p\mathrm{s}}N_\mathrm{s}}_{\text{SCE tipo N}}
 \label {eq: Velocidad de recombinación}   
\end{equation}
Este proceso de recombinación nos genera una corriente, conocida como \textbf{corriente de recombinación superficial}
\begin{equation}
    J_\mathrm{rec}=qR_\mathrm{s}=qs\delta
    \label{eq: corriente de recombinacion superficial}
\end{equation}
Esta expresión se utiliza como condición de contorno para resolver la ecuación de continuidad ya que las corrientes de difusión y recombinación han de ser iguales.\\\\ Por último, vamos a definir como \textbf{contacto metal/óhmico} a aquella disrrupción en la que $s\rightarrow{} \infty$ y \textbf{superficie ideal} a aquella en la que $s\rightarrow{} 0$
\section{Ecuación de Continuidad}
Las modificaciones de las concentraciones de portadores, nos generan corrientes, ya sean de arrastre o de difusión, así como los procesos de generación-recombinación, que influyen en estas últimas. Por ello queremos estudiar cómo se comportan dichas concentraciones cuando combinamos ambos efectos, corrientes (i.e campos eléctricos y gradientes de concentración) con generación-recombinación. En el Apéndice \ref{sec: Soluciones de la Ecuacion de Continuidad} se detallan todas las soluciones.
\begin{figure}[h]
    \centering
    \includegraphics{Standalone/T2_07_flujo_de_corriente_en_un.pdf}
    \caption{Variación de corriente al atravesar un elemento de volumen de un semiconductor.}
    \label{fig: flujo de corriente en un diferencial de volumen}
\end{figure}

Para construir la ecuación de continuidad vamos a combinar las ecuaciones \ref{eq: Densidad de corriente total} y \ref{eq: generacion recombinacion}, para combinarlas separamos entre $e^-$ y $h^+$. Primero calculamos la variación del flujo de portadores, recordamos que el flujo en este caso es $\Phi=\mathbf{J}\cdot \mathbf{\hat{n}}\,A$
\begin{equation}
    A\mathrm{d}x \pdiff[]{n}{t}=\left [\frac{J_n(x)}{q}-\frac{J(x+\mathrm{d}x)}{q}+ \left. \pdiff[]{n}{t} \right )_{g-r} \right ] A \mathrm{d}x
\end{equation}
Que extendiendo a 3D y eliminando las partes correspondientes al volumen
\begin{equation}
    \pdiff[]{n}{t}=-\frac{1}{q} \nabla_{\mathbf{r}}\cdot \mathbf{J} + \left. \pdiff[]{n}{t} \right )_{g-r}
\end{equation}
Mediante las ecuaciones \ref{eq: Densidad de corriente total} y \ref{eq: generacion recombinacion} obtenemos finalmente la ecuación de continuidad para $e^-$ y $h^+$
\begin{equation}
\pdiff[]{n}{t}=D_n \nabla_\mathbf{r}^2 n+ \mu_n \nabla_\mathbf{r} (n\mathbf{E})-\frac{n-n_0}{\tau_n}+g_\mathrm{L}
\label{eq: continuidad electrones}
\end{equation}

\begin{equation}
    \pdiff[]{p}{t}=D_p \nabla_\mathbf{r}^2 p \textcolor{red}{-} \mu_p \nabla_\mathbf{r} (p\mathbf{E})-\frac{p-p_0}{\tau_p}+g_\mathrm{L}
    \label{eq: continuidad huecos}
\end{equation}
Podemos acoplar las ecuaciones a través de la ecuación de Poisson 
\begin{equation}
    \nabla_\mathbf{r}\cdot \mathbf{E}=-\frac{e}{\epsilon} (n-p-N_\mathrm{D}^++N_\mathrm{A}^-)
    \label{eq: ecuacion de poisson huecos y electrones}
\end{equation}
Dichas ecuaciones, pueden resolverse mediante aproximaciones razonables.\newpage
Si además consideramos que $n=n_0\delta n$ y $p=p_0+\delta p$ las ecuaciones de continuidad son
\begin{equation}
    \pdiff[]{\delta n}{t}=D_n \nabla_\mathbf{r}^2 \delta n+ \mu_n \nabla_\mathbf{r} (\delta n\mathbf{E})-\frac{\delta n}{\tau_n}+g_\mathrm{L}
    \label{eq: continuidad exceso de electrones}
\end{equation}
\begin{equation}
    \pdiff[]{\delta p}{t}=D_p \nabla_\mathbf{r}^2 \delta p\textcolor{red}{-} \mu_p \nabla_\mathbf{r} (\delta p\mathbf{E})-\frac{\delta p}{\tau_p}+g_\mathrm{L}
    \label{eq: continuidad exceso de huecos}
\end{equation}
Una vez construidas todas las ecuaciones, vamos a explicar las implicaciones que tienen: si se crea un exceso de $e^-$ y $h^+$ en cierta posición de un semiconductor sobre el que se ha aplicado un campo $E_\mathrm{app}$, los excesos de $e^-$ y $h^+$ tienden a moverse en direcciones opuestas. Esto origina un desequilibrio de carga y un campo interno $E_\mathrm{int}$ que tiende a mantenerlos unido, de modo que $e^-$ y $h^+$ se desplazarán unidos por arrastre o difusión con una única movilidad o coeficiente de difusión efectivos. Es lo que conocemos como \textbf{transporte ambipolar}.
\begin{figure}[h]
    \centering
    \includegraphics{Standalone/T2_08_La_creacion_de_un_campo.pdf}
    \caption{La creación de un campo eléctrico interno hace que los $e^-$ y los $h^+$ tiendan a juntarse, mientras que el aplicado hace que tiendan a separarse}
    \label{fig:enter-label}
\end{figure}

\subsection{Derivación de la ecuación de transporte ambipolar}
Una vez hecho el análisis cualitativo, vamos a caracterizar las ecuaciones. 

\noindent Como las ecuaciones \ref{eq: ecuacion de poisson huecos y electrones}- \ref{eq: continuidad exceso de huecos}, son muy difíciles de tratar, vamos a hacer algunas aproximaciones, una de ellas es la \textbf{aproximación de cuasi neutralidad}.\\\\
\noindent El campo eléctrico total $\mathbf{E}=\mathbf{E_\mathrm{int}}+ \mathbf{E_\mathrm{app}}$, si lo aplicamos la ecuación \ref{eq: ecuacion de poisson huecos y electrones} $$\nabla \cdot \mathbf{E}=-\frac{e}{\epsilon} (n_0+\delta n-p_0 -\delta p -N_\mathrm{D}^++N_\mathrm{A}^-)$$
Podemos asociar al campo interno
\begin{equation}
    \nabla \cdot \mathbf{E_\mathrm{int}}=-\frac{e}{\epsilon}(\delta n-\delta p)
\end{equation}
Este término solo es interesante sí consideramos la ACN, dada por la siguiente ecuación
\begin{equation}
    \delta n =n-n_0 \cong p-p_0 =\delta p
\end{equation}
Por lo que el término que corresponde al campo eléctrico interno es despreciable. \newpage

\noindent Queremos eliminar, el término $\nabla \cdot \mathbf{E}$, para ello supongamos que el factor de generación sea el mismo tanto para electrones como para huecos $g_n=g_p\equiv g_\mathrm{L}$. Vamos a operar con las ecuaciones \ref{eq: continuidad electrones} y \ref{eq: continuidad huecos} multiplicando por $\mu_p p$ la primera y $\mu_n n$ la segunda y sumándolas, además asumiremos la aproximación de cuasi-neutralidad y $n=n_0+\delta n$, $p=p_0+\delta p$.
\begin{equation*}
\begin{aligned}
&\mu_p p \cdot \left [\pdiff[]{n}{t}=D_n \nabla_\mathbf{r}^2 n+ \mu_n \nabla_\mathbf{r} (n\mathbf{E})-\frac{n-n_0}{\tau_n}+g_\mathrm{L} \right ] \\
    &+ \ \mu_n n \cdot \left [ \pdiff[]{p}{t}=D_p \nabla_\mathbf{r}^2 p  \textcolor{red}{- } \mu_p \nabla_\mathbf{r} (p\mathbf{E})-\frac{p-p_0}{\tau_p}+g_\mathrm{L} \right ] \end{aligned}
\end{equation*}
Obteniendo la ecuación ambipolar de difusión
\begin{equation}
    \boxed{\pdiff[]{\delta n}{t}+ \mu ^* \mathbf{E} \nabla \delta n-D^* \nabla^2 \delta n +\frac{\delta n}{\tau^*}-g_\mathrm{L}=0}
    \label{eq: ecuacion ambipolar de difusion}
\end{equation}
Donde los términos $$\mu^*=\frac{\mu_n \mu_p n_0-\mu_n\mu_p p_0}{\mu_pp+\mu_nn} \hspace{1.5cm} D^*=\frac{D_n\mu_pp+\mu_nD_pn}{\mu_pp+\mu_nn} \hspace{1.5cm} \tau^*=\frac{\tau_n\tau_p(\mu_pp+\mu_nn)}{\mu_pp\tau_p+\mu_nn\tau_n}$$ 
Que en principio es una ecuación no-lineal, donde los coeficientes dependen de $\delta n$.
La ecuación \ref{eq: ecuacion ambipolar de difusion} presenta un término del campo eléctrico simplificado, ya que \textbf{E} está fuera del gradiente. En función del material, los coeficientes serán:
\begin{itemize}
    \item \textbf{Semiconductores intrínsecos}: $\mu^*$ nula, $D^*$ y $\tau^*$ son independientes del tiempo.
    \item \textbf{Semiconductores extrínsecos y débil nivel de inyección}: Los parámetros de la ecuación ambipolar coinciden con los de los portadores minoritarios. \textit{Un SCE tipo N.} $$n_0\gg p_0 \hspace{1cm} \delta n\ll n_0 \hspace{0.5cm} \Rightarrow \hspace{0.5cm} D^* \cong D_p \hspace{1.5cm} \tau^*\cong \tau_p \hspace{1.5cm} \mu^*\cong \mu_p$$
    Para calcular \textit{n} y \textit{p} en un semiconductor extrínseco:
    \begin{itemize}
        \item La ecuación de continuidad de minoritarios con término de \textbf{E} despreciable.
        \item $n=n_0+\delta n$ y $p=p_0+\delta m$
        \item \textbf{E}: La ecuación de continuidad de los mayoritarios sin despreciar \textbf{E}.
    \end{itemize}
\end{itemize}

\begin{wrapfigure}{r}{0.4\textwidth}
    \centering
    \includegraphics{Standalone/T2_09_Inyeccion_de_huecos_en_una.pdf}
    \caption{Inyección de huecos en una pequeña región de la superficie de un SCE tipo N.}
\end{wrapfigure}

\subsection{Tiempo de relacación dieléctrico}

Si despreciamos los efectos de los procesos de generación-recombinación y la difusión, vamos a estudiar cuánto tarda un exceso de carga en desaparecer.\\\\
Supongamos que tenemos una densidad de carga $\rho$ uniforme sobre una cierta distancia próxima a la superficie, inicialmente toma le valor $\rho = e\delta p$, además tenemos que considerar que el material es óhmico $(\mathbf{J}=\sigma \mathbf{E})$. Si utilizamos la ley de Gauss y la Ecuación de continuidad del campo eléctrico. $$\nabla \cdot \mathbf{J}=\sigma \nabla \cdot \mathbf{E}=\frac{\sigma \rho}{\epsilon} \Rightarrow \frac{\sigma \rho}{\epsilon}=-\diff[]{\rho}{t} \Rightarrow \rho (t)=\rho_0 \exp{-\frac{t}{\tau_\mathrm{d}}} \text{ con } \boxed{\tau_\mathrm{d}=\frac{\epsilon}{\sigma}}$$
\section{\textcolor{red}{Efecto Hall}}
El Efecto Hall aparece cuando sometemos a un semiconductor por el que pasa una corriente a un campo magnético y no es más que consecuencia de las fuerzas que ejercen los campos eléctrico y magnético.
\begin{figure}[h]
    \centering
    \includegraphics{Standalone/T2_10_efecto_Hall.pdf}
    \caption{Geometría del Efecto Hall en un semiconductor tipo N.}
    \label{fig:efecto Hall}
\end{figure}

\noindent Cuando aparece un campo magnético, la dinámica de la partícula viene dada por la ecuación de la fuerza de Lorentz.
\begin{equation}
    \mathbf{F}=q (\mathbf{E}+\mathbf{v} \times \mathbf{B})
    \label{eq: Fuerza de Lorentz}
\end{equation}
Por la geometría que se ve en la Figura \ref{fig:efecto Hall}\footnote{Los ejes son un poco ñeh pero no queria hacer otro dibujo.}, el equilibrio de fuerzas se alcanza cuando sólo tenemos componente Y en el campo eléctrico. Por lo que si trabajamos en módulo $$E_y=v_xB_z$$ que es lo que conocemos como \textbf{campo Hall}, este campo produce una diferencia de potencial llamada \textbf{potencial Hall}
\begin{equation}
    V_\mathrm{H}=+E_\mathrm{H}W=\frac{J_xB_z}{qn}W=|R_\mathrm{H}|J_xB_zW \Rightarrow \mathbf{E}=R_\mathrm{H}\mathbf{B}\times \mathbf{J} \Rightarrow R_\mathrm{H}=\frac{1}{qn}
    \label{eq: Potencial Hall}
\end{equation}
El factor $R_\mathrm{H}$ lo definimos como \textbf{coeficente Hall}, que será negativo o positivo en función de si el semiconductor es tipo N o tipo P.\\\\
Una aplicación del Efecto Hall, es la caracterización de semiconductores, si medimos $V_\mathrm{H}, J_x, B_z$ y por la geometría del semiconductor, podemos determinar la concentración de portadores extrayéndolo de la ecuación \ref{eq: Potencial Hall}. Además si medimos la conductividad, podemos también determinar la movilidad $\mu=\sigma |R_\mathrm{H}|$.\\\\
Hemos desarrollado para un SC tipo N pero análogamente se puede aplicar con la concentración de $h^+$.

\part{Dispositivos electrónicos fundamentales}
\pagestyle{fancy}
\fancyhead[R]{Electrónica Física: Tema 3}
\fancyhead[L]{4º Grado en Física}
\chapter{Diodos}
Una vez estudiado del comportamiento de los portadores de carga en equilibrio térmico, determinado el nivel de Fermi. Después hemos estudiado situaciones fuera del equilibrio térmico, en las que tenemos excesos en las concentraciones de portadores de carga. Queremos estudiar qué ocurre cuando hacemos contacto entre dos semiconductores, uno tipo P y uno tipo N.
\section{\textcolor{red}{Unión PN. Comportamiento estático y dinámico}}

\begin{wrapfigure}{r}{0.35\textwidth}
	\centering
    \begin{subfigure}[h]{0.3\textwidth}
	\includegraphics{Standalone/T3_01_Diodo.pdf}
	\caption{Representación de un diodo PN, con polaridad               positiva.}
        \label{fig: Diodo}
    \end{subfigure}
    \begin{subfigure}[h]{0.3\textwidth}
        \includegraphics{Standalone/T3_02_Transferencia_de_y_en_una.pdf}
        \caption{Transferencia de $e^-$ y $h^+$ en una unión PN idealmente uniforme.}
    \end{subfigure}
        \caption{Representación de una uníon PN.}
        \label{fig: Representacion union PN}
\end{wrapfigure}

Para poder formar una unión PN, no podemos hacerlo de cualquier manera, la estructura cristalina del semiconductor debe mantenerse. Por ello la unión PN nos dará un \textbf{único cristal}, dividido en dos regiones, una unión con un exceso de huecos, que denominamos \textbf{región P}; y otra con un exceso de electrones que denominamos \textbf{región N}.\\\\
La figura \ref{fig: Representacion union PN}, las concentraciones de portadores de cargas son uniformes en cada región pero hay cambio abrupto en la unión, ya que pasamos de tener huecos a electrones, esto se conoce como \textbf{unión abrupta}.\\\\
Si analizamos la figura \ref{fig: Diodo}, considerando $I=0$, inicialmente tenemos una gran concentración de electrones y huecos en sus respectivas regiones. Como el sistema tiende a estar en equilibrio aparecen fuerzas de difusión que hace que haya un flujo de electrones y huecos entre ambas regiones. Este movimiento de cargas nos genera una diferencia de potencial cerca de la unión metalúrgica y por ende un campo eléctrico, en la dirección de la región positiva (N) a la negativa (P).
\clearpage
\textbf{¿Por qué la región en N es positiva y en P negativa?}, como la difusión es un transporte de portadores de carga, al haber difusión, los electrones se irán depositando en una región cercana a la unión metalúrgica, por lo que la zona positiva es se encuentra más cerecana a la región N mientras que la negativa estará más cercana a a la región P. Como se muestra en la figura \ref{fig:Zona de carga espacial}

\begin{figure}[h]
    \centering
    \includegraphics{Standalone/T3_03_Zona_de_carga_espacial.pdf}
    \caption{La zona de carga espacial (ZCE), el campo eléctrico y las fuerzas que actúan sobre los portadores} 
    \label{fig:Zona de carga espacial}
\end{figure}
\begin{wrapfigure}{r}{0.4\textwidth}
    \centering
    \begin{subfigure}[h]{0.35\textwidth}
    \includegraphics{Standalone/T3_04_flujos_de_corriente_PN.pdf}
    \caption{Representación de flujos de corriente en una unión PN.}
    \label{fig: flujos de corriente PN}
    \end{subfigure}
    \begin{subfigure}[h]{0.35\textwidth}
        \includegraphics{Standalone/T3_05_representacion_IV_PN.pdf}
 \caption{Representación ideal I-V de un diodo de unión PN.}
    \label{fig: representacion IV PN}
    \end{subfigure}
\end{wrapfigure}

\noindent Cuando este sistema alcance el equilibrio, las furezas eléctricas y las de difusión se compensan, por lo que la corriente en el interior será nulo. 

\subsubsection{Descripción cualitativa del flujo de corriente}

En la figura \ref{fig: flujos de corriente PN}, vemos qué ocurre cuando existe una diferencia de potencial externa entre las regiones del semiconductor. 
\begin{itemize}
    \item $V_\mathrm{P} - V_\mathrm{N}=V_\mathrm{PN} >0$: Disminuye el campo de equilibrio, la contribución de la difusión es superior al arrastre por lo que I aumenta con V. Es lo que denominaremos \textbf{polarización directa}.
    \item $V_\mathrm{P} - V_\mathrm{N}=V_\mathrm{PN} <0$: El campo externo  tiene la misma dirección  que el campo de equilibrio, por lo que la difusión es menor que el arrastre. En este caso la intensidad es muy baja e independiente de V. La denominaremos \textbf{polarización inversa}.
\end{itemize}
\clearpage
\subsection{Unión PN en equilibrio} \label{sec: Union PN en equilibrio}
Cuando las regiones PN están aisladas y en equilibrio, demostramos en la sección 2.1.3 que el nivel de Fermi no varía con la posición. Esto se mantendrá cuando las regiones estén unidas y la unión esté en equilibrio.
\begin{figure}[h]
    \centering
    \begin{subfigure}[H]{0.5\textwidth}
\includegraphics{Standalone/T3_06_Niveles_de_energia_y_Fermi.pdf}
    \caption{Niveles de energía en la unión PN.}
    \label{fig:Niveles de energia y Fermi}
\end{subfigure}
\hfill
\begin{subfigure}[h]{0.4\textwidth}
\includegraphics{Standalone/T3_07_concentraciones_PN.pdf}
    \caption{Concentraciones de portadores en diferentes zonas.}
    \label{fig: concentraciones PN}
\end{subfigure}
\caption{Representación de los niveles de energía y de las densidaades de portadores.}
\end{figure}

\noindent Como el nivel de Fermi tiene que mantenerse constante en el equilibrio, al cambiar las densidades de portadores, necesariamente cambiarán los niveles de la banda de conducción y de valencia. Los electrones en la banda de conducción están sometidos a una barrera de potencial $V_\mathrm{B}$ que denominamos potencial de contacto. Es esta barrera de potencial la que mantiene el equilibrio en las densidades de electrones en la región P y la N y análogamente para los huecos, por lo que no tenemos corrientes en ZCE.
\subsubsection{\textcolor{red}{Potencial de contacto}}

Como acabamos de indicar el potencial de contacto \textbf{sólo aparece en el lado N}. Por estar en equilibrio, $I=0$, si usamos la ecuación \ref{eq: continuidad electrones}, tendrá la forma
\begin{equation*}
    -D_n\diff[]{n}{x}=n\,\mu_nE \ \text{con } \ D_n=\mu_n \frac{K_\mathrm{B}T}{e}
\end{equation*}
Despejando obtenemos
\begin{equation}
    E=-\frac{K_\mathrm{B} T}{e} \frac{1}{n} \diff[]{n}{x} \Rightarrow \boxed{V_\mathrm{B}=\frac{K_\mathrm{B} T}{e} \int_{n_{p0}}^{n_{n0}} \frac{\mathrm{d} n}{n}=\frac{k_\mathrm{B} T}{e} \ln{\frac{n_{n0}}{n_{p0}}}=\frac{k_\mathrm{B} T}{e} \ln{\frac{N_\mathrm{D}N_\mathrm{A}}{n_\mathrm{i}^2}}}
    \label{eq: Potencial de contacto}
\end{equation}
Donde $$n_{n0}=N_\mathrm{c}\exp \frac{(E_\mathrm{F}-E_\mathrm{C})_\mathrm{N}}{k_\mathrm{B}T} \hspace{2cm} p_{p0}=N_\mathrm{c}\exp \frac{(E_\mathrm{F}-E_\mathrm{C})_\mathrm{P}}{k_\mathrm{B}T}$$
si sustituimos en la ecuación \ref{eq: Potencial de contacto} finalmente obtenemos que el potencial de contacto siempre es positivo y lo podemos expresar en función de las diferencias de los niveles de Fermi de las regiones P y N.
\begin{equation}
    V_\mathrm{B}=\frac{E_\mathrm{FN}-E_\mathrm{FP}}{e}
\end{equation}
\newpage
\subsubsection{\textcolor{red}{Campo eléctrico y potencial. Aproximación de vaciamiento total}} \label{sec: Vaciamiento total}
Hasta ahora hemos estudiado el potencial de contacto en ausencia de campo eléctrico. Si aplicamos un campo eléctrico, éste podrá aumentar o disminuir. Para estudiar cómo afecta el campo externo haremos la siguiente aproximación: la densidad de portadores libres en la zona de transición es despreciable frente a la concentración de impurezas ionizadas, conocida como \textbf{aproximación de vaciamiento total}.\\

\begin{wrapfigure}{r}{0.27\textwidth}
    \includegraphics{Standalone/T3_08_densidad_de_carga_asimetrica.pdf}
    \caption{La densidad de carga tiene que ser la misma en ambas regiones aunque la zona de transición es asimétrica.}
    \label{fig: densidad de carga asimetrica}
\end{wrapfigure}
En la zona de carga espacial:
\begin{itemize}
    \item $0<x<x_n \rightarrow{} \rho=eN_\mathrm{D}$.   
    \item $x_p<x<0 \rightarrow{} \rho=-eN_\mathrm{A}$.
    \item Tenemos una condición extra, la condición de neutralidad $$N_\mathrm{D}x_n=N_\mathrm{A}x_p$$
\end{itemize}
Para estudiar cómo afecta el campo eléctrico al potencial y a las concentraciones de portadores. Para ello vamos a utilizar la ley de Gauss del campo eléctrico, vamos a desarrollar un modelo unidimensional.
\begin{equation}
    \diff[]{E}{x}=\frac{\rho}{\epsilon} \xrightarrow{\text{Campo conservativo}} \diff[2]{V_\mathrm{j}}{x}=-\frac{\rho}{\epsilon}
\end{equation}
Si dividimos en las regiones $x_p<x<0$ y $0<x<x_n$ obtenemos\\
\begin{equation}
    \begin{aligned}
        &x_p<x<0 \xrightarrow{} \diff[]{E}{x}=-\frac{eN_\mathrm{A}}{\epsilon} \xrightarrow[]{} E=-\frac{eN_\mathrm{A}(x-x_p)}{\epsilon}\\
        &0<x<x_n \xrightarrow{} \diff[]{E}{x}=\frac{eN_\mathrm{D}}{\epsilon}\xrightarrow[]{} E=\frac{eN_\mathrm{D}(x-x_n)}{\epsilon}\\
    \end{aligned}
    \label{eq: campo electrico en union PN}
\end{equation}

\begin{wrapfigure}{r}{0.27\textwidth}
    \includegraphics{Standalone/T3_09_Los_campos_son_lineales_como.pdf}
    \caption{Los campos son lineales como hemos visto en las ecuaciones.}
\end{wrapfigure}

Donde el campo tiene que tener las siguientes condiciones de contorno $E(x_n)=0$, $E(x_p)=0$ y tiene que ser continuo en $x=0$, dado por $N_\mathrm{D}x_n=-N_\mathrm{A}x_p$.\\

\noindent Como hemos visto en las ecuaciones \ref{eq: campo electrico en union PN} el campo eléctrico es lineal por lo que el potencial ha de ser cuadrático (no hay más que integrar en $x$ entre $x_p$ y $x_n$)
\begin{equation}
\begin{aligned}
    & V_\mathrm{j}(x_p)=\frac{eN_\mathrm{D}(x-x_n)^2}{2\epsilon}\\
    & V_\mathrm{j}(x_n)=-\frac{eN_\mathrm{D}(x-x_n)^2}{2\epsilon}+V_\mathrm{B}
    \end{aligned}
\end{equation}
Con lo que todas las relaciones para el campo, potencial y concentraciones son funcion de $x_n$ y $x_p$, que son parámetros no conocidos. Si incluimos la continuidad de $V_\mathrm{B}$ en $x=0$
\begin{equation}
    V_\mathrm{B}=\frac{e}{2\epsilon} \left(N_\mathrm{A}x_p^2+N_\mathrm{D}x_n^2\right)
\end{equation}
\clearpage
Podemos así obtener las expresiones para dichos puntos
\begin{equation}
x_p=-\left(\frac{2 \epsilon V_\mathrm{B}}{e} \frac{N_\mathrm{D}}{N_\mathrm{A}} \frac{1}{N_\mathrm{A}+N_\mathrm{D}}\right)^{1 / 2} \hfill x_n=\left(\frac{2 \epsilon V_\mathrm{B}}{e} \frac{N_\mathrm{A}}{N_\mathrm{D}} \frac{1}{N_\mathrm{A}+N_\mathrm{D}}\right)^{1 / 2}
\end{equation}
\noindent Gracias a esto podemos hallar la \textbf{anchura de la zona de transición}.
\begin{equation}
W=x_n-x_p=\left[\frac{2 \epsilon V_\mathrm{B}}{e\left(N_\mathrm{A}+N_\mathrm{D}\right)}\right]^{1 / 2}\left(\sqrt{\frac{N_\mathrm{A}}{N_\mathrm{D}}}+\sqrt{\frac{N_\mathrm{D}}{N_\mathrm{A}}}\right)
\label{eq: Anchura de la zona de transicion}
\end{equation}
La \textbf{carga por unidad de área en un lado de la ZCE}.
\begin{equation}
    Q=e N_D x_n=\left(\frac{2 \epsilon V_\mathrm{B} e N_\mathrm{A} N_\mathrm{D}}{N_\mathrm{A}+N_\mathrm{D}}\right)^{1 / 2}
    \label{eq: Carga por unidad de área en un lado de la ZCE}
\end{equation}
\subsubsection{Uniones asimétricas}
\noindent Volviendo a la figura \ref{fig: densidad de carga asimetrica}, si acercamos al origen el punto $x_p$. Considerando $N_\mathrm{A}\gg N_\mathrm{D} \Rightarrow x_p\approx 0$ por lo que la posición $x_n$ es
\begin{equation}
    x_n=\left (\frac{2\epsilon V_\mathrm{B}}{e} \frac{1}{N_\mathrm{D}} \right )^{1/2}
\end{equation}

\noindent Comentar que todo el potencial cae en el lado menos impurificado de la ZCE. Para $N_\mathrm{D}>>N_\mathrm{A}$ se hace de manera análoga.
\subsection{Unión PN bajo polarización}

\begin{wrapfigure}{r}{0.3\textwidth}
    \centering
    \includegraphics{Standalone/T3_10_Inverso_de_la_capacidad_al.pdf}
    \caption{Inverso de la capacidad al cuadrado frente al potencial.}
    \label{fig:enter-label}
\end{wrapfigure}

\noindent La región que presenta más resistividad, es la zona de transición, ya que toda la tensión cae en ella. En la zona de transición, en la zona de transición cae como $$V_\mathrm{j}(x_n)-V_\mathrm{j}(x_p)=V_\mathrm{B}-V$$
Introducir el potencial externo, nos suaviza o agudiza las curvas de la figura \ref{fig:Niveles de energia y Fermi} para modelizar este comportamiento hacemos desarrollo análogo a la de la sección anterior solo que cambiando $V_\mathrm{B}$ por $V_\mathrm{B}-V$.
\subsubsection{Capacidad de transición}
\noindent Como hemos visto a ecuación \ref{eq: Carga por unidad de área en un lado de la ZCE} nos da una expresión para la carga que tenemos en la zona de transición por lo que al haber además una diferencia de potencial entre los extremos de la zona, obtenemos que existe una capacidad por unidad de área en la zona de transición.
\begin{equation}
    C_\mathrm{j}=C_\mathrm{T}=\left | \diff[]{Q}{V} \right|=\left( 
\frac{\epsilon e N_\mathrm{A}N_\mathrm{D}}{2(N_\mathrm{A}+N_\mathrm{D})} \right)^{(1/2)} (V_\mathrm{B}-V)^{-1/2}=\frac{\epsilon}{W}
\label{eq: Capacidad de transicion}
\end{equation}
\newpage
\subsection{Portadores y corrientes en la unión PN polarizada}
Para hacer este desarrollo vamos a calcular las densidades en las zonas neutras, para ello utilizaremos las densidades de portadores minoritarios en los bordes de ZCE.\\\\
Vamos a realizar las siguientes aproximaciones:
\begin{itemize}
    \item Condiciones estacionarias, \textbf{ausencia de generación recombinación en ZCE}. Por lo que se conservarán las corrientes $J_p$ y $J_n$ en ZCE.   \hfill C$_1$
    \item \textbf{Débil polarización}, $J_\mathrm{p}\cong 0$ frente a las componentes de arraster y difusión.   \hfill C$_2$
    \item  Débil nivel de inyección $p_p-p_{p0}<<p_{p0}$ y $n_n-n_{n0}<<n_{n0}$.    \hfill C$_3$
\end{itemize}
Donde vamos a utilizar la siguiente notación: $n_{p,n}=n(x_{p,n})$ y $n_{n0,p0}$ para las densidades en el equilibrio.
Ahora sí podemos calcular las densidades de portadores minoritarios.\\\\ Haremos el desasrrollo para el lado N , por lo que utilizaremos los huecos, de la ecuación \ref{eq: continuidad huecos} 
\begin{equation*}
    \cancelto{0 \text{ por C}_3}{\pdiff[]{p}{t}}=D_p \nabla_\mathbf{r}^2 p- \mu_p \nabla_\mathbf{r} (p\mathbf{E})-\cancelto{0 \text{ por }\mathrm{C}_1}{\frac{p-p_0}{\tau_p}+g_\mathrm{L}}
\end{equation*}
entonces
\begin{equation*}
    \nabla_r(\mu_p p \mathbf{E}-D_p\nabla_r p)=0 \xrightarrow{\mathrm{C_1}}\cancelto{0\text{ por C}_2}{\frac{\mathbf{J}_p}{e}}=\mu_p p\mathbf{E}-D_p\nabla_r p
\end{equation*}
Separando en componentes (solo hacemos $x$), utilizamos la ley de Gauss del campo eléctrico y obtenemos
\begin{equation}
    -\mu_p p \diff[]{V_\mathrm{j}}{x}=D_p\diff[]{p}{x} \xrightarrow{\text{Integramos en x}} \frac{\mathrm{d}p}{p}=\frac{-\mu_p}{D_p} \mathrm{d}V_\mathrm{j} \Rightarrow p(x)=p_p \exp{\left ( -\frac{eV_\mathrm{j}(x)}{k_\mathrm{B}T} \right )}
    \label{eq: Densidad portadores minoritarios p}
\end{equation}
Como $V_\mathrm{j}(x_n)=V_\mathrm{B}-V$ la expresión para los electrones
\begin{equation}
    n(x)=n_n\exp{\left ( -\frac{e(V_\mathrm{B}-V-V_\mathrm{j}(x))}{k_\mathrm{B}T} \right )}
    \label{eq: Densidad portadores minoritarios n}
\end{equation}
Vamos a establecer ahora los valores en los límites para los valores de $p_p$ y $n_n$. 
\begin{equation}
          \begin{aligned}
         & \text{Si } x=x_n, \ V_\mathrm{j}=V_\mathrm{B}-V \xrightarrow{} p(x_n)=p_n=p_p\exp{\left(-\frac{e(V_\mathrm{B}-V)}{k_\mathrm{B}T} \right)}\\\\ & \text{Si }  x=x_p, \ V_\mathrm{j}=0 \xrightarrow{} n(x_p)=n_p=n_n\exp{\left(-\frac{e(V_\mathrm{B}-V)}{k_\mathrm{B}T} \right)}\end{aligned}  
         \label{eq: expresión de las concentraciones de los minoritarios en los puntos analogos}
\end{equation}
Obteniendo un resultado análogo al del equilibrio
\begin{equation}
    V_\mathrm{B}-V=\frac{k_\mathrm{B}T}{e} \ln \frac{p_p}{p_n}=\frac{k_\mathrm{B}T}{e} \ln \frac{n_n}{n_p}
\end{equation}
\newpage
Por lo que obtenemos las expresiones para las concentraciones de los portadores minoritarios en los bordes de la ZCE como funciones de las concentraciones en el equilibrio, el potencial al que sometemos el semiconductor y la temperatura.
\begin{equation}
    p_n=p_{n0} \exp{\frac{eV}{k_\mathrm{B}T}} \hspace{2cm} n_p=n_{p0} \exp{\frac{eV}{k_\mathrm{B}T}}
    \label{eq: Relaciones de Boltzmann}
\end{equation}
Que son las conocidas como \textbf{Relaciones de Boltzmann} que nos servirán para resolver la ecuación de continuidad en las zonas neutras.
\subsubsection{Distribuciones de portadores minoritarios en las zonas neutras}
\begin{figure}[h]
\centering
\includegraphics{Standalone/T3_11_concenctraciones_portadores_minoritarios.pdf}
    \caption{Concentración de portadores minoritarios en una unión PN en polarización directa. Apreciamos que cuando nos alejamos de la ZCE la densidad de portadores se vuelve constante. }
    \label{fig: concenctraciones portadores minoritarios}
\end{figure}
Queremos ahora estudiar cómo se distribuyen los portadores de carga minoritarios en las zonas neutras, para ello volvemos a utilizar la ecuación de continuidad \ref{eq: continuidad huecos}, en este caso nos encontramos en $x>x_n$, trabajamos sobre la concentración y la ecuación es 
\begin{equation}
     \diff[2]{p_n}{x}=\frac{p_n}{L_p^2}-\frac{p_{n0}}{L_p^2}
     \label{eq: portadores minoritarios zona neutra N}
\end{equation}
Que tiene como solución la ecuación \ref{eq: Cont Generacion en volumen}, donde $L_p=\sqrt{D_p\tau_p}$ es la \textbf{longitud de difusión} ahora vamos a establecer condiciones de contorno para obtener la solución final de la distribución. Si utilizamos la condición de \textbf{diodo largo} $L_\mathrm{Zona Neutra} \gg L_\mathrm{eq}$.
\begin{equation}
    p_n(x)=p_{n0}+\cancelto{0 \text{ ya que } p_n(x)|_{x\xrightarrow[]{} \infty}=p_{n0}}{A}e^{x/L_p}+Be^{-x/L_p} \Rightarrow  \hspace{1cm} p_n(x)=Be^{-x/L_p}
\end{equation}
Ahora, si calculamos dicha ecuación en el punto $x=x_n$
\begin{equation*}
    p_n(x_n)|_{x=x_n}=p_n(x_n)=p_{n0} \exp{\frac{eV}{k_\mathrm{B}T}}\xrightarrow{} B=p_{n0}\left(\exp{\left(\frac{eV}{k_\mathrm{B}T}\right)-1} \right) \exp{\frac{x_n}{L_p}}
\end{equation*}
Con lo que tenemos la siguiente expresión para la distribución de los portadores minoritarios
\begin{equation}
    p_n(x)-p_{n0}=p_{n0} \left( \exp{\left( \frac{eV}{k_\mathrm{B}T}\right)}-1  \right) \exp{-\frac{x-x_n}{L_p}}
    \label{eq: Distribucion de portadores minoritarios, huecos}
\end{equation}
Análogamente obtenemos la expresión de los electrones del lado P ($x>x_p$)
\begin{equation}
    n_p(x)-n_{p0}=n_{p0} \left( \exp{\left( \frac{eV}{k_\mathrm{B}T}\right)}-1  \right) \exp{\frac{x-x_p}{L_n}}
    \label{eq: Distribucion de portadores minoritarios, electrones}
\end{equation}
\newpage
Por último vamos a considerar que la conocida como condición de \textbf{cuasi-neutralidad} que la diferencia entre el exceso de portadores y la densidad de portadores en el equilibrio es constante en las zonas neutras
\begin{equation}
    \begin{aligned}
        &n_n(x)-n_{n0} \cong p_n(x)-p_{p0} \ \ \text{Zona N} \\
        &n_n(x)-n_{n0} \cong p_n(x)-p_{p0} \ \ \text{Zona P}
    \end{aligned}
\end{equation}
\begin{figure}[h]
    \centering
    \begin{subfigure}[h]{0.4\textwidth}
    \includegraphics{Standalone/T3_12_Concentracion_de_portadores_frente_a.pdf}
    \caption{Concentración de portadores frente a la posición en polarización directa.}
    \end{subfigure}
    \hspace{15mm}
    \begin{subfigure}[h]{0.4\textwidth}
        \includegraphics{Standalone/T3_13_Concentracion_de_portadores_frente_a.pdf}
    \caption{Concentración de portadores frente a la posición en polarización indirecta.}
    \end{subfigure}
    \caption{Concentración de portadores frente a la posición.}
    \label{fig:enter-label}
\end{figure}
\subsubsection{Cálculo de la corriente}
En la ZCE, consideraremos ausencia de generación recombinación por lo que podemos considerar que la densidad de corriente es \textbf{constante}, para ello debemos considerar
\begin{itemize}
    \item Suponemos que el campo eléctrico en los límites de ZCE se hace nulo. Por lo que podemos despreciar la contribución de la corriente de arrastre de los minoritarios en las zonas neutras.
    \item La corriente total es la suma de $\mathbf{J}_p$ y $\mathbf{J}_n$ que son constantes en la ZCE.
\end{itemize}
Vamos a calcular la corriente en las zonas neutras, para ello utilizamos la ecuación \ref{eq: correinte de difusion}, como de costumbre utilizaremos un modelo unidimensional
\begin{equation}
    \begin{aligned}
        &\text{Zona P   }J_n(x)=eD_n\diff[]{n_p(x)}{x}= \frac{eD_nn_{p0}}{L_n} \left( \exp{\left( \frac{eV}{k_\mathrm{B}T}\right)}-1  \right) \exp{\frac{x-x_p}{L_n}}\\\\
        & \text{Zona N   }J_p(x)=eD_p\diff[]{n_p(x)}{x}= \frac{eD_pn_{p0}}{L_p} \left( \exp{\left( \frac{eV}{k_\mathrm{B}T}\right)}-1  \right) \exp{\frac{x-x_p}{L_p}}
    \end{aligned}
    \label{eq: Corrientes de difusion zonas neutras}
\end{equation}
La corriente total será entonces
\begin{equation}
J=\underbrace{e\left[\textcolor{red}{\frac{n_{p o} D_n}{L_n}}+\textcolor{blue}{\frac{p_{n o} D_p}{L_p}}\right]}_{J_0=\textcolor{red}{J_{n0}}+\textcolor{blue}{J_{p0}}}\left( \exp{\left( \frac{eV}{k_\mathrm{B}T}\right)}-1  \right)=J_o\left( \exp{\left( \frac{eV}{k_\mathrm{B}T}\right)}-1  \right)
\end{equation}
Siendo $J_0$ la \textbf{corriente inversa de saturación.}
\newpage
\begin{wrapfigure}{r}{0.3\textwidth}
    \includegraphics{Standalone/T3_14_Corriente_ideal_PN.pdf}
    \caption{Corriente ideal de una unión PN.}
    \label{fig: Corriente ideal PN}
\end{wrapfigure}
Hasta ahora hemos considerado que la union PN era simétrica, es decir que el número de impurezas dondadoras y aceptadoras era más o menos igual.
\begin{itemize}
    \item Si tenemos mayor número de impurezas aceptadoras, la corriente en la unión es debida a los huecos.
    \item Si tenemos mayor número de impurezas donadoras, la corriente es debida a los electrones.
\end{itemize}
\noindent Ya hemos visto cómo son las corrientes en las zonas neutras de los portadores minoritarios. Si recordamos, la corriente total es constante por lo que para calcular la corriente de los portadores mayoritarios en las zonas neutras únicamente tenemos que restar las corrientes de los minoritarios a la corriente total.
\begin{wrapfigure}{r}{0.4\textwidth}
    \centering
    \includegraphics{Standalone/T3_15_Componentes_de_la_corriente_PN.pdf}
    \caption{Componentes ideales de la corriente en una unión PN.}
    \label{fig:Componentes de la corriente PN}
\end{wrapfigure}

\subsubsection{\textcolor{red}{Cuasiniveles de Fermi}}
\noindent Como vimos en la sección \ref{sec: Union PN en equilibrio}, el nivel de Fermi se mantiene constante respecto a la posición, en todo el cristal. Sin embargo cuando salimos del equilibrio, el nivel de Fermi sí que depende de la posición, por lo que tendremos que volver a hacer uso de los cuasiniveles de Fermi para poder calcular la corriente. En este caso, los niveles de Fermi sí que dependen de la posción, es decir si hacemos el cálculo para huecos tenemos un nivel y si hacemos el cálculo para electrones tendremos otro nivel. El cálculo es análogo a lo calculado en la sección \ref{sec: Union PN en equilibrio} para electrones
\begin{equation}
    n(x)=N_\mathrm{c}\exp{\left[ \frac{E_\mathrm{FN}(x)-E_\mathrm{C}(x)}{k_\mathrm{B}T}\right]}
    \label{eq: Concentracion de portadores cuasiniveles de Fermi}
\end{equation}
Para huecos
\begin{equation}
    p(x)=N_\mathrm{v}\exp{\left[ \frac{E_\mathrm{V}(x)-E_\mathrm{FP}(x)}{k_\mathrm{B}T}\right]}
    \label{eq: Concentracion de portadores cuasiniveles de Fermi}
\end{equation}
En la figura \ref{fig:Diagramas de energia PN y cuasiniveles} están representados los niveles de Fermi. Supongamos que estamos en el régimen de débil nivel de inyección, que para los minoritarios sí que suponem variaciones considerables mientras que para los mayoritarios apenas varían. Además apreciamos que en la Figura, los cuasiniveles son constantes.
\clearpage
\begin{figure}[h]
    \begin{subfigure}[h]{0.4\textwidth}
    \includegraphics{Standalone/T3_16_polarizacion_inversa.pdf}
    \caption{Polarización inversa.}
    \label{fig: polarizacion inversa}
    \end{subfigure}
    \begin{subfigure}[h]{0.4\textwidth}
\hspace{20mm}
    \includegraphics{Standalone/T3_17_polarizacion_directa.pdf}
    \caption{Polarización directa.}
    \label{fig: polarizacion directa}
    \end{subfigure}
    \caption{Diagramas de energías asociados a una unión PN.}
    \label{fig:Diagramas de energia PN y cuasiniveles}
\end{figure}

\section{Unión PN. Comportamiento dinámico}
\subsection{Modelo de pequeña señal}
Para poder realizar este modelo vamos a hacer las siguientes suposiciones
\begin{itemize}
    \item La respuesta a la pequeña señal siempre es \textbf{lineal}.
    \item La unión PN siempre tendrá \textbf{polarización directa}.
    \item Estamos en régimen de baja frecuencia, por lo que trataremos el problema como una sucesión de regímenes estacionarios.
    \item La señal que llega a la unión PN es es sinusoidal. $v(t)=V_0+v_i(t)=V_0+V_i\exp{j\omega t}$ y $J(t)=J_0+J_i\exp{j\omega t}$ 
\end{itemize}
Como condición inicial del problema tenemos alguna de las características estáticas y buscamos hayar $J(\omega)$ para poder calcular la impedancia.\\\\
A priori esperamos que haya dos contribuciones, una \textbf{contribución resistiva} y otra \textbf{contribución capacitiva.}
A partir de la ecuación de continuidad dependiente del tiempo obtenemos
\begin{equation}
\frac{1}{Z(\omega)}=Y(\omega)=\frac{J_1}{V_i}=\frac{e^2}{k_\mathrm{B}}\left[\frac{D_p p_{n 0}}{L_p}\left(1+j \omega \tau_p\right)^{1 / 2}+\frac{D_n n_{p 0}}{L_n}\left(1+j \omega \tau_n\right)^{1 / 2}\right] \exp{\frac{e V_0}{k_\mathrm{B} T}}
\label{eq: Admitancia}
\end{equation}
Obtenemos la solución general
\begin{equation}
Y(\omega)=G_{\mathrm{d} p}\left(1+j \omega \tau_p\right)^{1 / 2}+G_{\mathrm{d} n}\left(1+j \omega \tau_n\right)^{1 / 2} \end{equation}
\begin{equation*}
G_{\mathrm{d} p}=\frac{e^2}{k_\mathrm{B}} \frac{D_p p_{n 0}}{L_p} \exp{\frac{e V_0}{k_\mathrm{B}}} \quad G_{\mathrm{d} n}=\frac{e^2}{k_\mathrm{B}} \frac{D_n n_{p 0}}{L_n} \exp{\frac{e V_o}{k_\mathrm{B} T}}
\end{equation*}
\clearpage
\begin{wrapfigure}{r}{0.35\textwidth}
    \includegraphics{Standalone/T3_18_Corriente_ideal_y_real_PN.pdf}
    \caption{Corriente \textcolor{red}{ideal} y corrriente \textcolor{blue}{real} de una unión PN.}
    \label{fig: Corriente ideal y real PN}
\end{wrapfigure}


Reordenando términos obtenemos que la admintancia, tiene una parte real, que depende del punto de operación, asociado a la \textbf{contribución resistiva} y otro imaginario asociado a la \textbf{contribución capacitiva}
\begin{equation}
    Y(\omega)=\frac{1}{R_\mathrm{d}\left(\omega, V_o\right)}+j \omega C_\mathrm{d}\left(\omega, V_0\right)
    \label{eq: admitancia general}
\end{equation}
Si nos vamos al régimen de baja frecuencia, los coeficientes dejan de ser dependientes de la frecuencia y obtenemos $Y=G_\mathrm{d}+j\omega C_\mathrm{d}$
\begin{equation}
 \left.G_\mathrm{d}=\frac{1}{R_\mathrm{d}}=\frac{e^2}{K_B T}\left(\frac{D_p p_{n 0}}{L_p}+\frac{D_n n_{p 0}}{L_n}\right) e^{e V_0 / k_\mathrm{B}T}=\frac{\mathrm{d} J}{\mathrm{d} V}\right)_{V=V_0}
 \label{eq: conductancia }
 \end{equation}
 Que la pendiente de la característica $J-V$ en el punto de operación.
 \begin{equation}
 C_\mathrm{d}=\frac{e^2}{2 k_\mathrm{B} T}\left(\frac{D_p p_{n 0}}{L_p} \tau_p+\frac{D_n n_{p 0}}{L_n} \tau_n\right) e^{e V_0 / k_\mathrm{B} T}    
 \label{eq: Capacidad de difusion}
 \end{equation}
La capacidad de difusión que representa la variación de cargas en zonas neutras.
\subsubsection{\textcolor{red}{Circuito equivalente de pequeña señal y baja frecuencia}}
Como hemos visto, nuestro diodo, tiene asociadas contribuciones resistivas y capacitivas, esto hace que a la hora de estudiarlo en un circuito, podamos por reemplazarlo por una resistencia y dos condensadores conectados en paralelo.
\begin{itemize}
    \item \textbf{Capacidad de transición}: Que se da en la ZCE y que está representada por la ecuación \ref{eq: Capacidad de transicion}
    \item \textbf{Capacidad de difusión}: Que se encuentra en las zonas neutras \ref{eq: Capacidad de difusion}
    \item \textbf{Resistencia incremental}: Asociada al diodo y que corresponde al inverso de la ecuación \ref{eq: conductancia }
    \item Además podemos asociar una resistencia en serie asociada a las zonas neutras y a los contactos que será nuestro circuito equivalente real.    
\end{itemize}


\subsection{\textcolor{red}{Desviación del diodo real respecto al ideal}}
\noindent La figura \ref{fig: Corriente ideal PN} es un caso ideal, experimentalmente observamos que existe un potencial de rotura $V_\mathrm{BR}$ como se muestra en la figura \ref{fig: Corriente ideal y real PN}, a partir de este potencial, comenzamos a ver corrientes dadas por efecto tunel y por ionización por impacto de los átomos del semiconductor.
\clearpage

\subsubsection{\textcolor{red}{Procesos reales en polarización inversa}}
\begin{itemize}
    \item En el diodo real, sí que tenemos procesos de generación recombinación en la ZCE, la generación de pares, aparece cuando existe un campo eléctrico demasiado grande en los bordes, y en esta situación domina la corriente de arrastre por lo que contribuye a la corriente total.
    \item Cuando aplicamos potenciales altos, la separación entre las bandas de energía disiminuye (Figura \ref{fig: polarizacion inversa}) por lo que los electrones pueden pasar por efecto tunel de una banda de energía a otra, lo que contribuye a la corriente y nos genera la ruptura.
    \item A voltajes altos los electrones tienen mucha energía por lo que podrán colisionar con átomos de la red cristalina, ionizándolos y generando pares electrón hueco que hace que nos aumente la corriente.
\end{itemize}
\subsubsection{\textcolor{red}{Procesos reales en polarización directa}}
\begin{itemize}
    \item Si sometemos el diodo a voltajes positivos y no muy altos, el voltaje cae en la ZCE y apenas varía en las zonas neutras, tenemos excesos de portadores, por lo que aumenta la contribución por recombinación que es mayor a la ideal.
    \item Si aumentamos el voltaje, encontramos una zona que es prácticamente igual a la ideal, que denominamos \textbf{zona de cuasi-idealidad}.
    \item Cuando tenemos voltajes altos, comenzamos a entrar en el régimen de alta inyección por lo que los procesos de generación recombinación son más ineficientes lo que hace que perdamos intensidad, a estos procesos los podemos asociar una \textbf{resistencia en serie}. A voltajes muy muy altos, el diodo se comporta como una resistencia.
\end{itemize}

\section{Unión metal-semiconductor} \label{sec: Union MSC}
Hasta ahora hemos considerado que las zonas P y N estaban unidas como si fueran una única estructura, en una unión homogénea. Ahora bien, también podemos establecer una unión a través de un metal, conocida como \textbf{unión metal-semiconductor}.\\\\
\noindent Dentro de las uniones MSC, tenemos dos tipos de unión: \textbf{unión óhmica}, que apenas presenta resistencia frente al paso de corriente en cualquiera de los sentidos, por lo que se usa en los contactos de los dispositivos; y \textbf{unión rectificadora}, que funcina como una unión PN.
\\\\
Vamos a analizar las figura \ref{fig: Union MSC sin contacto} y a establecer las condicoines que tiene esta unión:
\begin{itemize}
    \item \textbf{Nivel de vacío}, es el nivel que tomaremos como referencia en estas uniones.
    \item \textbf{Función de trabajo} $\phi_\mathrm{M,SC}$, medida en voltios, nos define la energía mínima necesaria para arrancar un electrón de un sólido.
    \item \textbf{Afinidad electrónica} $\chi$, medida en eV's, que representa la energía que cede un electrón cuando llega a la capa de conducción.
\end{itemize}
\begin{figure}[h!]
    \centering
    \includegraphics{Standalone/T3_19_Union_MSC_sin_contacto.pdf}
    \caption{Diagrama de bandas de una unión metal semiconductor, antes de entrar en contacto. Los estados ocupados estan representados en las zonas grises, los vacíos en las vacías.}
    \label{fig: Union MSC sin contacto}
\end{figure}
\subsection{\textcolor{red}{Unión metal-semiconducor tipo N}}
\begin{wrapfigure}{r}{0.5\textwidth}
    \includegraphics{Standalone/T3_20_Union_metalsemiconductor_tipo_N.pdf}
        \caption{Diagrama de bandas de energía de una unión metal semiconductor tipo N.}
    \label{fig: Union metalsemiconductor tipo N}
\end{wrapfigure}
Antes del contacto el nivel de Fermi del semiconductor se encuentra por encima del nivel de Fermi del metal $\phi_\mathrm{M}>\phi_\mathrm{SC}$. Al hacer contacto, queremos que nuestro nivel de Fermi sea constante en todo el sistema cuando estemos en equilibrio térmico. Esto implica que los electrones del semiconductor, fluyan hasta el estado de energía menos energético de la banda de conducción.\\\\
Este flujo de electrones hace que, en el semiconductor, tengamos un un exceso de átomos cargados positivamente en una zona cercana a la unión.
\\\\
\noindent En la figura \ref{fig: Union metalsemiconductor tipo N} vemos que las bandas de energía se están doblando hacia arriba, esto es debido a que los electrones tienen cierta dificultad para pasar del metal al semiconductor por la presencia de las cargas positivas en los límites del semiconductor, ya que una vez llegan al límite de la zona cargada positivamente, tienden a permanecer estáticas. 
\\\\
Supondremos también que en la zona de carga espacial tendremos \hyperref[sec: Vaciamiento total]{vaciamiento total}. Además como nuestro nivel de referencia está en el nivel de vacío, el metal se encuentra a $-V_\mathrm{B}$ respecto al semiconductor. Los cálculos de $\rho, \ E, \ V, \ C$ son análogos a la unión $\mathrm{P}^+\mathrm{N}$, por lo tanto la anchura de la zona de carga espacial será
\begin{equation}
    d\equiv x_\mathrm{n}=\left(\frac{2\epsilon V_\mathrm{B}}{eN_\mathrm{D}} \right)^{1/2}
    \label{ec: Anchura de la ZCE unión P+N sin polarizar}
\end{equation}
\subsubsection{Polarización de la unión} \label{sec: Polarizacion Union metal semiconductor tipo N}
\noindent El criterio de signos que vamos a seguir para establecer la polarización de la corriente será el siguiente
\begin{table}[h]
\centering
\begin{tabular}{|c|c|c|}
    \hline
        Metal $(\mathrm{P}^+)$ & SC  (N)& Polarización\\
        \hline
        \hline
         -& + & $V_\mathrm{B}$\\
         +&  -& directa $(V>0)$\\
         -&  +& inversa $(V<0)$\\
         \hline
    \end{tabular}    
     \caption{Criterio de signos}
    \label{tab:Criterio de signos para la polarización}
\end{table}\\
\noindent Cuando la unión tiene polarización \textbf{directa}, vemos que la barrera de potencial sube respecto al potencial de contacto, por lo que la curvatura que tiene la barrera se atenúa como se muestra en la figura \ref{fig: Union MSC-N polarizacion directa}.\\\\
A esta unión con polarización directa, se le conoce como \textbf{unión rectificadora}, y se utiliza en los conocidos como \textbf{diodos Shottky}.
\begin{equation}
    J=J_0 \exp{\left(\frac{eV}{k_\mathrm{B}T}-1 \right)}
    \label{eq: union rectificadora}
\end{equation}

\begin{wrapfigure}{r}{0.5\textwidth}
        \includegraphics{Standalone/T3_21_Union_MSC_N_polarizacion_directa.pdf}
    \caption{Niveles de energía de una unión Metal-SC-N con polarización directa}
    \label{fig: Union MSC-N polarizacion directa}
\end{wrapfigure}

\noindent Esta unión elimina la capacidad de difusión que sí que presentaba (y dominaba) en la unión PN con polarización directa ya que no hay presentes corrientes de difusión en la unión, únicamente tenemos corrientes de arrastre.\\\\
Esto es posible gracias a la \textbf{emisión termoiónica}, que permite generar electrones únicamente ligados al potencial $\phi_\mathrm{M}$ que, nos dan corrientes de arrastre al ser movidos por el campo eléctrico externo. 
La corriente de difusión es prácticamente nula debido a la presencia del metal que actúa como un sumidero de portadores.\\\\
\noindent Hasta ahora hemos considerado uniones con $\phi_\mathrm{M}> \phi_\mathrm{SC}$. Si consideramos ahora $\phi_\mathrm{M}< \phi_\mathrm{SC}$ el diagrama de bandas antes de la unión es la representada en la figura \ref{fig: Union MSC sin contacto para phiM <phiSC}.
\clearpage
\begin{figure}[h!]
    \centering
    \includegraphics{Standalone/T3_22_Union_MSC_sin_contacto_para.pdf}
    \caption{Diagrama de bandas de una unión metal semiconductor, antes de entrar en contacto.}
    \label{fig: Union MSC sin contacto para phiM <phiSC}
\end{figure}

El análisis es análogo al caso anterior pero hay una difereencia, esta vez las bandas de energía se doblan hacia abajo, en lugar de hacia arriba.\\

\begin{wrapfigure}{r}{0.5\textwidth}
    \includegraphics{Standalone/T3_23_Union_metalsemiconductor_tipo_N_curva.pdf}
        \caption{Diagrama de bandas de energía de una unión metal semiconductor tipo N.}
    \label{fig: Union metalsemiconductor tipo N, curva hacia abajo}
\end{wrapfigure}
\noindent En equilibrio térmico, el nivel de Fermi sigue siendo constante por lo que cuando tenemos contacto, las bandas tienen la forma que se muestra \ref{fig: Union metalsemiconductor tipo N, curva hacia abajo}.\\\\ El metal sigue teniendo carga superficial, pero el semiconductor tiene carga libre cercana a su superficie ya que no tenemos una barrera de potencial para pasar del semiconductor al metal. Esto hace que las barreras en la superficie de contacto no cambien al polarizar, como sí pasaba en el caso de la figura \ref{fig: Union MSC-N polarizacion directa}.
\subsubsection{Polarización de la Unión}
\noindent Las barreras de potencial son constantes en la región cercana a la superficie, lo que implica que la concentración de portadores en la superficie sea constante y que se encuentre \textbf{siempre en equilibrio}. Esta unión apenas ofrece resistencia al paso de corriente en cualquiera de los sentidos.
\begin{figure}[h!]
    \centering
    \begin{subfigure}[h]{0.4\linewidth}
    \centering
        \includegraphics{Standalone/T3_24_Polarizacion_directa.pdf}
        \caption{Polarización directa}
    \end{subfigure}
    \begin{subfigure}[h]{0.4\textwidth}
    \centering
            \includegraphics{Standalone/T3_25_Polarizacion_inversa.pdf}
        \caption{Polarización inversa}
    \end{subfigure}
    \caption{Polarización de contacto óhmico.}
\end{figure}
\newpage
\noindent Esta geometría de las bandas genera estados superficiales que capturan electrones que vacían el semiconductor. Para solucionar este problema se opta por hacer que la unión sea gradual, para ello podemos alear el metal y el semiconductor; ó impurificar mucho el semiconductor, de tal manera que sólo tenemos contribución de los electrones (mayoritarios).
\subsection{\textcolor{red}{Unión metal-semiconductor tipo P} }\label{sec: Union MSC-P}
Volvamos a considerar que $\phi_\mathrm{M}<\phi_\mathrm{SC}$, antes de entrar en contacto vemos que tenemos una situación análoga a la anterior, pero en este caso vamos a considerar huecos.
\begin{figure}[h!]
    \centering
    \includegraphics{Standalone/T3_26_Union_MSC_sin_contacto.pdf}
    \caption{Diagrama de bandas de una unión metal semiconductor, antes de entrar en contacto.}
    \label{fig: Union MSC sin contacto}
\end{figure}

\noindent Cuando ponemos en contacto el metal con el semiconductor tipo P, y estamos en equilibrio térmico, el nivel de Fermi es constante y tendremos exceso de concentración de cargas positivas en la superficie del metal y de cargas negativas en el semiconductor. Por lo que tendremos un diagrama de bandas similiar al de la figura \ref{fig: Union metalsemiconductor tipo N, curva hacia abajo}.\\

\begin{wrapfigure}{r}{0.5\textwidth}
    \includegraphics{Standalone/T3_27_Union_metalsemiconductor_tipo_P_curva.pdf}
        \caption{Diagrama de bandas de energía de una unión metal semiconductor tipo P.}
    \label{fig: Union metalsemiconductor tipo P, curva hacia abajo}
\end{wrapfigure}

En el metal tenemos densidad de carga superficial mientras que en el semiconductor tenemos densidad de carga volúmica fija. Vamos a seguir considerando vaciamiento total en el semiconductor. El metal se encuentra a $+V_\mathrm{B}$, respecto al nivel de vacío,que es el potencial de contacto. El cálculo de $\rho, \ E, \ V, \ C$ es análogo a los de la unión P$\mathrm{N}^+$.\\\\
Como estamos tratando con huecos, para pasar del semiconductor al metal, tienen que tener una energía $eV_\mathrm{B}$ mientras que para pasar del semiconductor al metal, como el extremo de la banda de valencia está por debajo de la banda de conducción, tendrá que tener una energía $E_\mathrm{SC}-\phi_\mathrm{M}=\chi_\mathrm{SC}+E_\mathrm{g}-\phi_\mathrm{M}$.
\clearpage
\subsubsection{Polarización de la Unión}
El criterio de signos que vamos a utilizar en este caso es 
\begin{table}[h!]
\centering
\begin{tabular}{|c|c|c|}
    \hline
        Metal $(\mathrm{N}^+)$ & SC  (P)& Polarización\\
        \hline
        \hline
         +& - & $V_\mathrm{B}$\\
         -&  +& directa $(V>0)$\\
         +&  -& inversa $(V<0)$\\
         \hline
    \end{tabular}    
     \caption{Criterio de signos}
    \label{tab:Criterio de signos para la polarización union mscp}
\end{table}

\begin{wrapfigure}{r}{0.5\textwidth}
    \includegraphics{Standalone/T3_28_Union_metalsemiconductor_tipo_P_curva.pdf}
        \caption{Diagrama de bandas de energía de una unión metal semiconductor tipo P.}
    \label{fig: Union metalsemiconductor tipo P, curva hacia abajo}
\end{wrapfigure}

Si tenemos polarización directa, sólo la barrera para pasar del metal al semiconductor permanece constante mientras que la barrera para pasar del semiconductor al metal cambia a $e(V_\mathrm{B}-V)$. En este caso las bandas presentan la misma estructura que en la figura \ref{fig: Union metalsemiconductor tipo P, curva hacia abajo} solo que la curvatura es menor debido a la presencia del nuevo potencial. Además aparecen dos niveles de Fermi, uno en el semiconductor y otro en el metal.\\\\
\noindent La unión es rectificadora y sigue siendo la ecuación \ref{eq: union rectificadora} \\
\begin{center}
    \rule{8cm}{2pt} 
\end{center}

\noindent Vamos ahora a considerar que $\phi_\mathrm{M}>\phi_\mathrm{SC}$ como se ha hecho anteriormente, el diagrama de bandas antes del contacto es 
\begin{figure}[h!]
    \centering
    \includegraphics{Standalone/T3_29_Union_MSC_tipo_P_sin.pdf}
    \caption{Diagrama de bandas de una unión metal semiconductor, antes de entrar en contacto.}
    \label{fig: Union MSC tipo P sin contacto phim>phisc}
\end{figure}
\clearpage
\begin{wrapfigure}{r}{0.5\textwidth}
    \includegraphics{Standalone/T3_30_Union_metalsemiconductor_tipo_P_en.pdf}
        \caption{Diagrama de bandas de energía de una unión metal semiconductor tipo P en equilibrio.}
    \label{fig: Union metalsemiconductor tipo P en equilibrio}
\end{wrapfigure}

\noindent En el equilibrio, las barreras se doblan hacia abajo y tienen la forma que se muestra en la figura \ref{fig: Union metalsemiconductor tipo P en equilibrio}.\\\\
Las barreras en este caso se mantienen constantes cuando polarizamos el diodo. A este tipo de unión se le conoce como \textbf{unión óhmica}.\\\\\\\\\\\\\\
\section{Unión real}
\noindent En una unió real, tenemos estados superficiales, este tipo de estados tienen ocupación variable tanto en equilibrio como con polarización; además si no se realiza correctamente la unión, no podremos conocer el tipo de unión que tenemos en el diodo.\\\\
Cuando tenemos uniones óhmicas, utilizaremos un semiconductor muy impurificado ($\phi_\mathrm{M}>\phi_\mathrm{SC}$ y SC-N), que nos hará que funcione como un rectificador de corriente., si tenemos que la función de trabajo del metal es mayor que la del semiconductor, la zona de vaciamiento y la barrera que hay en las bandas es muy estrecha por lo que el efecto túnel tiene relevancia a la hora del paso de la corriente en ambas direcciones.
\pagestyle{fancy}
\fancyhead[R]{Electrónica Física: Tema 4}
\fancyhead[L]{4º Grado en Física}
\chapter{Transistor Bipolar}
Determinadas las características de las uniones PN y las uniones metal-semiconductor, podemos construir diodos, que son dispositivos hechos por una única unión entre un semiconductor tipo P y otro tipo N. Los \textbf{transistores} son dispositivos formados por varias uniones entre semiconductores y metal-semiconductor.\footnote{ Este capítulo sólo estudiaremos el transistor bipolar, mientras que en el siguiente estudiaremos el MOSFET.}
\section{Transistor bipolar} \label{sec: Transistor bipolar}
El transistor bipolar es un dispositivo electrónico que está formado por tres regiones diferenciadas, con dos uniones PN. A estas regiones se les conoce como \textbf{terminales} \footnote{Vamos a desarrollar un transistor bipolar PNP pero todo es análogo al NPN.}, estas regiones están dopadas con distintas concentraciones (se justificarán los distintsos motivos de los dopajes a lo largo del capítulo), además supondremos que tenemos \textbf{uniones abruptas}.
\begin{itemize}
    \item \textbf{Emisor} (E): Esta región es la que tiene mayor dopaje y es la que nos aporta los electrones.
    \item \textbf{Base} (B): Es la región con dopaje intermedio y nos controla la transición de los electrones entre emisor y colector.
    \item \textbf{Colector} (C): Es la región con menor dopaje. Recoge los portadores de carga emitidos por el emisor y es responsable de la corriente de salida.
\end{itemize}

\begin{figure}[h]
    \centering
    \begin{subfigure}[h]{0.3\linewidth}
    \centering
        \includegraphics{Standalone/T4_01_Transistor_NPN.pdf}
        \caption{Transistor NPN.}
        \label{fig: Transistor NPN}
    \end{subfigure}
\hfill
    \begin{subfigure}[h]{0.3\linewidth}
    \centering
    \includegraphics{Standalone/T4_02_Transistor_PNP.pdf}
            \caption{Transistor PNP.}
            \label{fig: Transistor PNP}
    \end{subfigure}
    \hfill
    \begin{subfigure}[h]{0.3\linewidth}
        \includegraphics{Standalone/T4_03_Concentraciones_de_portadores_en_un.pdf}
        \caption{Concentraciones de portadores en un transistor PNP.}
    \end{subfigure}
    \caption{Geometría de un transistor bipolar}
    \label{fig:Geometrias transistor}
\end{figure}
\newpage
\subsection{Funcionamiento cualitativo}
Ahora seguiremos desarrollando para un transistor PNP. En el equilibrio, el transistor bipolar presenta la siguiente estructura.

\begin{figure}[h]
    \centering
    \begin{subfigure}[h]{0.3\linewidth}
    \centering
        \includegraphics{Standalone/T4_04_Distribucion_espacial_equilibrio.pdf}
        \caption{Distribución espacial del transistor.}
        \label{fig: Distribucion espacial equilibrio}
    \end{subfigure}
\hfill
    \begin{subfigure}[h]{0.3\linewidth}
    \centering
    \includegraphics{Standalone/T4_05_Diagrama_de_bandas_equilibrio.pdf}
    \caption{Diagrama de bandas de energía.}
    \label{fig: Diagrama de bandas equilibrio}
    \end{subfigure}
    \hfill
    \begin{subfigure}[h]{0.3\linewidth}
        \includegraphics{Standalone/T4_06_Distribucion_concentraciones_equilibrio.pdf}
        \caption{Concentraciones de portadores en un transistor PNP.}
        \label{fig: Distribucion concentraciones equilibrio}
    \end{subfigure}
    \caption{Distribuciones y geometrías de un transistor bipolar PNP en equilibrio.}
    \label{fig: Distribuciones de un transistor equilibrio}
\end{figure}
\noindent La figura \ref{fig: Distribuciones de un transistor equilibrio} nos muestra las distribuciones espacial, de bandas de energía y de concentraciones de un transistor tipo PNP. Ahora bien, expliquemos la figura \ref{fig: Distribucion espacial equilibrio}
\begin{itemize}
    \item Las líneas \textcolor{red}{\textbf{rojas punteadas}}, nos indican los límites de los distintos terminales del transistor.
    \item Las líneas \textcolor{blue}{\textbf{azules punteadas}}, nos indican los límites de las zonas de carga espacial en la base y nos delimitan la \textbf{zona neutra de la base}.
    \item Las líneas \textbf{negras punteadas}, nos indican los límites de las zonas de carga espacial en el emisor y el colector respectivamente.
    \item Realizamos la aproximación de \hyperref[sec: Vaciamiento total]{\textbf{vaciamiento total}} en las ZCE's.
\end{itemize}
Vemos que los límites de los terminales los de la zona neutra de la base no coinciden. Esto lo solucionamos haciendo una asimetría entre las concentraciones de los terminales ya que nos hacen más repentinas las caídas/subidas que hemos visto en las figuras \ref{fig: Diagrama de bandas equilibrio} y \ref{fig: Distribucion concentraciones equilibrio}. Es decir, \textbf{la concentración de portadores, afecta a la anchura en las ZCE's}, concretamente es inversamente proporcional; a mayor dopaje, menor será la anchura \footnote{Anchura efectiva}.
\subsubsection{\textcolor{red}{Unión bajo polarización}}
Si existe una diferencia de potencial entre el colector y el emisor, ya no nos encontramos en equilibrio y la geometría de las uniones del transistor hace que aparezcan distintas polarizaciones en cada una de las uniones.\\\\
Si polarizamos un transistor PNP, $V_\mathrm{C}-V_\mathrm{E}=V_\mathrm{CE}<0$, la unión entre \textbf{emisor} y \textbf{base}, tiene polarización \textbf{directa} mientras que la unión \textbf{base}, \textbf{colector}, tiene polarización \textbf{inversa}.
\newpage
\begin{figure}[h]
    \centering
    \begin{subfigure}[h]{0.3\linewidth}
    \centering
        \includegraphics{Standalone/T4_07_Distribucion_espacial_polarizacion.pdf}
        \caption{Distribución espacial del transistor.}
        \label{fig: Distribucion espacial polarizacion}
    \end{subfigure}
\hfill
    \begin{subfigure}[h]{0.3\linewidth}
    \centering
    \includegraphics{Standalone/T4_08_Diagrama_de_bandas_polarizacion.pdf}
    \caption{Diagrama de bandas de energía.}
    \label{fig: Diagrama de bandas polarizacion}
    \end{subfigure}
    \hfill
    \begin{subfigure}[h]{0.3\linewidth}
        \includegraphics{Standalone/T4_09_Distribucion_concentraciones_polarizacion.pdf}
        \caption{Concentraciones de portadores en un transistor PNP.}
        \label{fig: Distribucion concentraciones polarizacion}
    \end{subfigure}
    \caption{Distribuciones y geometrías de un transistor bipolar PNP bajo polarización\\ $V_\mathrm{CE}<0$.}
    \label{fig: Distribuciones de un transistor polarizacion}
\end{figure}

\noindent Si consideramos que las uniones son independientes tenemos $W\gg L_{p\mathrm{B}}$, mientras que si no son independientes $W \ll L_{p\mathrm{B}}$\footnote{Linea punteada roja de la figura \ref{fig: Distribucion concentraciones polarizacion}.}.\\\\
Los diagramas de la figura \ref{fig: Distribuciones de un transistor polarizacion}, nos muestran que en un transistor PNP, los huecos inyectados por el emisor en la base llegan por difusión a la unión colectora (antes de recombinarse) y son arrastrados por el campo eléctrico por lo que aumenta la densidad de corriente.\\\\
En las uniones no independientes, la base pierde electrones ya que los inyecta en el emisor a través de la unión BE que está en directa y pierde electrones y huecos por recombinación, además inyecta electrones para mantener el exceso de huecos. \\\\
La base inyecta entonces los portadores mayoritarios, por lo que para anular las contribuciones de los minoritarios, hacemos que las uniones sean asimétricas ya que esto hace que sólo tengamos contribuciones de portadores mayoritarios en los terminales.Para mantener el exceso de electrones, utilizamos la \textbf{corriente de base} que es la corriente necesaria para mantener en directa la unión BE.
\subsubsection{Transistor NPN}
En este caso tendremos la misma disposición que antes pero trabajaremos con huecos y tendremos que polarizar con $V_\mathrm{CE}>0$, nuestro diagrama de bandas en este caso en lugar de ser dos pozos finitos, son dos barreras y cambiamos la posición de los cuasiniveles de Fermi de forma análoga a la figura \ref{fig: Union metalsemiconductor tipo P, curva hacia abajo}. 
\newpage
\subsection{\textcolor{red}{Regiones de polarización}}
\begin{wrapfigure}{r}{0.3\textwidth}
\centering
\includegraphics{Standalone/T4_10_Modos_de_operacion.pdf}
\caption{Modos de operación de los transistores en función de los voltajes entre terminales.}
\label{fig: Modos de operacion}
\end{wrapfigure}
Ya hemos visto que aunque tengamos una polarización global en un transistor, tenemos diferentes polarizaciones en las uniones, por ello en función de las distintas polarizaciones que tengamos entre los terminales, obersevaremos distintos fenómnos.\\\\
La figura \ref{fig: Modos de operacion}, nos muestra cómo en función de las diferentes polarizaciones de las uniones base-emisor, colector-base.  Vamos a estudiar cómo sería para un \textbf{transistor PNP} \label{sec: Reg Polarizacion trns PNP}, aunque solo vamos a centrarnos en las regiones de la base.

\subsubsection{Región Activa Directa}
\noindent Esta se da cuando tenemos polarización directa en la unión Base-Emisor y polarización inversa en la unión Colector-Base.\\
\begin{figure}[h]
\centering
    \begin{subfigure}[h]{0.3\linewidth}
    \centering
        \includegraphics{Standalone/T4_11_Region_activa_directa.pdf}
        \caption{$V_\mathrm{EB}=\text{cte}>0; V_\mathrm{CB}^\text{var}<0$}
    \end{subfigure}
    \hfill
        \begin{subfigure}[h]{0.3\linewidth}
    \centering
        \includegraphics{Standalone/T4_12_Region_activa_directa.pdf}
        \caption{$V_\mathrm{EB}^\text{(var)}>0; V_\mathrm{CB}=\text{cte}<0$}
    \end{subfigure}
    \hfill 
        \begin{subfigure}[h]{0.3\linewidth}
    \centering
        \includegraphics{Standalone/T4_13_Region_activa_directa.pdf}
        \caption{$V_\mathrm{EB}>0; V_\mathrm{CB}<0$}
    \end{subfigure}
    \caption{Concentración de portadores minoritarios de la base en la región activa directa.}
    \label{fig: Region activa directa}
\end{figure}


\begin{minipage}{0.45\textwidth}
\subsubsection{Región Activa Inversa}
  Decimos que el transistor se encuentra en la región activa inversa cuando tenemos polarización directa en la unión Colector-Base y polarización inversa en la unión Emisor-Base. $V_\mathrm{EB}=\text{cte}<0; V_\mathrm{CB}^\text{var}>0$
  \begin{center}
    \includegraphics{Standalone/T4_14_Region_activa_inverasa.pdf}
        \captionof{figure}{Concentración de portadores minoritarios de la base.}
        \label{fig: Region activa inverasa}
  \end{center}
\end{minipage}
\hfill
\begin{minipage}{0.45\textwidth}
\subsubsection{Región de Corte}
Decimos que el transistor se encuentra en la región de corte cuando tenemos polarización inversa en ambas uniones. En este caso tenemos que $V_\mathrm{EB}<0;V_\mathrm{CB}<0$ \vspace{13pt}
  \begin{center}
    \includegraphics{Standalone/T4_15_Region_de_corte.pdf}
        \captionof{figure}{Concentración de portadores minoritarios de la base.}
        \label{fig: Region de corte}
  \end{center}
\end{minipage}

\subsubsection{Región de Saturación}
Decimos que el transistor está en saturación cuando ambas uniones están en directa.

\begin{figure}[h]
    \centering
    \begin{subfigure}[h]{0.4\linewidth}
    \centering
        \includegraphics{Standalone/T4_16_Region_saturacion.pdf}
        \caption{$V_\mathrm{EB}^\text{(var)}>0; V_\mathrm{CB}^\text{(var)}>0;V_\mathrm{EB}>V_\mathrm{CB}$}
    \end{subfigure}
    \hspace{2cm}
        \begin{subfigure}[h]{0.4\linewidth}
    \centering
        \includegraphics{Standalone/T4_17_Region_saturacion.pdf}
        \caption{$V_\mathrm{EB}^\text{(var)}>0; V_\mathrm{CB}^\text{(var)}>0;V_\mathrm{EB}<V_\mathrm{CB}$}
    \end{subfigure}
    \caption{Concentración de portadores minoritarios de la base en la región saturación.}
    \label{fig: Region saturacion}
\end{figure}
Para un \textbf{transistor NPN} \label{sec: Regiones de polarizacion trns NPN}, observamos que la concentración de portadores minoritarios en la base sigue las mismas curvas que para el transistor PNP.
\section{Comportamiento estático: características corriente-tensión. Modelo de Ebers-Moll}
A la hora de estudiar la corriente, utilizaremos el siguiente sistema de coordenadas y criterio de signos (desarrollamos para un PNP).
\begin{figure}[h!]
    \centering
    \begin{subfigure}[h]{0.4\textwidth}
    \includegraphics{Standalone/T4_18_Sistema_de_coordenadas_en_un.pdf}
    \caption{Sistema de coordenadas.}
    \label{fig:Sistema de coordenadas en un transistor}        
    \end{subfigure}
    \hspace{2cm}
\begin{subfigure}[h]{0.4\textwidth}
    \centering
    \includegraphics{Standalone/T4_19_Transistor_PNP.pdf}
            \caption{Criterio de signos.}
            \label{fig: Transistor PNP}    
            \end{subfigure}
            \caption{Geometría del modelo de Ebers-Moll.}
            \label{fig: Modelo de Ebers-Moll}
\end{figure}

El \textbf{modelo de Ebers-Moll} tiene las siguientes características:
\begin{itemize}
    \item Es un modelo \textbf{unidimensional}.
    \item Nos encontramos en \textbf{bajo nivel de inyección}, por lo que las relaciones de Boltzmann son válidas.
    \item Tenemos \textbf{ausencia de generación-recombinación} en las ZCE's.
    \item  Consideraremos que las ZCE's \textbf{no varían con la polarización}.
    \item  Las corrientes serán \textbf{de difusión} de los minoritarios en los bordes de ZCE y serán \textbf{constantes} dentro de la ZCE. La suma de las corrientes es nula $\sum_i I_i=0$.
\end{itemize}
\newpage
\subsection{Distribuciones de minoritarios en las zonas neutras}
Supongamos que el emsior y el colector son largos, es decir que sus longitudes de difusión son mucho menores que la longitud física de los terminales. Por lo que las concentraciones de minoritarios son 
\begin{equation}
    n_\mathrm{E}(x^\prime)=n_\mathrm{E0}+n_\mathrm{E0}\left( \exp{\left( \frac{eV_\mathrm{EB}}{k_\mathrm{B}T}\right)-1}\right) \exp{\left( \frac{-x^\prime}{L_{n\mathrm{E}}} \right)} \ \ \text{Emisor}
    \label{eq: Minoritarios emisor}
\end{equation}
\begin{equation}
    n_\mathrm{C}(x^{\prime \prime})=n_\mathrm{C0}+n_\mathrm{C0}\left( \exp{\left( \frac{eV_\mathrm{CB}}{k_\mathrm{B}T}\right)-1}\right) \exp{\left( \frac{-x^\prime}{L_{n\mathrm{C}}} \right)} \ \ \text{Colector}
    \label{eq: minoritarios colector}
\end{equation}
Supongamos también que la base es corta, por lo que de la ecuación de continuidad \ref{eq: ecuacion ambipolar de difusion}, solo tendremos en cuenta la contribución de difusión.
\begin{equation*}
    D_{p \mathrm{B}} \diff[2]{ p_\mathrm{B}(x)}{x}=0 \quad p_\mathrm{B}(x)=A x+B
    \left\{\begin{aligned}
&x=0 \rightarrow p_\mathrm{B}(0)=p_{\mathrm{B} 0} \exp{\frac{e V_\mathrm{EB}}{k_\mathrm{B} T}} \\\\
&x=W \rightarrow p_\mathrm{B}(W)=p_{\mathrm{B} 0} \exp{\frac{e V_\mathrm{CB}}{k_\mathrm{B} T}}
\end{aligned}\right.
\end{equation*}
Resultando
\begin{equation}
    p_\mathrm{B}(x)=p_\mathrm{B0}\, \frac{W-x}{W}\, \exp{\left(\frac{e V_\mathrm{EB}}{k_\mathrm{B} T}\right)} +p_\mathrm{B0}\, \frac{x}{W}\, \exp{\left(\frac{e V_\mathrm{CB}}{k_\mathrm{B} T}\right)}
\end{equation}
\subsection{\textcolor{red}{Corrientes en los terminales. Ecuaciones de Ebers-Moll}}
Si consideramos que las corrientes son positivas al entrar en el transistor, la corriente de emisor será

\begin{equation}
\begin{aligned}
    I_\mathrm{E}=&e S\left(-\left.D_{p \mathrm{B}} \frac{\mathrm{d} p_\mathrm{B}(x)}{\mathrm{d} x}\right|_{x=0}-\left.D_{n \mathrm{E}} \frac{\mathrm{d} n_\mathrm{E}\left(x^{\prime}\right)}{\mathrm{d} x^{\prime}}\right|_{x^{\prime}=0}\right)=\\&e S\left[\left(D_{p \mathrm{B}} \frac{p_{\mathrm{B} 0}}{W}+D_{n \mathrm{E}} \frac{n_{\mathrm{E} 0}}{L_{n \mathrm{E}}}\right)\left(\exp{\left(\frac{e V_\mathrm{EB}}{k_\mathrm{B}T}\right)}-1\right)-\frac{D_{p \mathrm{B}}}{W} p_{\mathrm{B} 0}\left(\exp{\left(\frac{e V_\mathrm{CB}}{k_\mathrm{B}T}\right)}-1\right)\right] 
    \end{aligned}
\end{equation}

Mientras que para el colector la corriente es:

\begin{equation}
\begin{aligned}
 I_C=&e S\left(+\left.D_{p \mathrm{B}} \frac{\mathrm{d} p_\mathrm{B}(x)}{\mathrm{d} x}\right|_{x=W}-\left.D_{n \mathrm{E}} \frac{\mathrm{d} n_\mathrm{C}\left(x^{\prime \prime}\right)}{\mathrm{d} x^{\prime \prime }} \right|_{x^{\prime \prime}=0}\right)=\\&e S\left[\left(D_{p \mathrm{B}} \frac{p_{\mathrm{B} 0}}{W}+D_{n \mathrm{C}} \frac{n_{\mathrm{C} 0}}{L_{n \mathrm{C}}}\right)\left(\exp{\left(\frac{e V_\mathrm{CB}}{k_\mathrm{B}T}\right)}-1\right)-\frac{D_{p \mathrm{B}}}{W} p_{\mathrm{B} 0}\left(\exp{\left(\frac{e V_\mathrm{EB}}{k_\mathrm{B}T}\right)}-1\right)\right] 
 \end{aligned}
\end{equation}

Siendo $S$ la sección tranversal del terminal. Si utilizamos la ecuación de malla, obtenemos la expresión para $I_\mathrm{B}=-I_\mathrm{C}-I_\mathrm{E}$
\newpage
\noindent Si definimos unas corrientes conocidas como\\
\begin{minipage}{0.45 \textwidth}
    \begin{equation}
        I_\mathrm{CS}=eS\left(D_{p \mathrm{B}} \frac{p_{\mathrm{B} 0}}{W}+D_{n \mathrm{E}} \frac{n_{\mathrm{E} 0}}{L_{n \mathrm{E}}}\right)
    \end{equation}
    Corriente inversa de saturación de colector, con la unión emisora en cortocircuito.
\end{minipage}
\hfill
\begin{minipage}{0.45 \textwidth}
    \begin{equation}
        I_\mathrm{ES}=eS\left(D_{p \mathrm{B}} \frac{p_{\mathrm{B} 0}}{W}+D_{n \mathrm{E}} \frac{n_{\mathrm{E} 0}}{L_{n \mathrm{E}}}\right)
    \end{equation}
    Corriente inversa de saturación de emisor, con la unión colectora en cortocircuito.
\end{minipage}
\\\\\\
Obtenemos las \textbf{ecuaciones de Ebers-Moll}
\begin{equation}
    I_\mathrm{E}=I_\mathrm{ES} \left(\exp{\left(\frac{eV_\mathrm{EB}}{k_\mathrm{B}T} \right)}-1 \right)-\alpha_\mathrm{R} I_\mathrm{CS}\left( \exp \frac{eV_\mathrm{CB}}{k_\mathrm{B}T}-1 \right)
    \label{eq:Corriente de Emisor EbersMoll}
\end{equation}
\begin{equation}
    I_\mathrm{C}=-\alpha_\mathrm{F}I_\mathrm{ES} \left(\exp{\left(\frac{eV_\mathrm{EB}}{k_\mathrm{B}T} \right)}-1 \right)+ I_\mathrm{CS}\left(\exp \frac{eV_\mathrm{CB}}{k_\mathrm{B}T}-1 \right)
    \label{eq:Corriente de Colector EbersMoll}
\end{equation}
Siendo los parámetros $\alpha_i$ ganancias en corriente definidas por \\
\begin{minipage}{0.5\textwidth}
    \begin{equation}
\alpha_\mathrm{F}=\cfrac{D_{p \mathrm{B}} \cfrac{p_{\mathrm{B} 0}}{W}}{D_{p \mathrm{B}} \cfrac{p_{\mathrm{B} 0}}{W}+D_{p \mathrm{B}} \cfrac{p_{\mathrm{E} 0}}{L_{n\mathrm{E}}}}=-\left. \cfrac{I_\mathrm{C}}{I_\mathrm{E}}\right |_{V_\mathrm{CB}=0}
\label{eq: Ganancia directa en corriente}
    \end{equation}
    \textbf{Ganancia directa en corriente} \\ con la unión colectora en cortocircuito.
\end{minipage}
\hfill
\begin{minipage}{0.5\textwidth}
    \begin{equation}
\alpha_\mathrm{R}=\cfrac{D_{p \mathrm{B}} \cfrac{p_{\mathrm{B} 0}}{W}}{D_{p \mathrm{B}} \cfrac{p_{\mathrm{B} 0}}{W}+D_{p \mathrm{B}} \cfrac{p_{\mathrm{C} 0}}{L_{n\mathrm{C}}}}=-\left. \cfrac{I_\mathrm{E}}{I_\mathrm{C}}\right |_{V_\mathrm{EB}=0}
\label{eq: Ganancia inversa en corriente}
    \end{equation}
    \textbf{Ganancia inversa en corriente} \\ con la unión emisora en cortocircuito.
\end{minipage}
\\\\\\
El parámetro $\alpha_\mathrm{F}$ nos interesa que sea muy próximo a uno, es decir, este parámetro nos indica qué tan bueno es nuestro transistor. Para ello, si analizamos la ecuación \ref{eq: Ganancia directa en corriente} podemos reducir mucho $W$, lo cual hace que los diferentes efectos no ideales cobren mayor relevancia; o podemos \textbf{impurificar mucho el emisor},lo que justifica lo expuesto en la sección \ref{sec: Transistor bipolar}
\subsection{Puntos de operación. Características de entrada y salida de un transistor bipolar}
Las ecuaciones de Ebers-Moll \ref{eq:Corriente de Emisor EbersMoll} y \ref{eq:Corriente de Colector EbersMoll}, nos relacionan dos tensiones con dos corrientes. Además por las ecuaci ones de Kirchoff, si conocemos también dos de las corrientes y dos de las tensiones, conoceremos las restantes.\\\\
Con dos de las magnitudes involucradas (corrientes ó tensiones), \textbf{fijamos el punto de operación} del transistor.\\\\
El problema del transistor se reduce entonces a trabajar con dos gráficas, que relacionan tres variables (dependiente, independiente y una que actúa como parámetro).
\subsubsection{Referencia de tensión}
Tomamos como referencia el emisor $(V_\mathrm{CE},V_\mathrm{BE})$, por lo que nuestras tensiones son \textbf{características en emisor común}.
\subsubsection{Caracterísiticas de entrada}
En la entrada tomamos como variable independiente $V_\mathrm{BE}$ y $V_\mathrm{CE}$ como parámetro, por lo que $I_\mathrm{B}=I_\mathrm{B}(V_\mathrm{BE},V_\mathrm{CE})$ es la variable dependiente.
A partir de las ecuaciones de \hyperref[eq:Corriente de Emisor EbersMoll]{Ebers-Moll}
\begin{equation}
I_\mathrm{B}=-I_\mathrm{ES} (1-\alpha_\mathrm{F})\left(\exp{\left(\frac{-eV_\mathrm{BE}}{k_\mathrm{B}T} \right)}-1 \right)-I_\mathrm{CS}(1-\alpha_\mathrm{R})\left(\exp{\left(e\frac{V_\mathrm{CE}-V_\mathrm{BE}}{k_\mathrm{B}T}\right )}-1 \right)
\label{eq: Corriente de base caracteristica de entrada}
\end{equation}
\subsubsection{Características de salida}
En la salida tomamos como variable independiente $V_\mathrm{CE}$ y $I_\mathrm{B}$ como parámetro, por lo que $I_\mathrm{C}=I_\mathrm{C}(V_\mathrm{CE},I_\mathrm{B})$ es la variable dependiente.
A partir de las ecuaciones de \hyperref[eq:Corriente de Emisor EbersMoll]{Ebers-Moll}
\begin{equation}
I_\mathrm{C}=I_\mathrm{B} \frac{\alpha_\mathrm{F}}{1-\alpha_\mathrm{F}}+\frac{I_\mathrm{CO}}{1-\alpha_\mathrm{F}}\left(\exp{\left(e\frac{V_\mathrm{CE}-V_\mathrm{BE}}{k_\mathrm{B}T}\right )}-1 \right) \quad \text{con} \quad I_\mathrm{CO}=I_\mathrm{CS}(1-\alpha_\mathrm{F}\alpha_\mathrm{R})
\label{eq: Corriente de base caracteristica de entrada}
\end{equation}
\subsection{Regiones de polarización. Corrientes en los terminales}
\subsubsection{Región Activa Directa}
En la región activa directa, tenemos la unión emisora en directa $(V_\mathrm{BE}<0)$ mientras que la colectora está en inversa $(V_\mathrm{CB}<0)$, por la relación $V_\mathrm{CE}=V_\mathrm{CB}+V_\mathrm{BE}$, concluimos que $|V_\mathrm{CE}|>|V_\mathrm{BE}|$\\\\
\begin{minipage}{0.65\textwidth}
\begin{equation}
        I_\mathrm{B}=-I_\mathrm{ES} (1-\alpha_\mathrm{F})\left(\exp{\left(\frac{-eV_\mathrm{BE}}{k_\mathrm{B}T} \right)}-1 \right)+I_\mathrm{CS}(1-\alpha_\mathrm{R})
        \label{eq: Corriente de base en directa}
\end{equation}    
Que es la expresión de la corriente de base y vemos que siempre es negativa. La región activa directa en la base presenta un efecto no ideal y es la tensión umbral en la unión base-emisor.
\end{minipage}
\hfill
\begin{minipage}{0.3\textwidth}
\begin{center}
    \includegraphics{Standalone/T4_20_Caracteristica_I_V_ideal_de.pdf}
    \captionof{figure}{Característica I-V ideal de la base.}
\end{center}
\end{minipage}

\begin{equation}
            I_\mathrm{C}=I_\mathrm{B} \frac{\alpha_\mathrm{F}}{1-\alpha_\mathrm{F}}+\frac{I_\mathrm{CO}}{1-\alpha_\mathrm{F}}=\beta_\mathrm{F}I_\mathrm{B}-(\beta_\mathrm{F}+1)I_\mathrm{CO}\approx \beta_\mathrm{F}I_\mathrm{B} \quad \text{con} \quad \beta_\mathrm{F}=\left. \frac{I_\mathrm{C}}{I_\mathrm{B}}\right|_{V_\mathrm{CB}=0}=\frac{\alpha_\mathrm{F}}{1-\alpha_\mathrm{F}}
            \label{eq: Corriente de colector en directa}
\end{equation}
\begin{wrapfigure}{r}{0.3\textwidth}
\flushright
    \includegraphics{Standalone/T4_21_Caracterisitica_real_I_V_del.pdf}
    \caption{Caracterísitica real I-V del colector.}
\end{wrapfigure}

Siendo $\beta_\mathrm{F}$ la \textbf{ganancia directa en corriente}. La corriente de colector también es negativa en la región activa directa. Los efectos reales que presenta la corriente de colector son
\begin{itemize}
    \item \textcolor{red}{\textbf{Efecto Early}}, en el que observamos que las pendientes son no nulas en las corrientes de base, por lo que la suposición de que las ZCE son constantes no se puede aplicar a un transistor real.
    \item \textbf{Alta inyección}, conforme aumentan el valor de las corrientes, se van acumulando en un valor de $I_\mathrm{C}$.
\end{itemize}
\clearpage
\subsubsection{Región Activa Inversa}
En la región activa inversa, la unión emisora está en inversa ($V_\mathrm{BE}>0$), mientras que la colectora está en directa $(V_\mathrm{CB}>0)$. La caida de potencial $V_\mathrm{CE}>0$ por lo tanto obtenemos $|V_\mathrm{CE}|>|V_\mathrm{BE}|$ \\
\begin{minipage}{0.7\textwidth}
    La expresión para la corriente de base es
    \begin{equation}
        I_\mathrm{B}=I_\mathrm{ES} (1-\alpha_\mathrm{F})-I_\mathrm{CS}(1-\alpha_\mathrm{R})\left(\exp{\left(e\frac{V_\mathrm{CE}-V_\mathrm{BE}}{k_\mathrm{B}T}\right )}-1 \right)
        \label{eq: Corriente de base en inversa}
    \end{equation}
    En la región activa inversa, la corriente de base es negativa.
\end{minipage}
\hfill
\begin{minipage}{0.3\textwidth}
\begin{center}
    \includegraphics{Standalone/T4_22_Caracteristica_I_V_ideal_en.pdf}
    \captionof{figure}{Característica I-V ideal en la base.}
\end{center}
\end{minipage}

\noindent La corriente de colector la podemos expresar como 
\begin{equation}
    I_\mathrm{C}=I_\mathrm{B} \frac{\alpha_\mathrm{R}}{1-\alpha_\mathrm{R}}+\frac{I_\mathrm{CO}}{1-\alpha_\mathrm{R}} \left(\exp{\left(\frac{-eV_\mathrm{BE}}{k_\mathrm{B}T}\right )}-1 \right)=\beta_\mathrm{R}I_\mathrm{B}-(\beta_\mathrm{R}+1)I_\mathrm{EO}
    \label{eq: Corriente de colector en inversa}
\end{equation}
Siendo $I_\mathrm{EO}=I_\mathrm{ES}(1-\alpha_\mathrm{F}\alpha_\mathrm{R})$ y $\beta_\mathrm{R}$ la ganancia inversa en corriente cuya expresión es similar a la dada en la ecuación \ref{eq: Corriente de base en directa}. 
\subsubsection{Región de Saturación}
Ambas uniones están en directa $(V_\mathrm{BE}<0, V_\mathrm{CB}>0)$ y $V_\mathrm{CE}\approx 0$. $I_\mathrm{C}$ la determina el circuito en que esté el transistor $|\beta_\mathrm{F}I_\mathrm{B}>I_\mathrm{C\,max}| \xrightarrow{} I_\mathrm{C}=I_\mathrm{C\,max}=I_\mathrm{C}(V_\mathrm{CE}=0)$

\begin{wrapfigure}{r}{0.3\textwidth}
\flushright
    \includegraphics{Standalone/T4_23_Caracterisitica_real_I_V_del.pdf}
    \caption{Caracterísitica real I-V del colector.}
\end{wrapfigure}
\subsubsection{Región de Corte}

Ambas uniones están en inversa $(V_\mathrm{BE}>0, V_\mathrm{CB}<0)$ y $I_\mathrm{C}\approx0$. $V_\mathrm{CE}$ la determina el circuito en el que está el transistor.
Las expresiones para las corrientes de emisor, colector y base toman expresiones más sencillas que las descritas anteriormente siendo la corriente de emisor
\begin{equation}
    I_\mathrm{E}=-I_\mathrm{ES}(1-\alpha_\mathrm{F})
\end{equation}
La corriente de colector
\begin{equation}
    I_\mathrm{C}=-I_\mathrm{CS}(1-\alpha_\mathrm{R})
\end{equation}
Y para la corriente de base
\begin{equation}
    I_\mathrm{B}=-I_\mathrm{ES}(1-\alpha_\mathrm{F})+I_\mathrm{CS}(1-\alpha_\mathrm{R})
\end{equation}
\clearpage

\section{Comportamiento dinámico: señales débiles de baja frecuencia}
Supongamos que las señales de baja frecuencia tienen la siguiente forma, $$X_i(t)=X_{i0}+x_i(t)$$ donde las $X_{i0}$ representan los puntos de operación estacionarios y $x_i(t)$, representan las contribuciones dependientes del tiempo; la $i=\{\text{E, B, C} \}$, representa al terminal. Por lo que tenemos las siguientes expresiones para la intensidad y el voltaje de los terminales
\begin{equation}
\begin{aligned}
&I_\mathrm{E}(t)=I_{\mathrm{E} 0}+i_\mathrm{E}(t) ; \quad I_\mathrm{C}(t)=I_{\mathrm{C} 0}+i_\mathrm{C}(t) ; \quad I_\mathrm{B}(t)=I_{\mathrm{B} 0}+i_\mathrm{B}(t) \\\\
&V_{\mathrm{C} \mathrm{E}}(t)=V_{\mathrm{C} \mathrm{E} 0}+v_{\mathrm{C} \mathrm{E}}(t) ; \quad V_{\mathrm{C} \mathrm{B}}(t)=V_{\mathrm{C} \mathrm{B} 0}+v_{\mathrm{C} \mathrm{B}}(t) ; \quad V_{\mathrm{B} \mathrm{E}}(t)=V_{\mathrm{B} \mathrm{E} 0}+v_{\mathrm{B} \mathrm{E}}(t) \\\\
&i_\mathrm{E}+i_\mathrm{C}+i_\mathrm{B}=0 \quad v_{\mathrm{C} \mathrm{B}}=v_{\mathrm{C} \mathrm{E}}-v_{\mathrm{B} \mathrm{E}}
\end{aligned}
\label{eq: Ecuaciones de los terminales en baja señal}
\end{equation}
El modelo de pequeña señal nos permite establecer dos ecuaciones lineales ($f_1(i_1,i_2,v_1,v_2)=0$) y ($f_2(i_1,i_2,v_1,v_2)=0$) que relacionen entre las cuatro variables.\\\\
En este caso tenemos un \textbf{cuadripolo} ya que nuestro dispositivo esta siendo sometido a dos intensidades y a dos tensiones.\\\\
Tenemos seis combinaciones según qué variables escojamos como independientes y dependientes, ya que si nos fijamos en las relaciones lineales y en las últimas de las ecuaciones \ref{eq: Ecuaciones de los terminales en baja señal}, sólo nos hacen falta escoger dos de los 3 terminales, tomando el terminal que falta como referencia. Por lo que tenemos 18 posibilidades en total para determinar  cuatro coeficientes de las relaciones $f_i$. Es decir que sólo hace falta determinar 4 coeficientes independientes.\\\\
Vamos a determinar las relaciones $f_i$ considerando lo que se conoce como \textbf{parámetros híbridos}, si consideramos como variables dependientes $v_1$ y $i_2$; y $v_2$ y $i_1$ como independientes.

El sistema de ecuaciones que obtenemos es entonces
\begin{equation}
    \left\{ \begin{aligned}
        & v_1=h_{11}i_1+h_{12}v_2\\
        & i_2=h_{21}i_1+h_{22}v_2\\
    \end{aligned}\right.
\end{equation}
Donde el sentido físico de los parámetros híbridos es\footnote{Todos estos parámetros son para las componentes alternas. Recordar también que independientemente de los terminales que escojamos para resover tenemos que tener un terminal de referencia común.}:
\begin{itemize}
    \item $h_{11}$ \textbf{impedancia de entrada}. con la salida en cortocircuito.
    \item  $h_{12}$ \textbf{ganancia inversia en tensión}, con la entrada en circuito abierto.
    \item $h_{21}$ \textbf{ganancia directa en corriente}, con la salida en cortocircuito.
    \item $h_{22}$ \textbf{admitancia de salida}, con la entrada en circuito abierto.
\end{itemize}
\newpage

\begin{figure}[!ht]
\centering
\includegraphics[width=0.5\textwidth]{Standalone/T4_24_Representacion_circuital_transistor_BJT.pdf}
\caption{Representación circuital de un transistor}
\label{fig: Representacion circuital transistor BJT}
\end{figure}
Los parámetros dependen del punto de operación y de la frecuencia y por lo general serán números complejos, pero al trabajar en baja frecuencia, son reales e independientes de la misma.\\\\ La figura \ref{fig: Representacion circuital transistor BJT} es una representación equivalente de un transistor. En baja frecuencia y en la región activa, podemos sustituir las impedancia por resistencias, ya que el comportamiento del circuito\footnote{Podríamos escoger 1=B y 2=C} será lineal. Para el régimen de alta frecuencia, debemos considerar las contribuciones capacitivas de las uniones PN.
\subsection{Conmutación del BJT}
Cuando cambiamos bruscamente la polaridad del circuito de un transistor bipolar, no esperamos que los cambios sean inmediatos por lo que definimos una serie de tiempos que nos especifican hasta qué valor de frecuencia podemos utilizar el transistor.
\begin{itemize}
    \item $t_\mathrm{d}$ \textbf{tiempo de retardo}, es el tiempo que tarda la unión BE en pasar de inversa a borde de directa.
    \item $t_\mathrm{r}$ \textbf{tiempo de subida}, es el tiempo que tarda la unión BE de estar en el borde de directa al borde de saturación.
    \item $t_\mathrm{s}$ \textbf{tiempo de almacenamiento}, este es el que nos limita la velocidad de conmutación del transistor; y es el tiempo necesario para extraer la carga en exceso almacenada en la base cuando nos encontramos en saturación.
    \item $t_\mathrm{f}$ \textbf{tiempo de caída}, es el tiempo que tarda la unión BE en estar en el borde de saturación a estar en el bode de inversa, es decir de extraer la carga total.
\end{itemize}

\pagestyle{fancy}
\fancyhead[R]{Electrónica Física: Tema 5}
\fancyhead[L]{4º Grado en Física}

\chapter{Transistor de Efecto Campo Metal-Óxido-Semiconductor (MOSFET)}
Una vez desarrollado el transistor BJT, vamos a estudiar el transistor MOSFET. Este tipo de transistores, se basan en la estructura Metal-Aislante-Semiconductor (MIS), donde el aislante puede ser un óxido (entonces la denominamos MOS)\footnote{Los modelos que vamos a desarrollar son para dispositivos basados en MOS, aunque son igualmente aplicables para la estructura MIS}.\\\\ 
El MOSFET es un dispositivo de cuatro terminales \textbf{puerta} (G), \textbf{drenador} (D), \textbf{fuente} (S) y \textbf{sustrato} (B).\\

La idea fundamental detrás del MOSFET es controlar la corriente entre fuente y drenador a través de la tensión aplicada a la puerta. Destacar que el sustrato \textbf{siempre} tiene que ser del tipo contrario a la fuente y al drenador.
\begin{figure}[h]
\begin{subfigure}{0.45\linewidth}
\centering
       \includegraphics{Standalone/T5_01_Transistor_Mosfet.pdf}
    \caption{Esquema de un transistor MOSFET.}
    \label{fig:Transistor Mosfet} 
\end{subfigure}
\hspace{2cm}
\begin{subfigure}{0.45\linewidth}
\centering
    \includegraphics{Standalone/T5_02_Condensador_Mos.pdf}
    \caption{Esquema de un condensador MOS.}
    \label{fig:Condensador Mos}
\end{subfigure}
\caption{Representaciones de un MOSFET y un MOS}
\end{figure}
\newpage
\section{Estructura metal-óxido-semiconductor (MOS)}
En este caso, vamos a eliminar los terminales de fuente y drenador, por lo que tendremos un condensador representado en la figura \ref{fig:Condensador Mos}. Las zonas grises son las zonas donde se encuentran distribuidas las cargas.

\subsection{Situaciones de carga en el semiconductor. Análisis cualitativo de las bandas de energía}
Supongmaos que el semiconductor es tipo P. Cuando establecemos una diferencia de potencial entre los extremos del condensador MOS, observamos que las bandas de energía se comportan de manera análoga a lo discutido en la sección \ref{sec: Union MSC-P}.

\subsubsection{Acumulación}
\noindent Si tenemos $V<0$, como se muestra en la figura \ref{fig: Acumulacion MOS}, vemos que el campo eléctrico que nos genera la diferencia de potencial entre extremos del condensador, hace que dentro del semiconductor los huecos se desplacen hacia el óxido. Generando así, dentro del óxido, un campo eléctrico debido a la diferencia de potencial que se crea entre sus extremos. A este estado de la estructura MOS se le conoce como \textbf{acumulación}.
\subsubsection{Vaciamiento}
Si tenemos $V>0$ pero siendo pequeña, el campo eléctrico va en sentido contrario al anterior, por lo que tendremos un exceso de carga positiva en el contacto con el metal del óxido y los huecos que se encontraban en la parte superior, se nos desplazarán hacia el contacto inferior del semiconductor, por lo que se nos \textbf{inducirá una zona de carga espacial}, y alrededor de esta se nos formará una región cargada negativamente. A esta estado de la estructura MOS lo conocemos como \textbf{vaciamiento}.
\begin{figure}[h]
\centering
\begin{subfigure}{0.35\linewidth}
\centering
    \includegraphics{Standalone/T5_03_Acumulacion_MOS.pdf}
     \captionsetup{skip=1pt} % Ajusta el espacio vertical aquí
    \caption{Acumulación de cargas en un MOS.}
    \label{fig: Acumulacion MOS}
    \end{subfigure}
    \hspace{1mm}
    \begin{subfigure}{0.55\linewidth}
    \centering
    \includegraphics{Standalone/T5_04_Acumulacion_MOS.pdf}
    \caption{Vaciamiento de cargas. En la segunda imágen vemos la ZCE inducida.}
    \label{fig: Acumulacion MOS}
    \end{subfigure}
    \caption{Acumulación y vaciamiento}
\end{figure}
\newpage
\subsubsection{Inversión}
\begin{minipage}{0.6 \textwidth}
 \noindent Cuando tenemos un $V>0$ lo suficientemente grande, aumentará la zona de carga inducida, esto hace que las bandas de energía se doblen más hacia abajo, tanto que el cuasinivel de Fermi se encuentra por debajo del nivel de Fermi. Esto implica que la banda de conducción se encuentre más cercano al nivel de Fermi que la banda de conducción por lo que la zona cercana a la superficie del semiconductor se convierte en tipo N, es decir, hemos invertido el tipo de semiconductor. A este estado de la estructura MOS se le denomina \textbf{inversión}.\\\\
Para la estructura MOS con un semiconductor tipo N las bandas se tornan como lo desarrollado en la sección \ref{sec: Union MSC}
\end{minipage}
\hfill
\begin{minipage}{0.3 \textwidth}
    \begin{figure}[H]
\centering
    \includegraphics{Standalone/T5_05_Inversion_MOS.pdf}
    \caption{Inversión de cargas.}
    \label{fig: Inversion MOS}
\end{figure}
\end{minipage}
\subsection{Situacioness de carga en el semiconductor. Análisis cuantitativo de las bandas de energía}
Si en la figura \ref{fig: Union MSC sin contacto}, entre el metal y es semiconductor introducimos un aislante (óxido), tenemos que la afinidad eletrónica del aislante está más próxima al nivel de vacío que en el semiconductor. Además destacamos que en el óxido no tiene sentido establecer un nivel de Fermi.

\begin{figure}[h!]
    \centering
\includegraphics{Standalone/T5_06_Diagrama_de_bandas_de_los.pdf}
    \caption{Diagrama de bandas de los diferentes materiales aislados.}
    \label{fig:enter-label}
\end{figure}
Vamos a definir el potencial $e\varphi_\mathrm{F}=E_\mathrm{Fi}-E_\mathrm{F}$ que en el caso de semiconductores tipo P es positivo y en semiconductores tipo N es negativo. A partir de ahora supongamos que $\varphi_\mathrm{MS}<0$. Vamos a desarrollar los estados cuando la estructura MOS está en equilibrio y cuando no lo está.
\subsubsection{Equilibrio}
En este estado el nivel de Fermi de la unión se es constante y la corriente es nula (debido al aislante). Si nos refermos a los anchos de las zonas de carga espacial, la del metal es mucho menor que la del semiconductor, esto se debe a que en el metal la carga se distribuye por la superficie mientras que en el semiconductor se distribuye en un cierto volumen. Definimos los siguientes potenciales:
\begin{itemize}
    \item $V_\mathrm{OO^\prime}$: Caída de tensión en el óxido en equilibrio ($V_\mathrm{OO^\prime}>0$ en la figura).
    \item $\varphi_\mathrm{SO}$: caída de tensión en el semiconductor en equilibrio ($\varphi_\mathrm{SO}>0$ en la figura).
    \item  $\varphi(x)$: potencial en una posición genérica $x$ del semiconductor (no lo representamos por simplicidad).
\end{itemize}
Si nos damos cuenta, $\phi_\mathrm{MS}=-e(V_\mathrm{OO^\prime}+\varphi_\mathrm{SO})<0$. También vamos a recordar que los potenciales están medidos desde el nivel de vacío (nivel del semiconductor en equilibrio), por lo que cuanto más abajo se encuentre una energía, mayor será el potencial que la genera. El metal se encuentra a mayor potencial (y menor energía) que el semiconductor en $-\phi_\mathrm{MS}>0$.
%AQUÍ HAY QUE PONER LA FIGURA DE LA HOJA 33 DE LA PRESENTACIÓN
Dentro del aislante, el campo es constante (es un condensador), por lo que el potencial es lineal através de él.
\begin{figure}[h]
    \centering
    \includegraphics{Standalone/T5_07_Estructura_MOS_en_equilibrio.pdf}
    \caption{Estructura metal (izq) óxido (central) semiconductor (dcha), en equilibrio siendo un semiconductor de tipo P.}
    \label{fig:Estructura MOS en equilibrio}
\end{figure}
\subsubsection{Fuera del equilibrio}
\begin{itemize}
    \item \textbf{Banda plana}\\
    \indent \indent Definimos el potencial de banda plana $V_\mathrm{FB}=\phi_\mathrm{MS}/e$ como aquel potencial que compensa exactamente la diferencia de potencial de equilibrio. En este caso $I=0$ y $E_\mathrm{F}$ es constante y el potencial que antes era lineal en el óxido ahora también es \textbf{constante}. Esta estructura no presenta acumulaciones de cargas ni caídas de tensión.
    \begin{figure}[h!]
    \centering
    \includegraphics{Standalone/T5_08_Potrencial_de_banda_blana.pdf}
    \caption{Diagrama de bandas de una estructura MOS cuando el potencial de rejilla es el de banda plana.} 
    \label{fig: Potrencial de banda blana}
\end{figure}
    \item \textbf{Acumulación}\\
    \indent \indent En este caso el potencial aplicado $V_\mathrm{G}<\phi_\mathrm{MS}/e<0$. Cuando nos encontramos en acumulación, tenemos dos caídas de tensión una el óxido (el potencial es lineal) $V_\mathrm{OX}<0$ y una en el semiconductor $\varphi_\mathrm{S}$ por lo que el metal se encuentra a menor potencial que el semiconductor, resultando \begin{equation}
        V_\mathrm{G}-\frac{\phi_\mathrm{MS}}{e}=V_\mathrm{OX}+\varphi_\mathrm{S}<0
    \end{equation}
    \begin{figure}
        \centering
        \includegraphics{Standalone/T5_09_Estructura_MOS_en_acumulacion.pdf}
        \caption{Diagrama de bandas de una estructura MOS cuando se encuentra en acumulación.}
        \label{fig:Estructura MOS en acumulacion}
    \end{figure}
    \footnote{No incluimos el potencial $\phi_\mathrm{SO}$ por simpificar los esquemas} Como hemos dicho anteriormente, tendremos una acumulación de huecos en la frontera entre el óxido y el semiconductor.
    \item \textbf{Inversión}\\
    \indent \indent Cuando aplicamos un $V_\mathrm{G}>0$, lo suficientemente grande nos encontramos en inversión (será la condición necesaria para que exista corriente entre fuente y drenador). En este caso $V_\mathrm{OX}>0$, $\varphi_\mathrm{S}>2\varphi_\mathrm{F}>0$ y el metal se encuentra a mayor potencial que el semiconductor
    \begin{equation}
        V_\mathrm{G}-\frac{\phi_\mathrm{MS}}{e}=V_\mathrm{OX}+\varphi_\mathrm{S}>0
    \end{equation}
    Como dijimos antes, hay un exceso de electrones, tantos que superan al número de huecos en equilibrio \footnote{Si fuera un SCN sería al revés}. Si vamos aumentando $V_\mathrm{G}$ desde 0, tendremos vaciamiento hasta que \\ $V_\mathrm{G}>2\varphi_\mathrm{F}$ en ese caso tendremos inversión y ambas capas coexistirán en el semiconductor.
    \begin{itemize}
    \newpage
    \begin{figure}[h]
    \centering
    \includegraphics{Standalone/T5_10_Estructura_MOS_en_inversion.pdf}
    \caption{Diagrama de bandas de una estructura MOS en inversión.}
    \label{fig:Estructura MOS en inversion}
\end{figure}
        \item \textbf{Condición de inversión} Ya hemos dicho que necesitamos que $V_\mathrm{G}>\varphi_\mathrm{F}$. Si calculamos las concentraciones de portadores de carga obtenemos
        \begin{equation}
            \begin{aligned}
                &n(x)=N_\mathrm{C} \exp{\left( \frac{E_\mathrm{F}-E_\mathrm{C}+e\varphi(x)}{k_\mathrm{B}T} \right)}=n_\mathrm{i}\exp{\left( \frac{-e(\varphi_\mathrm{F}-\varphi(x))}{k_\mathrm{B}T} \right)}\\
                &p(x)=N_\mathrm{V} \exp{\left( \frac{E_\mathrm{V}-E_\mathrm{F}+e\varphi(x)}{k_\mathrm{B}T} \right)}=n_\mathrm{i}\exp{\left( \frac{e(\varphi_\mathrm{F}-\varphi(x))}{k_\mathrm{B}T} \right)}
            \end{aligned}
        \end{equation}
        Y si definimos
        \begin{equation}
        p_\mathrm{eq}=n_i\exp{\frac{e \varphi_\mathrm{F}}{k_\mathrm{B}T}}
        \end{equation}
        Por lo tanto como hemos dicho $n(x)\geq p_\mathrm{eq} \xrightarrow{} \varphi(x) \geq 2\varphi_\mathrm{F}$ que es una condición de inversión en una posición $x$ genérica, por lo tanto para que comience a haber inversión en el semiconductor $\varphi_\mathrm{S}\geq 2\varphi_\mathrm{F}$.\\\\ Si estudiamos las cargas nos encontramos 
        \begin{equation}
\begin{aligned}
& Q_\mathrm{a c} \propto \exp \left[e\left|\varphi_S\right| / k_B T\right] \\
& Q_\mathrm{i n v} \propto \exp \left[e\left|\varphi_S\right| / k_B T\right] \\
& Q_\mathrm{v a c} \propto\left|\varphi_S\right|^{1 / 2}
\end{aligned}
\label{eq: carga en funcion de los estados de carga}
\end{equation}
    \end{itemize}
\end{itemize}
\subsubsection{Tensión umbral}
Definimos la tensión umbral $V_\mathrm{T}$ como la tensión que hay que aplicar al metal para que comience a haber inversión en el semiconductor. A partir de la tensión umbral 
\begin{equation}
    V_\mathrm{T}=V_\mathrm{OX}-V_\mathrm{OO^\prime} +\varphi_\mathrm{S}-\varphi_\mathrm{SO}  \ \text{con} \ +\varphi_\mathrm{S}=2+\varphi_\mathrm{F}  \ \text{y} \ V_\mathrm{OO^\prime}+\varphi_\mathrm{SO}=-\frac{\phi_\mathrm{MS}}{e} \label{eq: potencial de contacto definicion}
\end{equation}
Recordamos que el óxido se comporta como un condensador 
\begin{equation}
V_\mathrm{OX}=E_\mathrm{OX}d=\frac{Q_\mathrm{G}}{\epsilon_\mathrm{OX}}d=\frac{\pm Q_\mathrm{S \ max}}{C_\mathrm{OX}} \ \text{con} \ C_\mathrm{OX}=\frac{\epsilon_\mathrm{OX}}{d}\label{eq: capacidad y potencial del oxido}
\end{equation}
Siendo $Q_\mathrm{S \ max}$ la carga máxima de la capa de vaciamiento en el semiconductor (cuando la caída de potencial de $2\varphi_\mathrm{F}$).\\\\
A partir de la carga máxima podemos calcular la anchura máxima como una función de $\varphi_\mathrm{F}$ en efecto 
\begin{equation}
    Q_\mathrm{S \ max}=-eN_\mathrm{A} W_\mathrm{max}= -eN_\mathrm{A} W(2\varphi_\mathrm{F})=-(2eN_\mathrm{A}\epsilon_\mathrm{S} 2\varphi_\mathrm{F})^{1/2} \xrightarrow{}  \boxed{W_\mathrm{max}=-\left( \frac{2\epsilon_\mathrm{S}|2\varphi_\mathrm{F}|}{eN_\mathrm{A}} \right)^{1/2}} \label{eq: anchura maxima}
\end{equation}
Obteniendo finalmente que el potencial de contacto lo podemos expresar como \begin{equation}
    V_\mathrm{T}=-\frac{ Q_\mathrm{S \ max}}{ C_\mathrm{OX}}+2\varphi_\mathrm{F}+\frac{\phi_\mathrm{MS}}{e}
    \label{eq: Potencial de contacto capacida oxido}
\end{equation}
\subsection{Estructura real MOS}
Hasta ahora, cuando hemos estudiado la estructura MOS hemos supuesto que el óxido es un medio aislante con cargas estáticas, es decir, que simplemente era una zona no conductora en la estructura. Sin embargo, sí que tenemos una presencia de cargas en el óxido, concretamente se distribuyen en la superficie de separación con el semiconductor, si las denotamos como $Q_\mathrm{SS}$ podemos determinar la caída de tensión en el óxido.
\begin{equation}
V_\mathrm{SS}=\frac{Q_\mathrm{SS}}{\epsilon_\mathrm{OX}}d=\frac{Q_\mathrm{SS}}{C_\mathrm{OX}}
    \label{eq: Caida de tension en el oxido}
\end{equation}
Esta caída de potencial, influye también en la tensión umbral de nuestro transistor, por lo que si tenemos en cuenta la contribución queda como 
\begin{equation}
V_\mathrm{T}=2\varphi_\mathrm{F}-\frac{Q_\mathrm{SS}}{C_\mathrm{OX}}+\frac{\phi_\mathrm{MS}}{e}-V_\mathrm{SS} \left \{\begin{aligned}
    &V_\mathrm{SS}>0 \xrightarrow{} V_\mathrm{T} \downarrow; \ Q_\mathrm{SS}>0 \ \text{atrae} \ e^-\\
    &V_\mathrm{SS}<0 \xrightarrow{} V_\mathrm{T} \uparrow ; \ Q_\mathrm{SS}<0 \ \text{atrae} \ h^-\\
\end{aligned} \right.
\label{eq: tension umbral real}
\end{equation}
El potencial de banda plana será
\begin{equation}
    V_\mathrm{FB}=\frac{\phi_\mathrm{MS}}{e}-V_\mathrm{SS} \ \text{genéricamente} \ V_\mathrm{G}=\textcolor{blue}{-\frac{Q_\mathrm{S}}{C_\mathrm{OX}}}+ \textcolor{red}{\varphi_\mathrm{S}+ \frac{\phi_\mathrm{MS}}{e}-V_\mathrm{SS}}= -\frac{Q_\mathrm{S}}{C_\mathrm{OX}}+ \varphi_\mathrm{S}+V_\mathrm{FB}
\end{equation}
Que tiene en cuenta la expresión del modelo \textcolor{red}{ideal} y la \textcolor{blue}{caída en el óxido}.
\section{Comportamiento dinámico. Capacidad de la estructura MIS}
Para estudiar la capacidad de la estructura MIS, supongamos que tenemos una estructura similar a la figura \ref{fig:Condensador Mos}, de área unidad y con $\phi_\mathrm{MS}=0; \ Q_\mathrm{SS}=0 \Rightarrow V_\mathrm{FB}=0$, con estas condiciones, obtenemos $$V_\mathrm{G}=V_\mathrm{OX}+\varphi_\mathrm{S} \hspace{2cm}
Q_\mathrm{S}=Q_\mathrm{S}(\varphi_\mathrm{S}) \hspace{2cm} 
Q_\mathrm{G}=\frac{V_\mathrm{OX}}{d}\epsilon_\mathrm{OX}$$
Utilizando la definición de capacidad, obtenemos
\begin{equation}
    C_\mathrm{OX}=\diff[]{Q_\mathrm{G}}{V_\mathrm{OX}}=\frac{\epsilon_\mathrm{OX}}{d} \hspace{3cm} C_\mathrm{SC}=-\diff[]{Q_\mathrm{S}}{\varphi_\mathrm{S}}
\end{equation}
\newpage
\begin{itemize}
    \item $C_\mathrm{OX}$: Es una capacidad constante y depende de la geometría de la estructura MIS. \textit{En la figura \ref{fig:Condensador Mos} es un condensador planoparalelo de espesor d y de dieléctrico el óxido}.
    \item $C_\mathrm{SC}$: Las cargas del semiconductor son función de la caída de tensión en él. En algunos casos depende de la frecuencia.
\end{itemize}
A partir de estas definiciones podemos definir una capacidad total, entre los extremos del condensador $C_\mathrm{T}$ que será
\begin{equation}
    C_\mathrm{T}=\left(\diff[]{(V_\mathrm{OX}+\varphi_\mathrm{S})}{Q_\mathrm{G}} \right)^{-1}=\left(\frac{1}{C_\mathrm{OX}}+\frac{1}{C_\mathrm{SC}} \right)^{-1}
\end{equation}
Es decir, que la estructura MIS la podemos representar como \textbf{dos condensadores en serie}.
\subsection{Distribuciones de carga en fucnión del estado de carga de la estructura}
Si nos fijamos en las ecuaciones \ref{eq: carga en funcion de los estados de carga}, ver cómo varían las distintas capacidades en función del potencial de inversión. En las figuras el lado izquierdo representa al metal y el derecho al semiconductor.\\\\
\begin{minipage}{0.55 \textwidth}
 Supongamos ahora que tenemos un semiconductor tipo P en \textbf{acumulación}, la capacidad que predomina es la del óxido por lo que
    \begin{equation}
    \begin{aligned}
    &Q_\mathrm{a c} \propto \exp \left[e\left|\varphi_S\right| / k_B T\right] \xrightarrow[]{} \\ &\xrightarrow[]{} C_\mathrm{SC} \uparrow \uparrow \gg C_\mathrm{OX} \xrightarrow[]{} C_\mathrm{T}\approx C_\mathrm{OX}
    \end{aligned}
\end{equation}
Como están en serie, la capacidad predominante será aquella que tiene menor valor. Destacar que en la figura \ref{fig: Dist Carga acumulacion} la carga cuando nos encontramos en acumulación, se encuentra en la superficie, esto se debe a que al curvarse hacia arriba las bandas, tendremos una acumulación de huecos en la banda de valencia.
\end{minipage}
\hfill 
\begin{minipage}{0.4\textwidth}
\centering
    \includegraphics{Standalone/T5_11_Dist_Carga_acumulacion.pdf}
        \captionof{figure}{Distribución de carga en acumulación}
        \label{fig: Dist Carga acumulacion}
\end{minipage}
\\\\\\
\begin{minipage}{0.55 \textwidth}
 Supongamos ahora que tenemos un semiconductor tipo P en \textbf{vaciamiento}, ahora la capacidad del semiconductor disminuye, pues sigue la condición 
    \begin{equation}
    \begin{aligned}
    & Q_\mathrm{v a c} \propto\left|\varphi_S\right|^{1 / 2} \xrightarrow[]{} C_\mathrm{SC} \propto |\varphi_\mathrm{S}|^{-1/2} \\ &\xrightarrow[]{} C_\mathrm{SC} \ \text{del orden de } C_\mathrm{OX} \xrightarrow[]{} C_\mathrm{T} \downarrow \ \text{si} \ V_\mathrm{G} \uparrow
    \end{aligned}
\end{equation}
En la figura \ref{fig: Dist Carga vaciamiento} vemos que la carga en el semiconductor se distribuye en el volumen y es de signo contrario al del tipo del semiconductor, es la zona de carga espacial inducida.
\end{minipage}
\hfill 
\begin{minipage}{0.4\textwidth}
\centering
    \includegraphics{Standalone/T5_12_Dist_Carga_vaciamiento.pdf}
        \captionof{figure}{Distribución de carga en vaciamiento}
        \label{fig: Dist Carga vaciamiento}
\end{minipage}
\newpage
\begin{minipage}{0.55 \textwidth}
 Supongamos ahora que tenemos un semiconductor tipo P en \textbf{inversión}, la capacidad total en este caso vuelve a predominar la del óxido ya que
    \begin{equation}
    \begin{aligned}
    &Q_\mathrm{a c} \propto \exp \left[e\left|\varphi_S\right| / k_B T\right] \xrightarrow[]{} \\ &\xrightarrow[]{} C_\mathrm{SC} \uparrow \uparrow \gg C_\mathrm{OX} \xrightarrow[]{} C_\mathrm{T}\approx C_\mathrm{OX}
    \end{aligned}
\end{equation}
En la figura \ref{fig: Dist Carga vaciamiento} vemos que la carga en el semiconductor se distribuye tanto en volumen como en superficie (capa de inversión) y es de signo contrario al del tipo del semiconductor.
\end{minipage}
\hfill 
\begin{minipage}{0.4\textwidth}
\centering
    \includegraphics{Standalone/T5_13_Dist_Carga_inversion.pdf}
        \captionof{figure}{Distribución de carga en vaciamiento}
        \label{fig: Dist Carga inversión}
\end{minipage}
\\\\
Por último vamos a graficar los estados de carga de la estructura MIS en función del tipo de semiconductor que la forme.
\begin{figure}[h!]
    \centering
    \includegraphics{Standalone/T5_14_Estados_de_carga_MOS_N.pdf}
    \caption{Estados de carga de una estructura MIS cuando el semiconductor es tipo N ($\varphi_\mathrm{F}<0$) }
    \label{fig:Estados de carga MOS-N}
\end{figure}
\begin{figure}[h!]
    \centering
    \includegraphics{Standalone/T5_15_Estados_de_carga_MOS_P.pdf}
    \caption{Estados de carga de una estructura MIS cuando el semiconductor es tipo P ($\varphi_\mathrm{F}>0$) }
    \label{fig:Estados de carga MOS-P}
\end{figure}
\section{Transistor MOSFET}
El transistor MOSFET está formado por una placa de metal y un semiconductor separados por una zona de óxido de semiconductor que actúa como aislante. Tiene dos regiones muy dopadas de tipo contrario al del semiconductor que forma su sustrato.
Los electrodos que están conectado a los terminales:
\begin{itemize}
    \item \textbf{Puerta} (G): que se conecta a la placa metálica. La corriente en la puerta es nula.
    \item \textbf{Fuente} (S) y \textbf{drenador} (D): Ambos son simétricos entre sí y están conectados a las zonas muy dopadas.
    \item \textbf{Sustrato} (B): Es un terminal que puede ser suprimido y se conecta a la placa de metal que se encuentra en la cara inferior del semiconductor.
\end{itemize}
\begin{figure}[h]
    \centering
    \begin{subfigure}[h]{0.45\linewidth}
    \centering
       \includegraphics{Standalone/T5_16_Transistor_Mosfet_CNSP.pdf}
    \caption{Esquema de un transistor MOSFET canal N y sustrato P.}
    \label{fig:Transistor Mosfet CNSP} 
    \end{subfigure}
    \begin{subfigure}[h]{0.45\linewidth}
    \centering
       \includegraphics{Standalone/T5_17_Transistor_Mosfet_CPSPN.pdf}
    \caption{Esquema de un transistor MOSFET canal P y sustrato N.}
    \label{fig:Transistor Mosfet CPSPN} 
    \end{subfigure}
    \caption{Transistores MOSFET.}
    \label{fig: Tipos de MOSFET}
\end{figure}
La \textbf{idea de funcionamiento} es: controlar la corriente que fluye entre S y D mediante la tensión aplicada a G, para ello es necesario que tengamos capa de inversión debajo del óxido.
En función del tipo de MOSFET que tengamos tendremos distintos rangos de valores para las tensiones y los terminales\newpage
\begin{minipage}{0.5\textwidth}
\textbf{MOSFET canal N y sustrato P} $$I_\mathrm{G}=0 \quad V_\mathrm{GS}>V_\mathrm{T} \quad V_\mathrm{DS}>0 \quad I_\mathrm{D}>0$$    
\end{minipage}
\hfill
\begin{minipage}{0.5\textwidth}
     \textbf{MOSFET canal P y sustrato N} $$I_\mathrm{G}=0 \quad V_\mathrm{GS}<V_\mathrm{T} \quad V_\mathrm{DS}<0 \quad I_\mathrm{D}<0$$
\end{minipage}
\subsection{Comportamiento estático del MOSFET}
Supongamos que en nuestro MOSFET es canal N y sustrato P, y las siguientes características: $\phi_\mathrm{MS}=0; \ Q_\mathrm{SS}=0 \Rightarrow V_\mathrm{FB}=0$. También nos encontramos con que la fuente y el sustrato están cortocircuitados. Las corrientes y los potenciales son 
$$I_\mathrm{G}=0 \quad V_\mathrm{GS}>V_\mathrm{T} \quad V_\mathrm{DS}>0 \quad I_\mathrm{D}>0$$
Tomamos como referencia la fuente, por lo que estará conectada a tierra junto con el sustrato y los límites del canal se encuentran en $y=0, \quad y=L$
\subsubsection{MOSFET en equilibrio}
Si tenemos que $V_\mathrm{GS}=V_\mathrm{DS}=0, \ I_\mathrm{D}=0$ y con las asumciones anteriores nos encontramos en \textbf{situación de banda plana}.
\begin{figure}[h!]
    \centering
        \includegraphics{Standalone/T5_18_Diagrama_de_bandas_de_un.pdf}
    \caption{Diagrama de bandas de un MOSFET en equiblibrio con $V_\mathrm{GS}=V_\mathrm{DS}=0, \ I_\mathrm{D}=0$}
    \label{fig:enter-label}
\end{figure}
\newpage
\noindent Si tenemos que  $V_\mathrm{GS}>0, \ V_\mathrm{DS}=0, \ I_\mathrm{D}=0$ nos encontramos que el nivel de Fermi es constante con la posición, es decir \textbf{si la corriente de drenador es nula el nivel de Fermi es constante con la posición}.

\begin{figure}[h]
    \centering
    \includegraphics[width=0.5\linewidth]{Standalone/T5_20_Diagrama_de_bandas_VGS_positivo.pdf}
    \caption{Diagrama de bandas de energía al aplicar un potencial de rejilla.}
    \label{fig:Diagrama de bandas para VGS>0}
\end{figure}

\begin{wrapfigure}{r}{0.5\textwidth}
    \flushright
    \includegraphics{Standalone/T5_21_Diagrama_de_bandas_VGS_VDS_positivos.pdf}
    \caption{Diagrama de bandas de energía al aplicar un potencial de rejilla y una tensión drenador-fuente positivas}
    \label{fig:Diagrama de bandas para VGS>0 y VDS>0}
\end{wrapfigure}

Finalmente si tenemos corriente de drenador, el nivel de Fermi ya no es constante, por lo que nos aparecen dos cuasiniveles de Fermi $E_{\mathrm{F}N}$ y $E_{\mathrm{F}P}$ como $J_P=0$ el nivel de Fermi es constante en el canal. Esto se debe a que en el canal no hay corriente de huecos por lo que el nivel de Fermi es constante con la posición.
El nivel de Fermi del sustrato, será 
\begin{equation}
\begin{aligned}
    &E_\mathrm{FN}=E_\mathrm{FP}-eV(x,y) \quad \text{con}\\ &V(x,y)=\left \{ \begin{aligned}
        &V(x,0)=0\\
        &V(x,L)=V_\mathrm{DS}
    \end{aligned}\right.
    \end{aligned}
    \label{eq: Nivel de Fermi sustrato}
\end{equation}
Por lo que $eV(x,y)$ es la separación entre cuasiniveles de Fermi. Además cuando $y=0$, la unión se encontrará en directa mientras que cuando $y=L$ la unión se encontrará en inversa.\\\\\\\\
Si estudiamos la corriente tendremos que 
\begin{equation}
    J_n=n\mu_n \nabla E_\mathrm{FN}(x,y)=-en\mu_n \nabla V(x,y)
    \label{eq: Corriente de arrastre sustrato}
\end{equation}
Debido a la presencia del óxido $I_\mathrm{G}=0 \Rightarrow J_x=0$ por lo que $$V(x,y)=V(y)$$ Si recordamos la definición de \textbf{condición de inversión}, obtenemos \begin{equation}
\varphi_\mathrm{S}\geq 2\varphi_\mathrm{F}+V(y)
    \label{eq: Condición de inversion MOSFET}
\end{equation}
\subsubsection{Potencial umbral}
Nos vamos a referir a la tensión umbral, como la tensión $V_\mathrm{GS}$ mínima a aplicar para que exista capa de inversión en el canal (con $V_\mathrm{DS}=0$)
\begin{equation*}
V_\mathrm{T}=V_\mathrm{OX}+2\varphi_\mathrm{F}+\frac{\phi_\mathrm{MS}}{e}-V_\mathrm{SS}=-\frac{Q_\mathrm{S max}}{\epsilon_\mathrm{OX}}d+2\varphi_\mathrm{F} \ \text{con} \ Q_\mathrm{S max}=-\sqrt{2eN_\mathrm{A}\epsilon_\mathrm{S}2\varphi_\mathrm{F}}
\end{equation*}
Concluyendo que 
\begin{equation}
V_\mathrm{T}=2\varphi_\mathrm{F}+K(2\varphi_\mathrm{F})^{1/2} \quad \text{con} \quad K^2=\frac{2eN_\mathrm{A}\epsilon_\mathrm{S}}{C^2_\mathrm{OX}}
\label{eq: tension umbral MOSFET canal n}
\end{equation}
Ahora podemos calcular la tensión necesaria para comenzar la inversión en una posición $y$ genérica (con $V_\mathrm{DS}\neq 0$) que tiene como expresión
\begin{equation}
V_\mathrm{GS}=2\varphi_\mathrm{F}+V(y)+K(2\varphi_\mathrm{F}+V(y))^{1/2}
\label{eq: Tension VGS generica para dar inversion}
\end{equation}
Se diferencia de la ecuación \ref{eq: tension umbral MOSFET canal n} en que añadimos el término de separación entre los cuasiniveles de Fermi.
\subsubsection{Características $I_\mathrm{D}-V_\mathrm{DS}$. Equilibrio, no saturación, saturación}
\begin{wrapfigure}{r}{0.35\textwidth}
\flushright
    \includegraphics{Standalone/T5_19_Caracteristica_ID_VDS.pdf}
    \caption{Característica $I_\mathrm{D}-V_\mathrm{DS}$ de un transistor MOSFET.}
    \label{fig: Caracteristica ID-VDS}
\end{wrapfigure}
Si tenemos una capa de inversión uniforme nos encontramos con que $V_\mathrm{DS}=0$, además la zona de vaciamiento se encontrará entorno al drenador.\\\\
Si nos encontramos al borde de obtener inversión, es decir $V_\mathrm{DS}=V_\mathrm{Dsat}$, el potencial aumenta conforme nos acercamos al drenador, la zona de vaciamiento aumenta. Si el potencial cae más que una cantidad $2\varphi_\mathrm{F}$, comienza a caer en la zona de vaciamiento.\\\\
Si tenemos $V_\mathrm{DS}>V_\mathrm{Dsat}$, vemos que existe una separación entre la capa de inversión y el drenador. El orígen de la corriente de saturación es que los pocos electrones que hay en $L^\prime$, y que pasan al drenador, son acelerados por el campo que es muy intenso en esa pequeña separación, por lo que mantiene la corriente constante. La expresión para la separación es 
\begin{equation}
\Delta L=L-L^{\prime}=\left[\frac{2 \varepsilon_\mathrm{S}\left(V_\mathrm{DS}-V_{\text {Dsat }}\right)}{e N_\mathrm{A}}\right]^{1 / 2}
\end{equation}
En la figura \ref{fig: Caracteristica ID-VDS}, podemos ver quecuando el potencial de puerta es menor que la tensión umbral, la intensidad de drenador crece de manera lineal. Y comienza a curvarse más y más a medida que crece $V_\mathrm{GS}$ hasta que alcanzamos el potencial de saturación en el que se mantiene constante la corriente.
\newpage
\subsubsection{Tensión de Saturación}
Dada una $V_\mathrm{GS}$ constante, el punto $y$ del canal que está justo al borde de inversión siempre cumple la condición dada por \ref{eq: Tension VGS generica para dar inversion}. La saturación comienza en $y=L \Rightarrow V(L)=V_\mathrm{DS}=V_\mathrm{DSat}$ De la expresión \ref{eq: Tension VGS generica para dar inversion} obtenemos
\begin{equation}
V_\mathrm{GS}=2\varphi_\mathrm{F}+V(L)+K(2\varphi_\mathrm{F}+V(L))^{1/2}=2\varphi_\mathrm{F}+V_\mathrm{DSat}+K(2\varphi_\mathrm{F}+V_\mathrm{DSat})^{1/2}
\label{eq: Tension de puerta al comienzo de saturacion}
\end{equation}
El potencial de puerta será entonces
\begin{equation}
    V_\mathrm{DSat}=V_\mathrm{GS}-2\varphi_\mathrm{F}+\frac{K^2}{2}\left(1-\sqrt{1+\frac{4V_\mathrm{GS}}{K^2}}\right) \ \text{con} \ K^2=\frac{2eN_\mathrm{A}\epsilon_\mathrm{S}}{C^2_\mathrm{OX}}
    \label{eq: Tension de saturacion}
\end{equation}
Si $K$ es pequeña tendremos $V_\mathrm{T}\approx 2\varphi_\mathrm{F}$ y $V_\mathrm{DSat}\approx V_\mathrm{GS}-V_\mathrm{T}$.
\subsection{Expresiones analíticas para características de salida: $I_\mathrm{D}-V_\mathrm{DS}$ con $V_\mathrm{GS}$ constante}
Utilizaremos la aproximación de canal gradual en el que $\partial^2_{xx} \varphi \gg \partial^2_{yy} \varphi \Rightarrow \partial_xE_x \gg \partial_yE_y \Rightarrow E_y\approx \text{cte}$. Si \textbf{no estamos en saturación} la corriente de drenador será

\begin{equation}
    I_\mathrm{D}=\mu_n \textcolor{red}{\frac{Z C_\mathrm{OX}}{L}} \left [ \left( V_\mathrm{GS}-2\varphi_\mathrm{F} -\frac{V_\mathrm{DS}}{2}\right)V_\mathrm{DS}-\frac{2}{3}K\left( (V_\mathrm{DS}+2\varphi_\mathrm{F})^{3/2}- (2\varphi_\mathrm{F})^{3/2} \right ) \right ]
    \label{eq: Corriente de drenador en no saturacion general}
\end{equation}
Siendo $Z$ la anchura del transistor y $L$ la longitud del canal. Que depende de un \textcolor{red}{\textbf{factor geométrico}} y observamos que si $I_\mathrm{D}(V_\mathrm{DS}=0)=0$ y $V_\mathrm{DS} \uparrow \Rightarrow I_\mathrm{D} \uparrow$ y $V_\mathrm{GS} \uparrow \Rightarrow I_\mathrm{D} \uparrow$.\\\\
Si nos encontramos en \textbf{saturación} la corriente de drenador será  
\begin{equation}
    I^{'}_\mathrm{D}=\frac{Z\mu_nC_\mathrm{OX}}{L^{'}} \left [ \left( V_\mathrm{GS}-2\varphi_\mathrm{F} -\frac{V_\mathrm{DSat}}{2}\right)V_\mathrm{DSat}-\frac{2}{3}K\left( (V_\mathrm{DSat}+2\varphi_\mathrm{F})^{3/2}- (2\varphi_\mathrm{F})^{3/2} \right ) \right ]
    \label{eq: Corriente de drenador en saturacion general}
\end{equation}
Justo en la unión de las dos zonas $L=L^{'}$, por lo que $I_\mathrm{D}(V_\mathrm{DSat})=I^{'}_\mathrm{D}(V_\mathrm{DSat}) \xrightarrow{}  V_\mathrm{GS} \uparrow \Rightarrow I_\mathrm{D} \uparrow$ que depende de $V_\mathrm{DS}$ a través de $L^{'}$. Recordamos que cuando estamos en saturación la anchura del canal disminuye.\\\\
Si tenemos que $V_\mathrm{DS}$ es pequeño ó $K$ es pequeño, no podremos encontrarnos en saturación y la corriente de drenador es
\begin{equation}
    I_\mathrm{D}=\mu_n \frac{Z C_\mathrm{OX}}{L} \left [ ( V_\mathrm{GS}-V_\mathrm{T})V_\mathrm{DS}-\frac{V_\mathrm{DS}^2}{2} \right ]
    \label{eq: Corriente de drenador para VDS pequeño o K pequeño}
\end{equation}
Si nos encontramos que $V_\mathrm{DS}$ es muy pequeño, el comportamiento de la corriente de drenador será lineal y lo podemos controlar con la tensión de rejilla. La corriente de drenador será
\begin{equation}
        I_\mathrm{D}=\mu_n \frac{Z C_\mathrm{OX}}{L} ( V_\mathrm{GS}-V_\mathrm{T})V_\mathrm{DS}
        \label{eq: Corriente de drenador para VDS muy pequeño}
\end{equation}
Si sólo tenemos $K$ pequeño y nos encontramos en saturación la corriente de drenador será
\begin{equation}
    I^{'}_\mathrm{D}=\frac{Z\mu_nC_\mathrm{OX}}{L^{'}} \frac{(V_\mathrm{GS}-V_\mathrm{T})^2}{2}
    \label{eq: Corriente de drenador para K pequeño en saturaicon}
\end{equation}
Por último vamos a definir los siguientes parámetros
\begin{itemize}
    \item \textbf{Conductancia de salida} 
    \begin{equation}
        g_\mathrm{D}= \left. \pdiff[]{I_\mathrm{D}}{V_\mathrm{DS}} \right |_{V_\mathrm{GS}=\mathrm{cte}} \quad g^{'}_\mathrm{D}= \left. \pdiff[]{I^{'}_\mathrm{D}}{V_\mathrm{DS}} \right |_{V_\mathrm{GS}=\mathrm{cte}}
        \label{eq: conductancia de salida}
    \end{equation}
    \item \textbf{Transconductancia}:
    \begin{equation}
        g_\mathrm{m}=\left. \pdiff[]{I_\mathrm{D}}{V_\mathrm{GS}} \right |_{V_\mathrm{DS}=\mathrm{cte}}=\frac{Z\mu_nC_\mathrm{OX}}{L}V_\mathrm{DS} \quad g^{'}_\mathrm{m}= \left. \pdiff[]{I^{'}_\mathrm{D}}{V_\mathrm{GS}} \right |_{V_\mathrm{DS}=\mathrm{cte}}=\frac{Z\mu_nC_\mathrm{OX}}{L^{'}}(V_\mathrm{GS}-V_\mathrm{T})
    \end{equation}
\end{itemize}
Siendo $g^{'}_\mathrm{m}$ el parámetro más importante del mosfet pues os indica si es un buen transistor o no.
\subsection{Comportamiento dinámico. Señales débiles de baja frecuencia}
Cuando tenemos señales débiles tenemos tres contribuciones para la corriente de drenador
\begin{equation}
i_\mathrm{d}=\left.\frac{\partial I_\mathrm{D}}{\partial V_\mathrm{G S}}\right|_{\substack{V_\mathrm{D S}=\mathrm{cte} \\ V_\mathrm{B S}=\mathrm{cte}}} v_\mathrm{g s}+\left.\frac{\partial I_\mathrm{D}}{\partial V_\mathrm{D S}}\right|_{\substack{V_\mathrm{G S}=\mathrm{cte} \\ V_\mathrm{B S}=\mathrm{cte}}} v_\mathrm{d s}+\left.\frac{\partial I_\mathrm{D}}{\partial V_\mathrm{B S}}\right|_{\substack{V_\mathrm{D S}=\mathrm{cte} \\ V_\mathrm{G S}=\mathrm{cte}}} v_\mathrm{b s}
\end{equation}
La expresión depende de las conductancias y de las transconductancias.
\begin{equation}
i_\mathrm{d}=g_\mathrm{m}v_\mathrm{gs}+g_\mathrm{d}v_\mathrm{ds}+g_\mathrm{mb}v_\mathrm{mb}  
\end{equation}
En este caso podemos representarlo por un circuito equivalente
\begin{figure}[h]
    \centering
    \begin{subfigure}[h]{0.45\linewidth}
    \centering
    \includegraphics{Standalone/T5_22_Circuito_MOSFET_baja_frecuencia.pdf}
    \caption{Representación circuital de un MOSFET con señal de baja frecuencia.}
    \end{subfigure}
    \hfill
    \begin{subfigure}[h]{0.45\linewidth}
    \centering
        \includegraphics{Standalone/T5_23_Circuito_MOSFET_alta_frecuencia.pdf}
        \caption{Representación circuital de un MOSFET con señal de alta frecuencia.}
    \end{subfigure}
    \caption{Representación circuital del MOSFET.}
    \label{fig: Representacion circuito MOSFET}
\end{figure}
\section{MOSFET real}
Hasta ahora no hemos tenido en cuenta el efecto de la caída de potencial en el metal y el semiconductor; ni tampoco el efecto de las cargas en el óxido $V_\mathrm{SS}$, para ello tenemos que sustituir el potencial de puerta $$V_\mathrm{G}\xrightarrow[]{} V_\mathrm{GS}-V_\mathrm{FB}=V_\mathrm{GS}-\left ( \frac{\phi_\mathrm{MS}}{e}-V_\mathrm{SS}\right )$$
También tenemos que tener en cuenta el efecto de la polarización del sustrato ($V_\mathrm{BS}\neq 0$) que hasta ahora lo consideramos nulo. Para el caso de un transistor canal N\footnote{Uniones PN$^+$ B-S y B-D en inversa, $I_\mathrm{B}=0$} $V_\mathrm{BS}<0$ por la forma en la que tenemos conectado el MOSFET.

En este caso tenemos que sustituir $$
V_\mathrm{GS} \xrightarrow[]{} V_\mathrm{GS}-V_\mathrm{BS}=V_\mathrm{GB} \quad  \quad  2\varphi_\mathrm{F} \xrightarrow[]{} 2\varphi_\mathrm{F}-V_\mathrm{BS}$$

Con estos cambios la tensión umbral resulta
\begin{equation}
    V_\mathrm{T}=V_x+K(2\varphi_\mathrm{F}-V_\mathrm{BS})^{1/2} \ \text{con} \ V_x=2\varphi_\mathrm{F}+\frac{\phi_\mathrm{MS}}{e}-V_\mathrm{SS}
\end{equation}

Ahora podemos volver a estudiar las corrientes de drenador, como expresión general en \textbf{no saturación}
\begin{equation}
    I_\mathrm{D}=\frac{Z \mu_n  C_\mathrm{OX}}{L} \left [ \left( V_\mathrm{GS}-V_x -\frac{V_\mathrm{DS}}{2}\right)V_\mathrm{DS}-\frac{2}{3}K\left( (V_\mathrm{DS}+2\varphi_\mathrm{F}-V_\mathrm{BS})^{3/2}- (2\varphi_\mathrm{F}-V_\mathrm{BS})^{3/2} \right ) \right ]
    \label{eq: Corriente de drenador en no saturacion general real}
\end{equation}

Si tenemos que $V_\mathrm{DS}$ ó $K$ pequeño, $V_\mathrm{T}=V_x=2\varphi_F-V_\mathrm{FB}$ recuperamos la ecuación \ref{eq: Corriente de drenador para VDS pequeño o K pequeño}. Mientras que si $V_\mathrm{DS}$ es muy pequeño recuperamos \ref{eq: Corriente de drenador para VDS muy pequeño}.\\\\
Si nos encontramos en \textbf{saturación} la corriente de drenador será 
\begin{equation}
    I^{'}_\mathrm{D}=\frac{Z \mu_n  C_\mathrm{OX}}{L^{'}} \left [ \left( V_\mathrm{GS}-V_x -\frac{V_\mathrm{DS}}{2}\right)V_\mathrm{DS}-\frac{2}{3}K\left( (V_\mathrm{DS}+2\varphi_\mathrm{F}-V_\mathrm{BS})^{3/2}- (2\varphi_\mathrm{F}-V_\mathrm{BS})^{3/2} \right ) \right ]
    \label{eq: Corriente de drenador en no saturacion general real}
\end{equation}
Mientras que si $K$ es pequeño recuperamos la ecuación \ref{eq: Corriente de drenador para K pequeño en saturaicon}.\\\\
Como estas expresiones son bastante complejas, podemos aproximarlas (de manera bastante precisa) mediante
\begin{equation}
    I_\mathrm{D}=\beta \left [(V_\mathrm{GS}-V_\mathrm{T})V_\mathrm{DS}-\frac{V^2_\mathrm{DS}}{2} \right ]
    \label{eq: corriete de drenador con ganancia}
\end{equation}
Para no saturación y para saturación
\begin{equation}
    I^{'}_\mathrm{D}=\frac{\beta}{2}(V_\mathrm{GS}-V_\mathrm{T})^2
\end{equation}

\section{\textcolor{red}{Tipos de MOSFET y representación}}
Como hemos mostrado en la figura \ref{fig: Tipos de MOSFET}, tenemos dos tipos de MOSFET que los denominamos canal N y canal P. Cada uno de ellos funciona con unas características distintas, vamos a ver en qué se diferencian
\begin{itemize}
    \item \textbf{MOSFET canal N y sustrato P} Este MOSFET tiene como caracteristica $$I_\mathrm{D}-V_\mathrm{DS}; \quad I_\mathrm{D}>0 \ V_\mathrm{DS}>0$$; $\quad$  y conducen corriente si $V_\mathrm{GS}>V_\mathrm{T}$. Además para $K\approx 0$ tenemos 
    \begin{equation}
K \approx 0\left\{\begin{array}{l}
V_\mathrm{DS}<V_\mathrm{GS}-V_\mathrm{T}=V_\mathrm{DSat} \quad \text{no hay saturación} \\
V_\mathrm{DS}>V_\mathrm{GS}-V_\mathrm{T}=V_\mathrm{DSat} \quad \text{hay saturación}
\end{array}\right.
\end{equation}
    \item \textbf{MOSFET canal N y sustrato P} Este MOSFET tiene como caracteristica $$I_\mathrm{D}-V_\mathrm{DS}: \quad I_\mathrm{D}<0V_\mathrm{DS}<0$$, y conducen corriente si $V_\mathrm{GS}<V_\mathrm{T}$. Además para $K\approx 0$ tenemos 
\begin{equation}
K \approx 0\left\{\begin{array}{l}
V_\mathrm{DS}>V_\mathrm{GS}-V_\mathrm{T}=V_\mathrm{DSat} \quad \text{no hay saturación} \\
V_\mathrm{DS}<V_\mathrm{GS}-V_\mathrm{T}=V_\mathrm{DSat} \quad \text{hay saturación}
\end{array}\right.
\end{equation}
\end{itemize}
La tensión umbrál toma la expresión
\begin{equation}
V_\mathrm{T}=V_x+K(2\varphi_\mathrm{F}-V_\mathrm{BS})^{1/2} \quad \text{con} \quad V_x=2\varphi_\mathrm{F}+\frac{\phi_\mathrm{MS}}{e}-V_\mathrm{SS}
\end{equation}
$V_\mathrm{T}$ puede ser positivo o negativo y es \textit{controlable} a través de $\phi_\mathrm{MS}$ y $V_\mathrm{SS}$.
\subsubsection{Transistores de realce y vaciamiento}
Decimos que un transistor es de \textbf{realce} cuando $V_\mathrm{T}>0$ para el canal N y $V_\mathrm{T}<0$ para canal P, tenemos ausencia de capa de inversión para $V_\mathrm{GS}=0$. \\\\
Decimos que un transistor es de \textbf{vaciamiento} cuando $V_\mathrm{T}<0$ para el canal N y $V_\mathrm{T}>0$ para canal P, tenemos presencia de capa de inversión para $V_\mathrm{GS}=0$.
\begin{figure}[h]
    \centering
    \includegraphics[width=1\linewidth]{Standalone/T5_24_Tipos_de_MOSFET.pdf}
    \captionof{table}{Tabla con los tipos de MOSFET, según si son de realce o vaciamiento y qué canal tienen.}
    \label{fig:Tipos de MOSFET}
\end{figure}
\part{Dispositivos electrónicos especializados}
\pagestyle{fancy}
\fancyhead[R]{Electrónica Física: Tema 6}
\fancyhead[L]{4º Grado en Física}

\chapter{Dispositivos Optoelectrónicos}
Una vez determinadas las propiedades de los dispositivos electrónicos fundamentales, basados esencialmente en las propiedades electromagnéticas de los semiconductores. Ya hemos estudiado qué sucede cuando irradiamos con ondas electromagnéticas un semiconductor, como vimos en la sección \ref{sec: procesos de generacion-recombinacion}; o las condiciones descritas en el apéndice \ref{sec: Recombinacion superficial}. Estudiaremos entonces dispositivos basados en las \textbf{propiedades ópticas}.

\section{Dispositivos detectores de luz. Fotodetectores}
\subsection{\textcolor{red}{Fotodiodos}}
Un \textbf{fotodiodo PN} es una unión PN en la que la corriente inversa aumenta con el flujo de fotones incidente. Este dispositivo presenta las siguientes características
\begin{itemize}
    \item En las proximidades de la unión PN existe un campo eléctrico intenso.
    \item Este campo es capaz de separar rápidamente los pares de portadores generados, \textbf{disminuyendo} así la probabilidad de que se produzcan \textbf{recombinaciones} que impidan su contribución a la fotocorriente.
\end{itemize}
Los que nos interesa para un correcto funcionamiento es que el flujo de fotones incida en la zona próxima a la unión (zona de transición) para conseguir la mayor eficiencia, aunque normalmente es estrecha.\\

\begin{minipage}{0.45\textwidth}
    \centering
    \includegraphics{Standalone/T6_01_Esquema_de_un_fotodiodo.pdf}
    \captionof{figure}{Esquema de un fotodiodo.}
\end{minipage}
\hfill
\begin{minipage}{0.45\textwidth}
    Si no iluminamos la intensidad de manera general viene dada por
    \begin{equation}
        I=I_0\left(\exp{\left( \frac{eV}{k_\mathrm{B}T}\right)}-1 \right )
        \label{eq: Intensidad fotodiodo}
    \end{equation}
    Mientras que si iluminamos tenemos que sumarle la contribución de la generación de fotones
    \begin{equation}
        I=-I_\mathrm{ph}+I_0\left(\exp{\left( \frac{eV}{k_\mathrm{B}T}\right)}-1 \right )
        \label{eq: Intensidad fotodiodo con fotocorriente}
    \end{equation}
\end{minipage}
\newpage

Si incidimos con luz de manera uniforme en toda la unión generaremos una cantidad $G_\mathrm{L}$ de pares $e^-h^+$ por unidad de volumen y tiempo. La corriente que obtenemos tendrá varias contribuciones
$$J=J_\mathrm{L1}+J_{n1}+J_{p1}++J_\mathrm{ph}$$
Siendo $J_\mathrm{L1}$ la más significativa si aumentamos mucho $W$. Las expresiones son 
\begin{equation*}
    J_\mathrm{L1}=eG_\mathrm{L}W \quad J_{n1}=eG_\mathrm{L}L_n+\frac{eD_nn_{p0}}{L_n} \quad  J_{p1}=eG_\mathrm{L}L_p+\frac{eD_np_{n0}}{L_p} \quad 
    J_\mathrm{ph}=eG_\mathrm{L}(W+L_n+L_p)
\end{equation*}

Se usan en el \textbf{tercer cuadrante de las características I-V} y su diseño está orientado a lograr que la sensibilidad a la luz sea máxima. De la siguiente manera
\begin{itemize}
    \item Para mejorar las prestaciones, aumentamos el tamaño (área) de la zona de transición.
    \item El semiconductor está expuesto a la luz a través de una cobertura de cristal, a veces en forma de lente.
    \item Por su diseño y construcción será especialmente sensible a la incidencia de la luz visible o infrarroja.
\end{itemize}
\subsection{Fotodiodos PIN}
Se construyen añadiendo una región intrínseca larga de longitud $W$, donde toda la región intríseca es de zonas de vaciamiento. Gracias a añadir dicha región, reducimos la capacidad de la unión y se mejora la velocidad de respuesta. Es más sensible a las variaciones de longitudes de onda ya que no las absorbe todas con la misma eficiencia.
\subsection{\textcolor{red}{Fotodiodos de avalancha}}
Es como un fotodiodo PN o PIN pero con polarización inversa muy intensa. En este caso tenemos generación de pares $e^-h^+$ en la zona de transición. Los portadores generados por los fotones, originan pares $e^-h^+$ adicionales por \textbf{ionización por impacto}. Obteniendo \textbf{amplificación de la fotocorriente} generada por la absorción de la luz (no es útil para el tratamiento de señales débiles pues nos genera un mayor ruido que otros fotoreceptores).
\subsection{Fototransistores}
La idea de los fototransistores consiste en amplificar $I_\mathrm{ph}$, haciendo que forme parte de la corriente de base.\\\\ Denominamos \textbf{base flotante} a la base que no está polarizada, por lo que $$I_\mathrm{B}=0 \Rightarrow I_\mathrm{C}=-I_\mathrm{E}$$ De las ecuaciones de Ebers-Moll para un transistor NPN tendremos que la corriente de colector es

\begin{equation}
    I_\mathrm{C}=-\alpha_\mathrm{F}I_\mathrm{E}-I_\mathrm{CO}\left(\exp{\left(\frac{-eV_\mathrm{CB}}{k_\mathrm{B}T} \right)}-1 \right)
\end{equation}
Si iluminamos la unión colectora le debemos añadir la fotocorriente.
Si consideramos $V_\mathrm{CB}>0$ tendremos que 
\begin{equation}
    I_\mathrm{C}=-\alpha_\mathrm{F}I_\mathrm{E}-I_\mathrm{CO}+I_\mathrm{ph} \Rightarrow \boxed{I_\mathrm{C}=\frac{I_\mathrm{CO}+I_\mathrm{ph}}{1-\alpha_\mathrm{F}}=(1+\beta_\mathrm{F})(I_\mathrm{CO}+I_\mathrm{ph})}
\end{equation}
\section{\textcolor{red}{Células solares}}
Las células solares son uniones PN, que pueden ser de Si, Si amorfo u otros semiconductores, sin polarizar y sólo con una resistencia $R_\mathrm{L}$ de carga. La expresión para su corriente es la dada por la ecuación \ref{eq: Intensidad fotodiodo con fotocorriente} y la caída de potencial será $V=-IR_\mathrm{L}$.
\begin{wrapfigure}{r}{0.3\textwidth}
    \flushright
    \includegraphics{Standalone/T6_02_Caracteristica_I_V_de_una.pdf}
    \caption{Característica I-V de una célula solar.}
\end{wrapfigure}

Nos encontramos en el cuarto cuadrante, si tenemos un circuito abierto, $I=0$ tendremos una tensión $V_\mathrm{OC}=\frac{k_\mathrm{B}T}{e} \ln{\left(1+\frac{I_\mathrm{ph}}{I_\mathrm{0}} \right)}$ mientras que si tenemos un circuito cerrado $V=0$ y $I_\mathrm{L}=-I_\mathrm{ph}$.\\\\ Si calculamos la \textcolor{red}{potencia de salida} mediante
\begin{equation}
    P=IV=VI_0 \left( \exp{\left( \frac{eV}{k_\mathrm{B}T}\right)}-1\right )-I_\mathrm{ph}V
\end{equation}
Podremos calcular la tensión para obtener la máxima potencia mediante la expresión $\partial_V P=0$ que resolviendo dicha ecuación obtenemos que la intenisad para maximizar la potencia es
\begin{equation}
    I_\mathrm{mp}=I_0\left( \exp{\left( \frac{eV_\mathrm{mp}}{k_\mathrm{B}T}\right)}-1\right )-I_\mathrm{ph}
\end{equation}
También podemos definir la \textbf{eficiencia de conversión} como
\begin{equation}
    \eta=\frac{P_\mathrm{m}}{P_\mathrm{in}}=\frac{V_\mathrm{mp}I_\mathrm{mp}}{P_\mathrm{in}}
\end{equation}
\subsection{Célula solar en heterounión}
Están construidas con materiales de GAP ancho que actúa como una ventana óptica para los fotones con energía $E_\mathrm{gP}<h\nu<E_\mathrm{gN}$, que hace que se absorban en la ZCE del semiconductor de GAP estrecho, proporcionando \textbf{mayor eficiencia}.\\\\
Las heterouniones permiten localizar la absorción de los fotones que itenteresan en la ZCE de la unión PN. Los semiconductores utilzados para hacer las heterouniones son los de los grupo III y V, de ahí que tengan un coste elevado.
\subsection{Célula solar de Si amorfo}
Este dispositivo es de fabricación más barata y con mucha mayor superficie. Encontramos estados electrónicos dentro del GAP, pero con muy baja movilidad. Se pueden modificar los bordes de movilidad con impurezas. El Si amorfo presenta un coeficiente de absorción muy elevado, por lo que las \textbf{capas son muy finas}.
Es menos eficiente que el Si cristalino, pero el coste es \textbf{mucho menor}.
\section{Dispositivos emisores de luz}
\subsection{Diodos emisores de luz. LED's}
Se basan en la \textbf{luminiscencia de inyección} y están formados por una unión PN en directa, esto provoca un exceso de portadores, por lo que tendremos recombinación que nos emite luz. Si recordamos la figura \ref{fig:Espectro electromagnetico} y la sección \ref{sec: Procesos banda a banda}, $h\nu\approx \text{GAP}$. No nos interesa sólo en el visible sino que tiene otras muchas aplicaciones como por ejemplo la comunicación a través de fibras ópticas que se utilizan en el \textbf{infrarrojo}.\\\\
Para poder crear dispositivos en el rango visible ha sido necesario ``fabricar'' semiconductores con GAP's elevados.\\\\
Para obtener luz Blanca, lo podemos hacer de dos maneras, mezclando tres LED's con diferentes colores (RGB). O utilizando un LED de InGaN que nos da \textbf{luz azul}, recubierto con una capa de fósforo que nos genera el \textbf{rojo} y el \textbf{verde}. \\\\
La ventaja de los LED's es que al funcionar en polarización directa consumen mucho menos que las bombillas corrientes y apenas se calienta, por lo que tienen mayor duración.
\\\\
En los LED's se da \textbf{recombinación espontánea}, donde los fotones son emitidos aleatoriamente y en todas las direcciones (no es luz coherente). Los parámetros característicos de los LED's son
\begin{itemize}
    \item \textbf{Eficiencia cuántica interna}: Es la fracción de la corriente que circula por el diodo que produce luminiscencia. Depende de la eficiencia de inyección del diodo y del porcentaje de recombianciones radiativas frente al total.
    \item \textbf{Eficiencia cuántica externa}: Es la fracción de fotones generados que son realmente emitidos al exterior. Muchos de los fotones son reabsorbidos por el semiconductor o no llegan a salir de él por fenómeno de refracción o de reflexión.
\end{itemize}
Para optimizar ambas eficiencias hay que diseñar adecuadamente los diodos.
\subsection{Diodo LED de heterounión}
El lado N (GaAl$_{0.7}$As$_{0.3}$) (de mayor GAP) inyecta electrones que e lado P (GaAl$_{0.6}$As$_{0.4}$) (de menor GAP), donde se producen recombinaciones radiativas.\\\\
Como $h\nu \approx E_\mathrm{gP}<E_\mathrm{gN}$, los fotones son emitidos a través del semiconductor tipo N sin apenas absorción.\\\\
El material tipo N actúa como una ventana óptica y \textbf{la eficiencia cuántica externa mejora}.
\subsection{Diodo LASER}
Los LASER's son unas fuentes de luz monocromática y coherente. Los LASER funcionan a través de la \textbf{emisión estimulada}, los fotones emitidos están en fase y con la misma energía. Para que la emisión estimulada predomine sobre la emiisión espontánea, incorporamos la \textbf{cavidad resonante}.\\\\ Para que la emisión estimulada supere a la absorción la concentracion del nivel superior $n_2$ debe ser superior a la del inferior $n_1$.
\begin{itemize}
    \item La \textbf{inversión de ocupación} nos hace emitir fotones en el rango de $E_\mathrm{g}<h\nu<E_\mathrm{Fn}-E_\mathrm{Fp}$.
    \item Los modos resonantes de la \textbf{cavidad resonante óptica} vienen dados por \begin{equation}
        N \left(\frac{\lambda}{2} \right)=L
    \end{equation}
\end{itemize}
Por último cabe mencionar que el diodo láser de doble heteroestructura favorece la formación de la cavidad resonante. Existe una \textbf{corriente umbral} para que comience la emisión láser.
\pagestyle{fancy}
\fancyhead[R]{Electrónica Física: Tema 7}
\fancyhead[L]{4º Grado en Física}

\chapter{Nanodispositivos}
\section{La nanoescala. Propiedades dependientes del tamaño}
Los avances en tecnologías de fabricación permiten obtener dispositivos del orden del nanómetro. Es decir, tenemos dimensiones del orden de la longitud de onda de De Broglie. A estas escalas observamos
\begin{itemize}
    \item \textbf{Fenómenos de naturaleza cuántica}: cuantización de la energía, de la resistencia, etc...
    \item \textbf{Transporte balístico}: Los electrones apenas sufren scattering en el medio, por lo que la expresión de la velocidad como la movilidad por el campo ya no es válida ya que el scattering es muy pequeño.
    \item  Los dispositivos funcionan con principios diferentes a los dispositivos basados en materiales masivos.
    \item Buena parte de los dispositivos se basan en \textbf{heterouniones} de semiconductores con \textbf{diferente GAP}. Lo que nos permite \textbf{confinar electrones}.
\end{itemize}
El confinamiento de los electrones lo podemos modelar según las dimensiones cuantizadas: Pozo Cuántico (2D), Hilo Cuántico (1D), Punto Cuántico (0D).
\section{Nanoestructuras semiconductoras de baja dimensionalidad}
\subsection{\textcolor{red}{Ejemplo 1. Pozo cuántico de AlGaAs-GaAs-AlGaAs (pozo rectangular)}}
\noindent Se fabrica mediante crecimiento epitaxial a través de haces moleculares.\\\\
Lo consideramos un pozo cuadrado infinito de anchura $a$ (sólo en la dirección $Z$). Mientras que en el plano $XY$ consideraremos que es una partícula libre
\newpage
\begin{minipage}{0.45\textwidth}
\centering
    \includegraphics{Standalone/T7_01_Representacion_de_un_pozo_cuadrado.pdf}
    \captionof{figure}{Representación de un pozo cuadrado en las bandas de un semiconductor.}
\end{minipage}
\hfill
\begin{minipage}{0.45\textwidth}
    En la dirección $Z$ tenemos estados ligados por lo que la energía está cuantizada
    \begin{equation}
        E_n=\frac{\hbar^2 \pi^2}{2m^*a^2}n^2 \ \text{con} \ n=1,2...
    \end{equation}
    Mientras que en el plano XY la energía es característica de la partícula libre, por lo tanto la energía total es 
    \begin{equation}
        E=\frac{\hbar^2}{2m^*}(k_x^2+k_y^2)+E_n
    \end{equation}
    Lo que hace que para cada valor de \textit{n} tengamos una \textbf{subbanda de energía}.
\end{minipage}
\\
El número de estados de energía $$Z(E)=\frac{m^*}{\pi \hbar^2 a}$$. Según el número de subbandas para las que haya estados a una energía dada tendremos un número de estados múltipo entero de $Z(E)$. En 3 dimensiones la densidad de estados depende de la energía como $\sqrt{E}$.
\subsection{Ejemplo 2. Heterounión modulada en dopaje (pozo triangular)}
Para construirlo utilizamos un semiconductor de GAP ancho muy dopado y un semiconductor de GAP estrecho sin dopar. Los $e^-$ que pasan al GaAs quedan confinados en una región muy estrecha en la que el fondo de la banda de conducción tiene forma triangular.\\\\
Tenemos confinamiento en la dirección perpendicular a la heterounión. Movimiento libre en las direcciones paralelas a la misma, por lo que tenemos \textbf{subbandas de energía}.\\\\
Otras posibilidades para construir el pozo triangular son: Si-SiGe, InAlAs-InGaAs, AlGaN-GaN.\\\\
El pozo triangular es la base de los \textbf{transistores MODFET} (transistor de efecto de campo modulado en dopaje). El transporte tiene lugar en la dirección paralela a la unión (partículas libre). Al disminuir la interacción con impurezas la movilidad es mayor y se consiguen transistores más rápidos, capaces de alcanzar frecuencias de funcionamiento más altas.
\subsubsection{Heterouniones simétricas - pozos cuánticos - pozos cuánticos múltiples}
Si la anchura de la heterounión es grande, podremos tener dos pozos triangulares, lo cual hace que tengamos dos canales con $e^-$ junto a las heterouniones.
\\\\
Para anchuras pequeñas, tenemos un pozo rectangular lo cual hace que tengamos un único canal de $e^-$ ocupando el pozo.\\\\
En ambos casos los electrones pueden moverse muy rápidamente en la dirección paralela a las heterobuniones al estar lejos de las impurezas.\\\\
Para algunas aplicaciones se frabrican estructuras con muchos pozos cuánticos (PCM) independientes (suficientemente separados).
\subsection{Ejemplo 3. Superredes (minibandas)}
Es una estructura periódica con muchos pozos cuánticos de anchura $a$ separados por barreras \textbf{muy estrechas} de anchura $b$, por lo tanto, los estados permitidos no están confinados en los pozos, sino que los potenciales/estados de los pozos interaccionan entre sí.\\\\
La interacción entre pozos proporciona estados que se extienden a lo largo de toda la superred, cuyas energías forman bandas de energía permitidas separadas por bandas prohibidas.\\\\
Es similar a lo que ocurre con el potencial periódico de un cristal que da lugar a las bandas de energía, pero las bandas son mucho más estrechas (\textit{a} y \textit{b} son mucho mayores que distancias interatómicas, por lo que tenemos periodicidad conrtrolable y podemos ``diseñar bandas a medida'').  Debido a la estrechez de las bandas podemos tener \textbf{tranporte por efecto túnel}, a través de la superred.
\subsection{Ejemplo 4. Puntos cuánticos (0D) e hilos cuánticos (1D)}
Requieren técnicas de fabricación más complejas y menos controlables que el crecimiento epitaxial.
\begin{itemize}
    \item \textbf{Puntos cuánticos}: Típicamente son nanocristales que contienen muchos átomos (entre $10^4 \ \text{y} \ 10^6$)\footnote{El premio Nobel de Química 2023 se concedió a la síntesis química de puntos cuánticos.}. Las tres dimensiones están cuantizadas y tendremos \textbf{confinamiento físico o electrónico}. Los puntos cuánticos los podemos modelizar como una caja de potencial cerrada
    \begin{equation}
        E_{n_1, \, n_2, \, n_3 }=\frac{\hbar^2 \pi^2}{2m^*}\left(\frac{n_1^2}{a_x^2}+\frac{n_2^2}{a_y^2}+\frac{n_3^2}{a_z^2} \right) \ \text{con} \ n_1,n_2,n_3=1,2,3...
    \end{equation}
    \begin{itemize}
        \item Tienen propiedades similares a los átomos individuales.
        \item Por lo que tienen niveles de energía no bandas y presentan una energía de ionización.
        \item Presentan efectos interesantes como el \textbf{blocaje de Coulomb}, es un efecto que se da cuando la presencia de un $e^-$ en el pnto cuántico bloquea la entrada de otros. Lo que nos permite utilizar trnasistores de un solo $e^-$.
        \item Tienen diversas utilidades, en optoeléctronica (detectores, láseres de punto cuántico, emisión fluorescente de distintos colores) o nano electrónica (transistores de un solo $e^-$)
    \end{itemize}
    \newpage
    \item \textbf{Hilos cuánticos}: El movimiento libre sólo es en una dimensión (un ejemplo son los nanotubos de carbono). Tienen menos aplicaciones que los puntos cuánticos, ya que son más difíciles de fabricar controladamente, una posible aplicación es la de interconexiones de circuitos integrados.
\end{itemize}
\section{Propiedades de transporte de nanoestructuras}
\subsection{Tranporte paralelo a las heterouniones}
Como hemos dicho antes, la movilidad mejora al reducir el scatterring con el medio masivo con impurezas. Tiene utilidad en transistores de efecto campo.\\\\
Está limitado por los mecanismos de scatterring y la rugosidad de la superficie. La tecnología para desarrollar dispositivos con mejor movilidad está en constante mejora.
\subsubsection{Conductancia cuántica}
La conductancia cuántica la definimos como 
\begin{equation}
    G=\frac{2e^2}{h}n \ \text{con} \ n=1,2...
\end{equation}
La conductancia aumenta según las condiciones en las que se encuentre el nanodispositivo y de la características de entarda:
\begin{itemize}
    \item Si el nanodispositivo se encuentar a muy baja temperatura, conseguimos transporte balístico.
    \item Si aplicamos tensión en las puertas, afectamos a la anchura del canal lo que hace aumenta el número de subbandas qeu contribuyen a la conducción. A medida que tenemos mayor tensión tenemos un mahyor número de subbandas lo que tenemos incrementos de \textit{G}.
    \item Se observan efectos similares al \textbf{efecto Hall cuántico} (\textit{plateaus} en la tensión y resistencia Hall).
\end{itemize}
\subsection{\textcolor{red}{Diodo túnel resonante. Transporte perpendicular a las heterouniones}}
Es un dispositivo que muestra resistencia incremental negativa, por lo que es útil para aplicaciones de alta frecuencia. Lo modelizamos como dos barreras finitas de potencial, de tal manera que hace que entre barreras se formeun pozo finito. Este pozo \textbf{sólo tiene un nivel permitido} que denominamos como \textbf{nivel resonante}.\\\\
Cuando aumentamos la tensión, bajamos una de las barreras y el pozo no es plano sino \textbf{lineal}, con esto enfrentamos estados lo que nos genera una probabilidad de conseguir efecto túnel hasta llegar a una resonancia.
\newpage
\subsection{Blocaje de Coulomb. Transistor de un único electrón (SET)} 
En este dispositivo se prentende controlar la carga hasta el límite de un único electrón. Se utiliza un punto cuántico rodeado por dos barreras de potencial o túnel que lo separan de dos electrodos.\\\\
Esto nos genera una estructura con capacidad muy pequeña $(\unit[\sim 10^{-16}]{F})$ que sirve para poder detectar cambios de potencial asociados a la carga de un único electrón $\Delta V=e/C \sim \unit[1.6]{mV}$\\\\
En el caso de un condensador con \textbf{capacidad elevada}, la corriente por efecto túnel entre los electrodos suele ser lineal con la tensión.\\\\
Si el condensador es de \textbf{capacidad pequeña}, el paso de un electrón a través de una barrera requiere la energía que supone añadir una carga \textit{e} a condensador que es $$E=\frac{e^2}{2C}$$ luego hará falta aplicar una tensión \textit{V} que cumpla que $$eV>\frac{e^2}{2C}$$ Para $|V|=\frac{e}{2C}$ no puede circular corriente hacia o desde el punto crítico, a esto lo denominamos \textbf{blocaje de Coulomb}. \\\\
Para poder observar y controlar estos efectos la energía térmica $k_\mathrm{B}T$ ha de ser menor que la energía del condensador, lo cual supone que $T<e^2/2Ck_\mathrm{B}$, típicamente menor de $\unit[10]{K}$. Para trabajar a temperatura ambiente \textit{C} tendría que ser aún menor, lo que implica un menor tamaño de los dispositivos.
\subsubsection{Transistor de un único electrón (SET)}
La tensión de puerta (este terminal es un aislante, no hay corriente) $V_\mathrm{G}$ modifica la capacidad del punto cuántico y controla el número de electrones que puede haber en el punto cuántico.\\\\
Cada vez que el punto cuántico admite un electrón más hay un pulso de corriente entre fuente y drenador (\textbf{oscilaciones de Coulomb}), lo que resulta muy útil para aplicaciones lógicas.
\section{Propiedades ópticas de la nanoestructuras}
Las nanoestructuras tienen propiedades distintas a los semiconductores masivos, al presentar niveles y bandas de energía que depende de la dimensionalidad y el tamaño de las nanoestructuras. Esto \textbf{permite una cierta selección} mediante la fabricación con las dimensones adecuadas.\\\\
Las transiciones más probables se dan entre electrones y huecos con igual número cuántico \textit{n}.\\\\
Tienen mayor importancia los \textbf{efectos excitónicos} al estar más próxicmos electrones y huecos debido al confinamiento. La energía de enlace electrón-hueco es mayor, lo que les confiere mayor estabilidad y mayor relevancia en propiedades ópticas.
\subsubsection{Puntos cuánticos}
Son nanocristales obtenidos por síntesis química. Mediante su tamaño es posible desplazar en frecuencia su espectro de absorvión, lo que nos permite modificar el color de los nanocristales al aumentar el borde de absorción óptico. Una aplicación importane es la creación de pantallas.
\appendix
\thispagestyle{empty}
\chapter{Soluciones de la Ecuación de Continuidad} \label{sec: Soluciones de la Ecuacion de Continuidad}
\pagestyle{fancy}
\fancyhead[R]{Electrónica Física: Apéndice A}
\fancyhead[L]{4º Grado en Física}
La Ecuación de Continuidad que hemos desarrollado, depende altamente de las condiciones iniciales. Normalmente trabajaremos en condiciones estacionarias y despreciando el término del campo eléctrico.
\begin{equation}
    \pdiff[]{\delta n}{t}=D_n \nabla_\mathbf{r}^2 \delta n+ \mu_n \nabla_\mathbf{r} (\delta n\mathbf{E})-\frac{\delta n}{\tau_n}+g_\mathrm{L}
\end{equation}
Para ello vamos a estudiar las soluciones de diferentes casos de interés:
\begin{itemize}
    \item \textbf{Presencia de generación externa en volumen}: Para obtener esta solución, vemos que la ecuación diferencial a resolver es lineal con coeficientes constantes, por lo tanto la solución será una $\mathcal{S}=\mathcal{S_\mathrm{Homog}}+\mathcal{S_\mathrm{Part}}$ de la ecuación \ref{eq: continuidad exceso de electrones} toma entonces la forma
    \begin{equation}
        D_n\pdiff[2]{\delta n}{x} -\frac{\delta n}{\tau_n} + G=0 \hspace{0.5cm} \xrightarrow{\text{Solución}} \hspace{0.5cm} \delta n (x)=\underbrace{G\tau_n}_{\mathcal{S_\mathrm{Part}}}+\underbrace{A\exp{\frac{x}{L_n}}+B\exp{-\frac{x}{L_n}}}_{\mathcal{S_\mathrm{Homog}}}
        \label{eq: Cont Generacion en volumen}
    \end{equation}
    \item \textbf{Ausencia de generación externa}: Es una ecuación diferencial lineal, homogénea, (es la versión homogénea de la ecuación anterior)  $D_n\partial^2_x\delta n-\delta n/\tau_n=0$, y tiene como solución 
    \begin{equation}
        \delta n (x)= A\exp{\frac{x}{L_n}}+B\exp{-\frac{x}{L_n}}
        \label{eq: Cont No generacion en volumen}
    \end{equation}
    \item \textbf{Ausencia de generación-recombinación y presencia de generación externa en volumen}: Para resolver este caso vamos a suponer que la distancia del material es menor que la distancia de recombinación. 
    \begin{equation}
         D_n\pdiff[2]{\delta n}{x} + G=0 \xrightarrow{\text{Integrando dos veces}} \delta n(x)=-\frac{G}{2D_n}x^2+Ax+B
         \label{eq: Cont No gen-rec}
    \end{equation}
    \item \textbf{Ausencia de generación-recombinación y de generación externa}: Esta ecuación representa un caso en el que sólo tenemos difusión. En este caso, al eliminar la constante obtenemos como solución una recta 
    \begin{equation}
    \delta n (x)=Ax+B
    \label{eq: Cont No gen-rec no generacion en volumen}
    \end{equation}
\end{itemize}
\newpage
\section{Condiciones de contorno}
En general vamos siempre a considerar que $\delta n(x)$ es una función continua, para ello vamos a utilizar las definiciones de que hicimos en la sección 2.2.3 \footnote{Definimos \textbf{contacto metal/óhmico} a aquella disrrupción en la que $s\rightarrow{} \infty$ y \textbf{superficie ideal} a aquella en la que $s\rightarrow{} 0$.}. Distinguimos dos condiciones
\begin{itemize}
    \item \textbf{Valor fijo del exceso en $x_0$}, el valor de $\delta n(x_0)=\delta n_{x_0}$ es constante y conocido.
    \item \textbf{Contacto óhmico en $x_0$}: no tenemos exceso en $x_0$. $\delta n(x_0)=0$
\end{itemize}
\subsection{Recombinación superficial en $x_0$} \label{sec: Recombinacion superficial}
Primero ilustraremos los flujos de corriente que hay en las distintas regiones de la superficie del semiconductor
\begin{figure}[h]
    \centering
    \begin{subfigure}[h]{0.3\linewidth}
\includegraphics{Standalone/ApA_01_flujo_de_corriente_i.pdf}
        \caption{Diagrama del flujo de corriente en la región del límite izquierdo de un semiconductor}
        \label{fig: flujo de corriente i}
    \end{subfigure}
    \hfill
    \begin{subfigure}[h]{0.3\linewidth}
            \includegraphics{Standalone/ApA_02_flujo_de_corriente_g.pdf}
        \caption{Diagrama de flujo de corriente en una región interna del semiconductor}
        \label{fig: flujo de corriente g}
    \end{subfigure}
    \hfill
    \begin{subfigure}[h]{0.3\linewidth}
    \includegraphics{Standalone/ApA_03_flujo_de_corriente_d.pdf}
    \caption{Diagrama del flujo de corriente en la región del límite derecho de un semicondutor}
    \label{fig: flujo de corriente d}
    \end{subfigure} 
    \caption{Diagrama de flujos de corriente en distintas regiones de un semiconductor}
    \label{fig: flujo de corriente}
\end{figure}

Recordar que las corrientes generadas por la recombinación superficial, son \textbf{corrientes de difusión}. Vamos a fijar el sentido de la corriente de izquierda a derecha, y estudiaremos dichas condiciones en el punto $x=x_0$.
\subsubsection{Caso general. Posición en el interior de la muestra}
\noindent Está representado por la figura \ref{fig: flujo de corriente g}. Por conservación de la corriente 
\begin{equation}
J_\mathrm{dI}(x_0)=J_\mathrm{rec}(x_0)+J_\mathrm{dII}(x_0) \left \{
\begin{aligned}
    &e^- \xrightarrow{} eD_n \left. \diff[]{n_\mathrm{I}}{x}\right|_{x=x_0}=-es\delta n (x_0)+eD_n\left. \diff[]{n_\mathrm{II}}{x}\right|_{x=x_0} \\
    &h^+ \xrightarrow{} -eD_p \left. \diff[]{p_\mathrm{I}}{x}\right|_{x=x_0}=es\delta p (x_0)-eD_p\left. \diff[]{p_\mathrm{II}}{x}\right|_{x=x_0}
\end{aligned}
\right.
\end{equation}
Vemos que tenemos signo + en los $e^-$ y signos - en los $h^+$ en las corrientes de difusión.
\subsubsection{Casos en los límites izquierdo y derecho.}
Están representados por la figura \ref{fig: flujo de corriente i} y \ref{fig: flujo de corriente d}. En estos casos sólo tenemos una corriente de difusión, que en el caso izquierdo describen
\begin{equation}
   0=J_\mathrm{rec}(x_0)+J_\mathrm{d}(x_0) \left \{
\begin{aligned}
    &e^- \xrightarrow{} 0=-es\delta n (x_0)+eD_n\left. \diff[]{n}{x}\right|_{x=x_0} \\
    &h^+ \xrightarrow{} 0=es\delta p (x_0)-eD_p\left. \diff[]{p}{x}\right|_{x=x_0}
\end{aligned} \right.
\end{equation}
Mientras que en el caso derecho
\begin{equation}
   J_\mathrm{rec}(x_0)=J_\mathrm{d}(x_0) \left \{
\begin{aligned}
    &e^- \xrightarrow{} eD_n\left. \diff[]{n}{x}\right|_{x=x_0}=-es\delta n (x_0)+eD_n \\
    &h^+ \xrightarrow{} -eD_n\left. \diff[]{p}{x}\right|_{x=x_0}=es\delta p (x_0)-eD_p
\end{aligned} \right.
\end{equation}



\subsection{Generación superficial en $x_0$}
La siguiente figura nos ilustra los distintos flujos de corriente que se dan en la superficie de un semiconductor.
\begin{figure}[h]
    \centering
    \begin{subfigure}[h]{0.3\linewidth}
\includegraphics{Standalone/ApA_04_flujo_de_corriente_i_gen.pdf}
        \caption{Diagrama del flujo de corriente en la región del límite izquierdo de un semiconductor}
        \label{fig: flujo de corriente i gen}
    \end{subfigure}
    \hfill
    \begin{subfigure}[h]{0.3\linewidth}
            \includegraphics{Standalone/ApA_05_flujo_de_corriente_g_gen.pdf}
        \caption{Diagrama de flujo de corriente en una región interna del semiconductor}
        \label{fig: flujo de corriente g gen}
    \end{subfigure}
    \hfill
    \begin{subfigure}[h]{0.3\linewidth}
    \includegraphics{Standalone/ApA_06_flujo_de_corriente_d_gen.pdf}
    \caption{Diagrama del flujo de corriente en la región del límite derecho de un semicondutor}
    \label{fig: flujo de corriente d gen}
    \end{subfigure} 
    \caption{Diagrama de flujos de corriente en distintas regiones de un semiconductor}
    \label{fig: flujo de corriente gen}
\end{figure}

\noindent En este caso la generación viene representada por el coeficiente \textit{M} que es el número de pares $e^- \, h^+$ que se generan por unidad de superficie y tiempo y vemos que es independiente del exceso de concentración $J_\mathrm{gen}(x_0)=qM$.
\subsubsection{Caso general. Posición en el interior de la muestra}
\noindent Está representado por la figura \ref{fig: flujo de corriente g gen}. Por conservación de la corriente 
\begin{equation}
J_\mathrm{dI}(x_0)+J_\mathrm{gen}(x_0)=J_\mathrm{dII}(x_0) \left \{
\begin{aligned}
    &e^- \xrightarrow{} eD_n \left. \diff[]{n_\mathrm{I}}{x}\right|_{x=x_0}+(-e)M=eD_n\left. \diff[]{n_\mathrm{II}}{x}\right|_{x=x_0} \\
    &h^+ \xrightarrow{} -eD_p \left. \diff[]{p_\mathrm{I}}{x}\right|_{x=x_0}+eM (x_0)=-eD_p\left. \diff[]{p_\mathrm{II}}{x}\right|_{x=x_0}
\end{aligned}
\right.
\end{equation}
Vemos que tenemos signo + en los $e^-$ y signos - en los $h^+$ en las corrientes de difusión.
\newpage
\subsubsection{Casos en los límites izquierdo y derecho.}
Están representados por la figura \ref{fig: flujo de corriente i} y \ref{fig: flujo de corriente d}. En estos casos sólo tenemos una corriente de difusión, que en el caso izquierdo describen
\begin{equation}
   J_\mathrm{gen}(x_0)=J_\mathrm{d}(x_0) \left \{
\begin{aligned}
    &e^- \xrightarrow{} eD_n \left. \diff[]{n}{x}\right|_{x=x_0}=-eM \\
    &h^+ \xrightarrow{} -eD_n \left. \diff[]{p}{x}\right|_{x=x_0}=eM
\end{aligned} \right.
\end{equation}

Mientras que en el caso derecho
\begin{equation}
   0=J_\mathrm{gen}(x_0)+J_\mathrm{d}(x_0) \left \{
\begin{aligned}
    &e^- \xrightarrow{} 0=-eM+eD_n\left. \diff[]{n}{x}\right|_{x=x_0} \\
    &h^+ \xrightarrow{} 0=eM-eD_p\left. \diff[]{p}{x}\right|_{x=x_0}
\end{aligned} \right.
\end{equation}
\subsection{Superficie (ideal) sin recombinación ni generación superficial}
Supongamos que tenemos un semiconductor que esta siendo iluminado en un lado, mientras que en otro no. Por ser una superficie ideal, supondremos que no hay más que las corrientes generadas por el exceso de portadores.

\begin{figure}[h]
    \centering
    \begin{subfigure}[h]{0.3\linewidth}
\includegraphics{Standalone/ApA_07_flujo_de_corriente_i_gen.pdf}
        \caption{Diagrama del flujo de corriente en la región del límite izquierdo de un semiconductor}
        \label{fig: flujo de corriente i gen}
    \end{subfigure}
    \hfill
    \begin{subfigure}[h]{0.3\linewidth}
            \includegraphics{Standalone/ApA_08_flujo_de_corriente_g_id.pdf}
        \caption{Diagrama de flujo de corriente en una región interna del semiconductor}
        \label{fig: flujo de corriente g id}
    \end{subfigure}
    \hfill
    \begin{subfigure}[h]{0.3\linewidth}
    \includegraphics{Standalone/ApA_09_flujo_de_corriente_d_id.pdf}
    \caption{Diagrama del flujo de corriente en la región del límite derecho de un semicondutor}
    \label{fig: flujo de corriente d id}
    \end{subfigure} 
    \caption{Diagrama de flujos de corriente en distintas regiones de un semiconductor}
    \label{fig: flujo de corriente id}
\end{figure}

Este último caso únicamente nos exige buscar que $\partial_x \delta n (x)$ sea continua en $x=x_0$, para el caso descrito en la figura \ref{fig: flujo de corriente g id} y nula en los casos descritos por las figuras \ref{fig: flujo de corriente id}\ref{fig: flujo de corriente d id}

\end{document}